% Étude des écosystèmes : La forêt et la savane — SVTEEHB 3ème — Cours signé OnBuch+
% Programme officiel MINESEC. Compiler avec : tectonic N14-etude-ecosystemes-foret-savane.tex
\newcommand{\DOCMATIERE}{SVTEEHB}
\newcommand{\DOCNIVEAU}{3ème}
\newcommand{\DOCMODULE}{SVTEEHB – classe de 3ème}
\newcommand{\DOCLECON}{N14}
\newcommand{\DOCTITRE}{Étude des écosystèmes : La forêt et la savane}
% OnBuch+ — préambule « cours signés » ENRICHI (compilé avec tectonic / XeLaTeX)
% Système visuel « Pop » d'OnBuch+ : fond crème (jamais blanc), bordures encre
% épaisses, ombres « dures » décalées, titres Archivo Black en capitales.
% Version MATHÉMATIQUES : graphes (pgfplots/tikz), boîtes pédagogiques
% (définition, propriété, méthode, attention, le savais-tu, exercice résolu,
% exercice, TP, bilan), page de garde et signature OnBuch+ ancrée en bas.
%
% Macros attendues dans main.tex AVANT \input{preamble} :
%   \DOCMATIERE  \DOCNIVEAU  \DOCMODULE  \DOCLECON (numéro)  \DOCTITRE
\documentclass[11pt]{article}

\usepackage{fontspec}
\usepackage[french]{babel}
\usepackage[a4paper,margin=2cm,top=3.4cm,bottom=2.8cm,headheight=1.4cm,footskip=1.3cm]{geometry}
\usepackage{amsmath,amssymb,amsfonts}
\usepackage{mathtools} % \xmapsto, \coloneqq... souvent utilisés par les modèles
\usepackage{cancel} % simplifications barrées (bilans de forces)
% \abs{z} (module, valeur absolue) et \norm{u} (norme), utilisés par les modèles
\providecommand{\abs}{}\let\abs\relax\DeclarePairedDelimiter{\abs}{\lvert}{\rvert}
\providecommand{\norm}{}\let\norm\relax\DeclarePairedDelimiter{\norm}{\lVert}{\rVert}
\usepackage[table]{xcolor} % [table] : fournit aussi \rowcolors (alternance de lignes)
\usepackage{pagecolor}
\usepackage{tikz}
\usetikzlibrary{arrows.meta,positioning,calc,shapes.geometric,shapes.misc,shapes.arrows,decorations.pathmorphing,decorations.pathreplacing,decorations.markings,patterns,shadows,mindmap,trees,fit,backgrounds,matrix,intersections,angles,quotes}
\usepackage{pgfplots}
\usepackage[version=4]{mhchem} % équations nucléaires \ce{^{235}_{92}U}
\usepackage[european,siunitx]{circuitikz} % schémas électriques
\pgfplotsset{compat=1.18}
\usepgfplotslibrary{fillbetween,groupplots} % aires entre courbes (calcul intégral)
\usepackage{enumitem}
\usepackage{fancyhdr}
% [most] charge toutes les bibliothèques tcolorbox (listings, minted, raster,
% documentation...) et gonfle inutilement la mémoire TeX sur un document long
% (~300 boîtes) : on ne charge que ce qui est réellement utilisé.
\usepackage[skins,breakable]{tcolorbox}
\usepackage{graphicx}
\usepackage{colortbl}
\usepackage{booktabs}
\usepackage{tabularx}
\usepackage{multirow}
\usepackage{diagbox} % case d'en-tête diagonale des tableaux à double entrée
% Colonnes extensibles centrée / gauche / droite (C, L, R), souvent utilisées par les modèles
\newcolumntype{C}{>{\centering\arraybackslash}X}
\newcolumntype{L}{>{\raggedright\arraybackslash}X}
\newcolumntype{R}{>{\raggedleft\arraybackslash}X}
\usepackage{array}
\usepackage{multicol}
\usepackage{siunitx}
% parse-numbers=false : accepte aussi \SI{2,5 \times 10^{-3}}{...}
\sisetup{locale=FR,output-decimal-marker={,},group-minimum-digits=4,parse-numbers=false}
\DeclareSIUnit\ohm{\ensuremath{\symup{\Omega}}} % U+2126 absent de la police
\DeclareSIUnit\division{div} % réglages d'oscilloscope (ms/div, V/div)
\DeclareSIUnit\tr{tr} % tours (vitesse de rotation, ex. \SI{1500}{\tr\per\minute})
\DeclareSIUnit\molar{M} % concentration molaire (ex. \SI{3}{\molar}, \SI{1}{\micro\molar})
\DeclareSIUnit\an{an} % année (le modèle confond parfois avec \year, un compteur TeX) — ex. \SI{5}{\kilo\meter\per\an}
\DeclareSIUnit\FCFA{FCFA} % franc CFA (ex. \SI{500}{\FCFA\per\kilo\gram})
\DeclareSIUnit\degreeBrix{\textdegree Brix} % teneur en sucre d'un sirop/jus (réfractométrie)
\DeclareSIUnit\month{mois} % le modèle utilise parfois \month, un compteur TeX, comme unité de durée
\DeclareSIUnit\microliter{\micro\liter} % le modèle écrit parfois \microliter en un seul token
\DeclareSIUnit\million{millions} % dénombrement (ex. \SI{15}{\million\per\mL} spermatozoïdes)
\usepackage{qrcode}
% \degree : commande fréquemment utilisée par les modèles pour un angle ou une
% température en dehors de \SI{}{} (non fournie par les paquets chargés).
\providecommand{\degree}{^{\circ}}
% \qed (fin de démonstration), utilisé par les modèles sans amsthm
\providecommand{\qed}{\hfill$\square$}
% \barre : double barre des valeurs interdites dans un tableau de variations (tkz-tab)
\providecommand{\barre}{\ensuremath{\|}}
% environnement proof (amsthm non chargé) utilisé par les modèles
\ifdefined\proof\else\newenvironment{proof}[1][Démonstration]{\par\noindent\textit{#1.}\ }{\qed\par}\fi
\usepackage{lastpage}
\usepackage{hyperref}
\hypersetup{hidelinks}

% ============================================================
% Palette « Pop » OnBuch+ — mêmes hex que la classe P (plus_ui.dart)
% ============================================================
\definecolor{poporange}{HTML}{F59321}
\definecolor{popdark}{HTML}{E07A0C}
\definecolor{popcream}{HTML}{FFF6E8}
\definecolor{popink}{HTML}{1C1714}
\definecolor{poppurple}{HTML}{7A5AE0}
\definecolor{popgreen}{HTML}{2E9E6A}
\definecolor{poppink}{HTML}{E0457B}
\definecolor{popblue}{HTML}{2F7DE1}
\definecolor{popgold}{HTML}{C98A12}
\definecolor{popmuted}{HTML}{8A7F76}
\definecolor{popline}{HTML}{EADFD2}
\definecolor{popgray}{HTML}{8A7F76}
\colorlet{popgrey}{popgray}
% Alias tolérants (le modèle nomme parfois les couleurs différemment)
\colorlet{popwhite}{white}
\colorlet{popblack}{popink}
\colorlet{popred}{poppink}
\colorlet{popyellow}{popgold}
% Teintes claires (fonds de tableaux/encadrés) — le modèle nomme parfois
% les couleurs avec un suffixe L (ex. popblueL) sans les avoir définies.
\colorlet{poporangeL}{poporange!20}
\colorlet{popdarkL}{popdark!20}
\colorlet{poppurpleL}{poppurple!20}
\colorlet{popgreenL}{popgreen!20}
\colorlet{poppinkL}{poppink!20}
\colorlet{popblueL}{popblue!20}
\colorlet{popgoldL}{popgold!20}
\colorlet{popredL}{popred!20}
\colorlet{popgrayL}{popgray!20}
% Teintes claires pour fonds de tableaux / zones
\colorlet{popblueL}{popblue!10!white}
\colorlet{poporangeL}{poporange!14!white}
\colorlet{popgreenL}{popgreen!12!white}
\colorlet{poppurpleL}{poppurple!12!white}
\colorlet{poppinkL}{poppink!12!white}
\colorlet{popgoldL}{popgold!14!white}

\pagecolor{popcream}

% ============================================================
% Typographie
% ============================================================
\defaultfontfeatures{Ligatures=TeX}
\setmainfont{PlusJakartaSans-Medium.ttf}[Path=../../../fonts/,BoldFont=PlusJakartaSans-Bold.ttf]
% Maths en sans-serif (Fira Math) avec chiffres et lettres droites de Plus
% Jakarta, pour que les formules se fondent dans le texte.
\usepackage[math-style=ISO,bold-style=ISO]{unicode-math}
\setmathfont{FiraMath-Regular.otf}[Path=../../../fonts/]
\setmathfont{PlusJakartaSans-Medium.ttf}[Path=../../../fonts/,range={up/num,up/{Latin,latin},\mathup}]
\usepackage{icomma} % virgule décimale sans espace en mode math
% Caractères mathématiques Unicode absents des polices de texte (Plus Jakarta,
% Archivo Black) : rendus par leur équivalent mathématique, en texte comme en
% formule — sinon ils s'affichent en carré vide (ex. « Arithmétique dans ℤ »).
\usepackage{newunicodechar}
\newunicodechar{ℝ}{\ensuremath{\mathbb{R}}}
\newunicodechar{ℤ}{\ensuremath{\mathbb{Z}}}
\newunicodechar{ℂ}{\ensuremath{\mathbb{C}}}
\newunicodechar{ℕ}{\ensuremath{\mathbb{N}}}
\newunicodechar{ℚ}{\ensuremath{\mathbb{Q}}}
\newunicodechar{𝔻}{\ensuremath{\mathbb{D}}}
\newunicodechar{α}{\ensuremath{\alpha}}
\newunicodechar{β}{\ensuremath{\beta}}
\newunicodechar{γ}{\ensuremath{\gamma}}
\newunicodechar{δ}{\ensuremath{\delta}}
\newunicodechar{ε}{\ensuremath{\varepsilon}}
\newunicodechar{θ}{\ensuremath{\theta}}
\newunicodechar{λ}{\ensuremath{\lambda}}
\newunicodechar{μ}{\ensuremath{\mu}}
\newunicodechar{σ}{\ensuremath{\sigma}}
\newunicodechar{τ}{\ensuremath{\tau}}
\newunicodechar{φ}{\ensuremath{\varphi}}
\newunicodechar{ψ}{\ensuremath{\psi}}
\newunicodechar{ω}{\ensuremath{\omega}}
\newunicodechar{Σ}{\ensuremath{\Sigma}}
% majuscules produites par \MakeUppercase dans les titres (α-aminés -> Α)
\newunicodechar{Α}{\ensuremath{\alpha}}
\newunicodechar{Β}{\ensuremath{\beta}}
\newunicodechar{Γ}{\ensuremath{\Gamma}}
\newunicodechar{Δ}{\ensuremath{\Delta}}
\newunicodechar{Ε}{\ensuremath{\varepsilon}}
\newunicodechar{Θ}{\ensuremath{\Theta}}
\newunicodechar{Λ}{\ensuremath{\Lambda}}
\newunicodechar{Μ}{\ensuremath{\mu}}
\newunicodechar{Τ}{\ensuremath{\tau}}
\newunicodechar{Φ}{\ensuremath{\Phi}}
\newunicodechar{Ψ}{\ensuremath{\Psi}}
\newunicodechar{Ω}{\ensuremath{\Omega}}
\newunicodechar{↦}{\ensuremath{\mapsto}}
\newunicodechar{∈}{\ensuremath{\in}}
\newunicodechar{∉}{\ensuremath{\notin}}
\newunicodechar{∧}{\ensuremath{\wedge}}
\newunicodechar{⇔}{\ensuremath{\Leftrightarrow}}
\newunicodechar{⟺}{\ensuremath{\Longleftrightarrow}}
\newunicodechar{⇒}{\ensuremath{\Rightarrow}}
\newunicodechar{⟹}{\ensuremath{\Longrightarrow}}
\newunicodechar{∘}{\ensuremath{\circ}}
\newunicodechar{∪}{\ensuremath{\cup}}
\newunicodechar{∩}{\ensuremath{\cap}}
\newunicodechar{≡}{\ensuremath{\equiv}}
\newunicodechar{⊂}{\ensuremath{\subset}}
\newunicodechar{⊥}{\ensuremath{\perp}}
\newunicodechar{∃}{\ensuremath{\exists}}
\newunicodechar{∀}{\ensuremath{\forall}}
\newunicodechar{∛}{\ensuremath{\sqrt[3]{\,}}}
\newunicodechar{∜}{\ensuremath{\sqrt[4]{\,}}}
\newunicodechar{⃗}{\ensuremath{{}^{\to}}}
\newunicodechar{ⁿ}{\ifmmode^{n}\else\textsuperscript{\ensuremath{n}}\fi}
\newunicodechar{ˣ}{\ifmmode^{x}\else\textsuperscript{\ensuremath{x}}\fi}
\newunicodechar{ᵇ}{\ifmmode^{b}\else\textsuperscript{\ensuremath{b}}\fi}
\newunicodechar{ʳ}{\ifmmode^{r}\else\textsuperscript{\ensuremath{r}}\fi}
\newunicodechar{ᵘ}{\ifmmode^{u}\else\textsuperscript{\ensuremath{u}}\fi}
\newunicodechar{ʸ}{\ifmmode^{y}\else\textsuperscript{\ensuremath{y}}\fi}
\newunicodechar{ᵃ}{\ifmmode^{a}\else\textsuperscript{\ensuremath{a}}\fi}
\newunicodechar{ᵉ}{\ifmmode^{e}\else\textsuperscript{\ensuremath{e}}\fi}
\newunicodechar{ᵏ}{\ifmmode^{k}\else\textsuperscript{\ensuremath{k}}\fi}
\newunicodechar{ᵖ}{\ifmmode^{p}\else\textsuperscript{\ensuremath{p}}\fi}
\newunicodechar{ᴺ}{\ifmmode^{N}\else\textsuperscript{\ensuremath{N}}\fi}
\newunicodechar{ᶜ}{\ifmmode^{c}\else\textsuperscript{\ensuremath{c}}\fi}
\newunicodechar{ʰ}{\ifmmode^{h}\else\textsuperscript{\ensuremath{h}}\fi}
\newunicodechar{ᵠ}{\ifmmode^{q}\else\textsuperscript{\ensuremath{q}}\fi}
\newunicodechar{ₙ}{\ifmmode_{n}\else\textsubscript{\ensuremath{n}}\fi}
\newunicodechar{ₐ}{\ifmmode_{a}\else\textsubscript{\ensuremath{a}}\fi}
\newunicodechar{ᵢ}{\ifmmode_{i}\else\textsubscript{\ensuremath{i}}\fi}
\newunicodechar{ᵣ}{\ifmmode_{r}\else\textsubscript{\ensuremath{r}}\fi}
\newunicodechar{ₖ}{\ifmmode_{k}\else\textsubscript{\ensuremath{k}}\fi}
\newunicodechar{ᵦ}{\ifmmode_{\beta}\else\textsubscript{\ensuremath{\beta}}\fi}
\newunicodechar{ₓ}{\ifmmode_{x}\else\textsubscript{\ensuremath{x}}\fi}
\newfontfamily\popdisplay{ArchivoBlack-Regular.ttf}[Path=../../../fonts/]
\newfontfamily\poplogo{SpaceGrotesk-Bold.ttf}[Path=../../../fonts/]
\newfontfamily\popsemi{PlusJakartaSans-SemiBold.ttf}[Path=../../../fonts/]
\newfontfamily\popxbold{PlusJakartaSans-ExtraBold.ttf}[Path=../../../fonts/]
\newfontfamily\popxboldls{PlusJakartaSans-ExtraBold.ttf}[Path=../../../fonts/,LetterSpace=3.0]

\setlist[enumerate]{leftmargin=1.7em,itemsep=5pt,topsep=4pt}
\setlist[itemize]{leftmargin=1.7em,itemsep=4pt,topsep=4pt,label={\color{poporange}\textbullet}}
\setlength{\parindent}{0pt}
\setlength{\parskip}{5pt}
\renewcommand{\arraystretch}{1.25}

% pgfplots : style Pop par défaut
\pgfplotsset{
  every axis/.append style={axis line style={popink,line width=1pt},tick style={popink},
    grid style={popline,line width=0.5pt},label style={font=\small},tick label style={font=\footnotesize}},
  popcurve/.style={poporange,line width=1.8pt,no markers,smooth},
}
% Même style disponible dans un \draw TikZ simple (hors pgfplots)
\tikzset{popcurve/.style={poporange,line width=1.8pt}}
\tikzset{popgrid/.style={popline,very thin}}
\colorlet{popcurve}{poporange} % le nom du style est parfois utilisé comme couleur

% ============================================================
% Pill / squiggle
% ============================================================
\newcommand{\poppill}[3]{%
  \tikz[baseline=-0.55ex]\node[draw=popink,line width=1.2pt,rounded corners=2.6pt,%
    fill=#2,inner xsep=6.5pt,inner ysep=3pt]{%
    \popxboldls\fontsize{7}{7}\selectfont\color{#3}\MakeUppercase{#1}};%
}
\newcommand{\popsquiggle}[1]{%
  \tikz[baseline=-0.4ex]{
    \draw[#1,line width=2.6pt,line cap=round]
      plot [smooth] coordinates {(0,0.09) (0.34,0.01) (0.68,0.21) (1.02,0.01) (1.36,0.10)};
  }%
}
% Mot-clé surligné (marqueur orange sous le texte)
\newcommand{\cle}[1]{{\popxbold\color{popdark}#1}}

% ============================================================
% En-tête / pied
% ============================================================
\pagestyle{fancy}
\fancyhf{}
\renewcommand{\headrulewidth}{0pt}
\renewcommand{\footrulewidth}{0pt}
\lhead{\raisebox{-0.15ex}{\poplogo\fontsize{13}{13}\selectfont\color{popink}OnBuch\color{poporange}+}}
\rhead{\poppill{\DOCMATIERE\ \textbullet\ \DOCNIVEAU\ \textbullet\ Leçon \DOCLECON}{white}{popink}}
\lfoot{\popxbold\fontsize{7.6}{7.6}\selectfont\color{popmuted}\MakeUppercase{Cours signé OnBuch+}\ \textbullet\ {\popsemi\DOCTITRE}}
\rfoot{\tikz[baseline=-0.6ex]\node[rounded corners=8.5pt,fill=popink,minimum height=17pt,minimum width=17pt,inner xsep=5pt,inner sep=0]{\popxbold\fontsize{8}{8}\selectfont\color{popcream}\thepage\,/\,\pageref*{LastPage}};}

% ============================================================
% Titres
% ============================================================
\newcommand{\chaptitre}[2]{%
  \begin{tcolorbox}[enhanced,colback=poporange,colframe=popink,boxrule=2.6pt,arc=11pt,
      drop shadow=popink,left=15pt,right=15pt,top=12pt,bottom=11pt,before skip=3pt,after skip=18pt]
    \poppill{#2}{popink}{popcream}\par\vspace{9pt}
    {\raggedright\hyphenpenalty=10000\exhyphenpenalty=10000\popdisplay\fontsize{22}{24}\selectfont\color{popcream}\MakeUppercase{#1}\par}
  \end{tcolorbox}
}
% Section numérotée : pastille orange avec le numéro + titre Archivo
\newcounter{popsec}
\newcommand{\coursec}[1]{%
  \stepcounter{popsec}\setcounter{popsub}{0}%
  \par\vspace{16pt}\noindent
  \begin{minipage}{\linewidth}
  \tikz[baseline=-0.7ex]\node[draw=popink,line width=1.6pt,fill=poporange,rounded corners=4pt,minimum size=22pt,inner sep=0]{\popdisplay\fontsize{12}{12}\selectfont\color{popcream}\thepopsec};%
  \hspace{8pt}{\popdisplay\fontsize{15}{17}\selectfont\color{popink}\MakeUppercase{#1}}\par
  \vspace{4pt}\noindent{\color{poporange}\rule{36pt}{3.6pt}}
  \end{minipage}\par\nopagebreak\vspace{8pt}
}
\newcounter{popsub}
\newcommand{\courssub}[1]{%
  \stepcounter{popsub}%
  \par\vspace{9pt}\noindent{\popxbold\fontsize{11.5}{13}\selectfont\color{popink}{\color{poporange}\thepopsec.\thepopsub}\hspace{6pt}#1}\par\nopagebreak\vspace{4pt}
}

% ============================================================
% Boîtes pédagogiques — même recette : fond blanc, bordure épaisse colorée,
% ombre dure encre, badge plein. Titre optionnel [#1].
% ============================================================
\tcbset{popbase/.style={breakable,enhanced,colback=white,boxrule=2.3pt,arc=9pt,
  shadow={3pt}{-3pt}{0pt}{popink},left=14pt,right=14pt,top=11pt,bottom=11pt,before skip=11pt,after skip=15pt,
  fontupper=\popsemi\fontsize{10.6}{14.5}\selectfont\color{popink},
  fontlower=\popsemi\fontsize{10.6}{14.5}\selectfont\color{popink}}}
% Coche dessinée (le glyphe ✓ n'existe pas dans les polices du document)
\newcommand{\popcheck}[1]{\tikz[baseline=-0.5ex]\draw[#1,line width=1.6pt,line cap=round,line join=round](-0.13,0.0)--(-0.04,-0.09)--(0.13,0.1);}
\providecommand{\coche}{\popcheck{popgreen}}
\providecommand{\resultat}[1]{\par\smallskip\noindent\fcolorbox{popgreen}{popgreenL}{\parbox{\dimexpr\linewidth-2\fboxsep-2\fboxrule\relax}{\textbf{Résultat.} #1}}\par\smallskip}
\newcommand{\popboxhead}[3]{\poppill{#1}{#2}{white}\ifx\relax#3\relax\else\hspace{6pt}{\popxbold\fontsize{10.6}{12}\selectfont\color{popink}#3}\fi\par\vspace{7pt}}

\newtcolorbox{aretenir}[1][]{popbase,colframe=popgreen,before upper={\popboxhead{À retenir}{popgreen}{#1}}}
\newtcolorbox{exemplebox}[1][]{popbase,colframe=poporange,before upper={\popboxhead{Exemple}{poporange}{#1}}}
\newtcolorbox{definition}[1][]{popbase,colframe=popblue,before upper={\popboxhead{Définition}{popblue}{#1}}}
\newtcolorbox{propriete}[1][]{popbase,colframe=poppurple,before upper={\popboxhead{Propriété}{poppurple}{#1}}}
\newtcolorbox{methode}[1][]{popbase,colframe=popgold,before upper={\popboxhead{Méthode}{popgold}{#1}}}
\newtcolorbox{attention}[1][]{popbase,colframe=poppink,before upper={\popboxhead{Attention}{poppink}{#1}}}
\newtcolorbox{experience}[1][]{popbase,colframe=popblue,colback=popblueL,before upper={\popboxhead{Expérience}{popblue}{#1}}}
% « Le savais-tu ? » : carte violette pleine (contexte, culture, Cameroun)
\newtcolorbox{savaistu}[1][]{popbase,colback=poppurple,colframe=popink,
  fontupper=\popsemi\fontsize{10.4}{14.2}\selectfont\color{white},
  before upper={\poppill{Le savais-tu\,?}{white}{poppurple}\ifx\relax#1\relax\else\hspace{6pt}{\popxbold\color{white}#1}\fi\par\vspace{7pt}}}
% Exercice résolu : énoncé en haut, \tcblower puis la solution sur fond crème
\newcounter{popexo}\newcounter{popexr}
\newtcolorbox{exoresolu}[1][]{popbase,colframe=poporange,colbacklower=popcream,
  segmentation style={popink,line width=1.2pt,dashed},
  before upper={\stepcounter{popexr}\popboxhead{Exercice résolu \thepopexr}{poporange}{#1}},
  before lower={\poppill{Solution}{popink}{popcream}\par\vspace{6pt}}}
% Exercice (non résolu en place) — la correction est dans la partie Corrigés
\newtcolorbox{exercice}[1][]{popbase,colframe=popink,
  before upper={\stepcounter{popexo}\popboxhead{Exercice \thepopexo}{popink}{#1}}}
% Corrigé
\newtcolorbox{corrige}[1][]{popbase,colframe=popgreen,colback=popgreenL,
  before upper={\popboxhead{Corrigé}{popgreen}{#1}}}
% Objectifs / prérequis / bilan
\newtcolorbox{objectifs}{popbase,colframe=popink,colback=poporangeL,
  before upper={\popboxhead{Objectifs de la leçon}{popink}{}}}
\newtcolorbox{prerequis}{popbase,colframe=popink,colback=popblueL,
  before upper={\popboxhead{Prérequis}{popblue}{}}}
\newtcolorbox{bilan}{popbase,colframe=popink,colback=white,boxrule=2.8pt,
  before upper={\popboxhead{Fiche bilan}{poporange}{L'essentiel à connaître pour l'examen}}}
% Figure encadrée avec légende
\newtcolorbox{popfigure}[1][]{enhanced,breakable=false,colback=white,colframe=popline,boxrule=1.4pt,arc=8pt,
  left=10pt,right=10pt,top=10pt,bottom=8pt,before skip=10pt,after skip=12pt,halign=center,
  lower separated=false,halign lower=center,
  fontlower=\popsemi\fontsize{9}{11.5}\selectfont\color{popmuted},
  #1}
\newcommand{\legende}[1]{\par\vspace{6pt}{\popsemi\fontsize{9}{11.5}\selectfont\color{popmuted}#1}}

% ============================================================
% Signature OnBuch+
% ============================================================
\newcommand{\poplogoinline}[1]{{\poplogo\fontsize{#1}{#1}\selectfont\color{popink}OnBuch\color{poporange}+}}
\newcommand{\signatureonbuch}{%
  \begin{tcolorbox}[enhanced,colback=popink,colframe=popink,boxrule=2.2pt,arc=10pt,
      drop shadow=poporange,left=14pt,right=14pt,top=10pt,bottom=10pt,before skip=0pt,after skip=0pt]
    \tikz[baseline=-0.7ex]\node[circle,fill=popgreen,draw=popcream,line width=1.3pt,minimum size=22pt,inner sep=0]{\popcheck{white}};%
    \hspace{9pt}\begin{minipage}[c]{0.62\linewidth}
      {\popxbold\fontsize{12}{13}\selectfont\color{popcream}Signé\ \textbullet\ {\poplogo OnBuch\color{poporange}+}}\par\vspace{2pt}
      {\raggedright\popsemi\fontsize{8}{10.5}\selectfont\color{popline}\MakeUppercase{Équipe pédagogique OnBuch+}\\\MakeUppercase{Conforme au programme officiel MINESEC}\par}
    \end{minipage}\hfill
    {\poplogo\fontsize{22}{22}\selectfont\color{popcream}OnBuch\color{poporange}+}
  \end{tcolorbox}%
}

% ============================================================
% Page de garde
% #1 titre · #2 matière · #3 niveau · #4 module · #5 n° leçon · #6 durée ·
% #7 description / objectifs (texte ou itemize)
% ============================================================
\newcommand{\pagegarde}[7]{%
  \thispagestyle{empty}
  \newgeometry{a4paper,margin=1.7cm,top=1.6cm,bottom=1.4cm}%
  \begin{tikzpicture}[remember picture,overlay]
    \fill[poporange!18] ([xshift=-3.2cm,yshift=-3.6cm]current page.north east) circle (4.6cm);
    \fill[poppurple!14] ([xshift=2.4cm,yshift=7.6cm]current page.south west) circle (3.8cm);
  \end{tikzpicture}%
  {\poplogo\fontsize{30}{30}\selectfont\color{popink}OnBuch\color{poporange}+}\hfill
  \raisebox{0.6ex}{\poppill{Cours signé \textbullet\ Premium}{popink}{popcream}}\par
  \vspace{1pt}\popsquiggle{poppurple}\par
  \vspace{8pt}
  \begin{tcolorbox}[enhanced,colback=poporange,colframe=popink,boxrule=2.8pt,arc=13pt,
      drop shadow=popink,left=18pt,right=18pt,top=12pt,bottom=13pt,after skip=10pt]
    \poppill{#2 \textbullet\ #3}{popink}{popcream}\hspace{5pt}\poppill{#4}{white}{popink}\par\vspace{8pt}
    {\popdisplay\fontsize{14}{14}\selectfont\color{popink}LEÇON #5}\par\vspace{4pt}
    {\raggedright\hyphenpenalty=10000\exhyphenpenalty=10000\popdisplay\fontsize{25}{27}\selectfont\color{popcream}\MakeUppercase{#1}\par}
    \vspace{8pt}
    {\popxbold\fontsize{9.5}{9.5}\selectfont\color{popink}\MakeUppercase{Durée conseillée : #6}}
  \end{tcolorbox}
  \begin{tcolorbox}[enhanced,colback=white,colframe=popink,boxrule=2.2pt,arc=10pt,
      drop shadow=popink,left=16pt,right=16pt,top=10pt,bottom=10pt,after skip=8pt]
    \poppill{Dans cette leçon}{poppurple}{white}\par\vspace{5pt}
    \popsemi\fontsize{9.3}{12.6}\selectfont\color{popink}#7
  \end{tcolorbox}
  \noindent\begin{minipage}[t]{0.487\linewidth}
    \begin{tcolorbox}[enhanced,colback=popgreen,colframe=popink,boxrule=2.2pt,arc=10pt,
        drop shadow=popink,left=11pt,right=11pt,top=10pt,bottom=10pt,height=3.1cm,valign=center]
      \tcbox[colback=white,colframe=popink,boxrule=1.3pt,arc=5pt,boxsep=0pt,nobeforeafter,
        left=3pt,right=3pt,top=3pt,bottom=3pt,valign=center]{\qrcode[height=1.9cm]{https://play.google.com/store/apps/details?id=com.luvvix.onbuch&hl=fr}}%
      \hspace{8pt}\begin{minipage}[c]{0.5\linewidth}
        {\popdisplay\fontsize{11}{12.5}\selectfont\color{white}\MakeUppercase{Télécharge OnBuch}}\par\vspace{4pt}
        {\raggedright\popsemi\fontsize{8.4}{11}\selectfont\color{white}Gratuit \textbullet\ Google Play\\Tuteur IA, annales, concours, résultats\par}
      \end{minipage}
    \end{tcolorbox}
  \end{minipage}\hfill
  \begin{minipage}[t]{0.487\linewidth}
    \begin{tcolorbox}[enhanced,colback=poppurple,colframe=popink,boxrule=2.2pt,arc=10pt,
        drop shadow=popink,left=11pt,right=11pt,top=10pt,bottom=10pt,height=3.1cm,valign=center]
      \tikz[baseline=-0.7ex]\node[circle,fill=white,draw=popink,line width=1.3pt,minimum size=1.6cm,inner sep=0]{\popdisplay\fontsize{24}{24}\selectfont\color{poppurple}+};%
      \hspace{8pt}\begin{minipage}[c]{0.6\linewidth}
        {\popdisplay\fontsize{11}{12.5}\selectfont\color{white}\MakeUppercase{Passe à OnBuch+}}\par\vspace{4pt}
        {\raggedright\popsemi\fontsize{8.4}{11}\selectfont\color{white}Tous les cours signés, Tuteur IA illimité, fiches PDF — onglet OnBuch+ de l'app\par}
      \end{minipage}
    \end{tcolorbox}
  \end{minipage}
  \vspace{8pt}
  \popcontenu
  \vfill
  \signatureonbuch
  \restoregeometry
  \newpage
}

% Rangée « Ce cours contient » (page de garde)
\newcommand{\poptile}[3]{% #1 couleur #2 gros texte #3 légende
  \begin{tcolorbox}[enhanced,colback=#1,colframe=popink,boxrule=2pt,arc=9pt,shadow={2.5pt}{-2.5pt}{0pt}{popink},
      left=6pt,right=6pt,top=9pt,bottom=9pt,halign=center,valign=center,height=1.75cm,width=\linewidth,nobeforeafter]
    {\popdisplay\fontsize{12.5}{13}\selectfont\color{white}\MakeUppercase{#2}}\par\vspace{3pt}
    {\centering\popsemi\fontsize{7.8}{9.5}\selectfont\color{white}#3\par}
  \end{tcolorbox}}
\newcommand{\popcontenu}{%
  \par\noindent{\popxboldls\fontsize{7.5}{7.5}\selectfont\color{popmuted}CE COURS SIGNÉ CONTIENT}\par\vspace{7pt}
  \noindent\begin{minipage}[t]{0.235\linewidth}\poptile{popblue}{Cours}{complet, illustré, pas à pas}\end{minipage}\hfill
  \begin{minipage}[t]{0.235\linewidth}\poptile{poporange}{Méthodes}{et exercices résolus}\end{minipage}\hfill
  \begin{minipage}[t]{0.235\linewidth}\poptile{poppink}{Exercices}{type examen + corrigés}\end{minipage}\hfill
  \begin{minipage}[t]{0.235\linewidth}\poptile{popgreen}{Fiche bilan}{l'essentiel à retenir}\end{minipage}\par}

% Sommaire
\newtcolorbox{sommaire}{enhanced,colback=white,colframe=popink,boxrule=2.2pt,arc=10pt,
  drop shadow=popink,left=18pt,right=18pt,top=13pt,bottom=13pt,before skip=6pt,after skip=20pt,
  before upper={\poppill{Sommaire}{poppurple}{white}\par\vspace{9pt}\popsemi\fontsize{10.6}{14.5}\selectfont\color{popink}}}

% Signature de fin de cours — poussée tout en bas de la dernière page
\newcommand{\findecours}{%
  \par\vfill\nopagebreak
  \begin{center}\popsquiggle{poporange}\end{center}\vspace{4pt}
  \signatureonbuch
}
\AtBeginDocument{\def\llbracket{\mathopen{[\mkern-3.5mu[}}\def\rrbracket{\mathclose{]\mkern-3.5mu]}}} % intervalles d'entiers (glyphe ⟦ absent de la police)
% Glyphes absents de Fira Math : redéfinis à partir de glyphes disponibles
\AtBeginDocument{%
  \def\setminus{\mathbin{\backslash}}\let\smallsetminus\setminus
  \def\vdots{\mathord{\vbox{\baselineskip3.2pt\lineskiplimit0pt\kern4pt\hbox{.}\hbox{.}\hbox{.}}}}%
  \def\ddots{\mathinner{\mkern1mu\raise6.5pt\hbox{.}\mkern2mu\raise3.6pt\hbox{.}\mkern2mu\raise0.7pt\hbox{.}\mkern1mu}}%
  \def\triangle{\mathord{\scalebox{0.9}{$\symup{\Delta}$}}}%
  \def\nabla{\mathord{\raisebox{\depth}{\scalebox{1}[-1]{$\symup{\Delta}$}}}}%
  \def\checkmark{\popcheck{popdark}}%
  \def\star{\mathbin{\ast}}\def\bigstar{\mathord{\ast}}%
  \def\diamond{\mathbin{\scalebox{0.6}{$\Diamond$}}}%
  \def\bigcup{\mathop{\mathchoice{\raisebox{-0.3em}{\scalebox{1.7}{$\cup$}}}{\scalebox{1.25}{$\cup$}}{\cup}{\cup}}\displaylimits}%
  \def\bigcap{\mathop{\mathchoice{\raisebox{-0.3em}{\scalebox{1.7}{$\cap$}}}{\scalebox{1.25}{$\cap$}}{\cap}{\cap}}\displaylimits}%
  \def\lesseqqgtr{\mathrel{\lessgtr}}\def\bigoplus{\mathop{\oplus}\displaylimits}%
}
\providecommand{\ssi}{si et seulement si} % abréviation fréquente
\AtBeginDocument{\providecommand{\wideparen}{\overparen}} % arc de cercle

\begin{document}

\pagegarde{Étude des écosystèmes : La forêt et la savane}{SVTEEHB}{3ème}{SVTEEHB – classe de 3ème}{N14}{5 séances (fiche de progression harmonisée)}{Cette leçon étudie les deux grands écosystèmes du Cameroun, la forêt et la savane : leur biodiversité, les relations d'interdépendance entre leurs composants, les menaces liées aux activités humaines et les moyens de restaurer et conserver la biodiversité, notamment par la pratique du reboisement.
\vspace{4pt}
\begin{multicols}{2}\small
\begin{itemize}
\item À la fin de cette leçon, je sais définir un écosystème et identifier ses composants biotiques et abiotiques.
\item Je sais décrire la biodiversité de la forêt et de la savane camerounaises (espèces caractéristiques, stratification, adaptations).
\item Je sais mettre en évidence l'interdépendance des êtres vivants dans un écosystème à partir de chaînes et de réseaux alimentaires.
\item Je sais interpréter des documents (tableaux, courbes, cartes) montrant l'évolution d'un écosystème ou l'effet d'une perturbation.
\item Je sais citer et expliquer les activités humaines qui détruisent les écosystèmes (feux de brousse, déforestation, braconnage).
\item Je sais proposer des mesures de restauration et de conservation de la biodiversité (reboisement, aires protégées, sensibilisation).
\end{itemize}\end{multicols}}

\begin{sommaire}
\begin{multicols}{2}
\begin{enumerate}
\item Notion d'écosystème et les grands écosystèmes du Cameroun
\item Biodiversité de la forêt et de la savane
\item Interdépendance des êtres vivants dans un écosystème
\item Les activités humaines destructrices des écosystèmes
\item Restauration et conservation de la biodiversité
\item TP N°10 : Pratique du reboisement
\item Activité d'intégration
\item Méthodes et exercices résolus
\item Exercices
\item Corrigés des exercices
\item Fiche bilan
\end{enumerate}
\end{multicols}
\end{sommaire}

\begin{objectifs}
\begin{itemize}
\item À la fin de cette leçon, je sais définir un écosystème et identifier ses composants biotiques et abiotiques.
\item Je sais décrire la biodiversité de la forêt et de la savane camerounaises (espèces caractéristiques, stratification, adaptations).
\item Je sais mettre en évidence l'interdépendance des êtres vivants dans un écosystème à partir de chaînes et de réseaux alimentaires.
\item Je sais interpréter des documents (tableaux, courbes, cartes) montrant l'évolution d'un écosystème ou l'effet d'une perturbation.
\item Je sais citer et expliquer les activités humaines qui détruisent les écosystèmes (feux de brousse, déforestation, braconnage).
\item Je sais proposer des mesures de restauration et de conservation de la biodiversité (reboisement, aires protégées, sensibilisation).
\item Je sais pratiquer les étapes d'un reboisement (TP N°10) : choix des plants, pépinière, mise en terre, suivi.
\item Je sais rédiger et justifier une démarche, une réponse ou une conclusion à partir de documents.
\end{itemize}
\end{objectifs}

\begin{prerequis}
\begin{itemize}
\item Notion de milieu de vie et d'adaptation des êtres vivants (4ème).
\item Chaîne alimentaire simple : producteur, consommateur, décomposeur (4ème).
\item Rôle de la photosynthèse dans la production de matière organique par les végétaux (début d'année de 3ème).
\item Notion de reproduction et de cycle de vie des êtres vivants (4ème).
\end{itemize}
\end{prerequis}

\newpage

\coursec{Notion d'écosystème et les grands écosystèmes du Cameroun}

\courssub{Situation d'accroche et problématique}

Imagine un instant : nous sommes en janvier, en plein cœur de la saison sèche. Au Nord, dans la région de l'\textbf{Adamaoua}, le ciel est d'un bleu implacable. Soudain, une colonne de fumée noire s'élève à l'horizon : un \textbf{feu de brousse} parcourt la savane à une vitesse fulgurante, emportant herbes sèches, arbustes et tout petit animal trop lent pour fuir. Au même moment, à des centaines de kilomètres au Sud, dans la \textbf{forêt dense humide} du Dja ou du Campo, le bruit sourd des \textbf{tronçonneuses} déchire le chant des oiseaux. Des arbres séculaires — iroko, sapelli, ayous — s'effondrent pour laisser place à des champs de cacao ou alimenter le commerce du bois d'œuvre.

Pourtant, dans les deux cas, les villageois tirent la même sonnette d'alarme : \emph{« Les pluies ne viennent plus à temps », « Le gibier a disparu », « La terre ne donne plus comme avant »}.

\medskip
\noindent Comment expliquer que la destruction de la végétation entraîne la raréfaction de la pluie, la fuite des animaux et l'appauvrissement du sol ? Comment la forêt et la savane fonctionnent-elles pour maintenir la vie ? Pourquoi les êtres vivants qui y habitent dépendent-ils les uns des autres ? Et surtout, \textbf{que peut faire un élève de 3ème, toi, pour protéger ces milieux} ?

\medskip
\noindent C'est à ces questions que nous allons répondre en découvrant la notion d'\textbf{écosystème} et en explorant les deux grands écosystèmes qui font la richesse naturelle du Cameroun.

\courssub{Qu'est-ce qu'un écosystème ?}

\begin{experience}[Observation d'un milieu proche : la mare de l'école]
\textbf{Matériel :} Un seau, une épuisette, des bocaux transparents, une loupe binoculaire, un thermomètre, du papier pH, un GPS (ou smartphone), un carnet de terrain.
\textbf{Protocole :}
\begin{enumerate}
    \item Se rendre à la mare la plus proche (ou au jardin botanique, au champ en jachère).
    \item Relever les \textbf{paramètres physiques} : température de l'eau, température de l'air, pH de l'eau, profondeur approximative, nature du fond (vase, sable, cailloux), ensoleillement.
    \item Prélever délicatement de l'eau, de la vase et observer au fond des bocaux : noter \textbf{tous les êtres vivants} visibles (insectes, larves, têtards, grenouilles, nénuphars, algues, moustiques, vers).
    \item Relever des indices de présence : traces, nids, crottes, coquilles vides.
\end{enumerate}
\textbf{Observations types (exemple réel) :}
\begin{itemize}
    \item \textbf{Milieu (abiotique) :} Eau douce, pH 6,8, T° 24 °C, fond vaseux riche en matière organique, lumière traversante.
    \item \textbf{Êtres vivants (biotiques) :} Nénuphars (\emph{Nymphaea}), algues filamenteuses, moustiques (larves + adultes), grenouilles (\emph{Xenopus} ou \emph{Ptychadena}), libellules, escargots d'eau (\emph{Biomphalaria}), bactéries invisibles dans la vase.
\end{itemize}
\textbf{Interprétation :} Tu observes deux ensembles indissociables : \textbf{le milieu physique et chimique} (l'eau, la vase, la lumière, la chaleur, les sels minéraux) et \textbf{l'ensemble des êtres vivants} qui l'occupent. Mais ce n'est pas une simple coexistence : les nénuphars \textbf{produisent} de l'oxygène et de la matière organique grâce à la lumière ; les larves de moustiques \textbf{consomment} les bactéries et algues ; les grenouilles \textbf{mangent} les moustiques ; les bactéries \textbf{décomposent} les cadavres et les déchets pour rendre les sels minéraux à l'eau. \textbf{Tout est lié.}
\end{experience}

\begin{definition}[Écosystème, Biotope, Biocénose]
Un \textbf{écosystème} est un \textbf{système écologique} formé par l'association, dans un espace et un temps donnés, de deux composants indissociables en interactions constantes :
\begin{itemize}
    \item Le \textbf{biotope} (du grec \emph{bios} = vie, \emph{topos} = lieu) : c'est le \textbf{milieu de vie physique et chimique}. Il est défini par des \textbf{facteurs abiotiques} : climat (température, pluviométrie, vent), nature du sol (texture, pH, humus), eau (disponibilité, salinité, pH), lumière (intensité, durée), topographie (altitude, pente).
    \item La \textbf{biocénose} (du grec \emph{bios} = vie, \emph{koinos} = commun) : c'est l'\textbf{ensemble des êtres vivants} (populations d'espèces différentes) qui habitent ce biotope. Elle comprend les \textbf{producteurs} (végétaux chlorophylliens), les \textbf{consommateurs} (animaux) et les \textbf{décomposeurs} (champignons, bactéries).
\end{itemize}
L'écosystème n'est pas une simple addition : \textbf{biotope + biocénose + interactions = écosystème}.
\end{definition}

\begin{attention}[Piège à éviter : Écosystème $\neq$ Milieu de vie]
Ne confonds pas \textbf{écosystème} et \textbf{milieu de vie} (ou habitat).
\begin{itemize}
    \item Le \textbf{milieu de vie} (ou habitat) d'une espèce, c'est \textbf{l'endroit précis} où elle vit (ex. : la mare pour la grenouille, le tronc d'un arbre pour une termitière).
    \item L'\textbf{écosystème}, c'est \textbf{l'ensemble fonctionnel} : le milieu \textbf{ET} tous les êtres vivants \textbf{ET} leurs relations (nutrition, reproduction, abri, compétition).
\end{itemize}
\emph{Astuce mnémotechnique :} L'écosystème est un \textbf{système} (ça tourne, ça s'échange) ; le milieu de vie est un \textbf{lieu} (un adressage).
\end{attention}

\begin{exemplebox}[La mare de l'école : un écosystème complet]
Reprenons notre observation.
\begin{center}
\begin{tabularx}{\linewidth}{|l|X|X|}\hline
\textbf{Composant} & \textbf{Exemples observés} & \textbf{Rôle fonctionnel} \\ \hline
\textbf{Biotope (Abiotique)} & Eau douce, vase, lumière solaire, chaleur (24 °C), $\text{O}_2$ dissous, $\text{CO}_2$, sels minéraux (nitrates, phosphates) & Support physique, source d'énergie (lumière), réservoir de matière \\ \hline
\textbf{Biocénose (Biotique)} & \textbf{Producteurs :} Nénuphars, algues, cyanobactéries & Photosynthèse : $\text{CO}_2 + \text{H}_2\text{O} \xrightarrow{\text{lumière}} \text{Matière organique} + \text{O}_2$ \\ \cline{2-3}
& \textbf{Consommateurs :} Larves de moustiques (primaires), Grenouilles, Libellules (secondaires/tertiaires) & Nutrition hétérotrophe : mangent d'autres organismes \\ \cline{2-3}
& \textbf{Décomposeurs :} Bactéries putréfiantes, champignons microscopiques dans la vase & Minéralisation : Matière organique $\to$ Sels minéraux ($\text{NO}_3^-$, $\text{PO}_4^{3-}$) \\ \hline
\end{tabularx}
\end{center}
\textbf{Interactions observées :} Grenouille $\to$ mange Moustique ; Moustique (larve) $\to$ mange Algues/Bactéries ; Bactéries $\to$ décomposent Grenouille morte $\to$ Sels minéraux $\to$ absorbés par Nénuphars. \textbf{C'est un cycle.}
\end{exemplebox}

\begin{popfigure}
\begin{tikzpicture}[scale=0.9, every node/.style={font=\small}]
    % Biotope box
    \node[draw=popblue, thick, rounded corners, fill=popblue!10, minimum width=5cm, minimum height=3.5cm] (biotope) at (-4,0) {};
    \node[popblue, font=\bfseries\large] at (-4, 1.5) {BIOTOPE (Abiotique)};
    \begin{scope}[shift={(-4,0)}]
        \foreach \y/\txt in {1.0/{Climat : Pluie, Température, Vent}, 0.4/{Sol : Texture, pH, Humus}, -0.2/{Eau : Disponibilité, Qualité}, -0.8/{Lumière : Intensité, Durée}, -1.4/{Relief : Altitude, Pente}} {
            \node[anchor=west, text width=4.5cm] at (-2.2, \y) {\textbullet\ \txt};
        }
    \end{scope}
    % Biocenose box
    \node[draw=popgreen, thick, rounded corners, fill=popgreen!10, minimum width=5cm, minimum height=3.5cm] (biocenose) at (4,0) {};
    \node[popgreen, font=\bfseries\large] at (4, 1.5) {BIOC\'ENOSE (Biotique)};
    \begin{scope}[shift={(4,0)}]
        \node[anchor=west, text width=4.5cm] at (-2.2, 1.0) {\textbullet\ \textbf{Producteurs} (V\'eg\'etaux chlorophylliens)};
        \node[anchor=west, text width=4.5cm] at (-2.2, 0.3) {\textbullet\ \textbf{Consommateurs} (Animaux : herbivores, carnivores, omnivores)};
        \node[anchor=west, text width=4.5cm] at (-2.2, -0.4) {\textbullet\ \textbf{D\'ecomposeurs} (Champignons, Bact\'eries)};
    \end{scope}
    % Interactions arrows
    \draw[poporange, very thick, -{Stealth[length=3mm]}] (-1.5, 0.5) -- node[above, font=\small, popdark, align=center]{Facteurs limitants\\(S\'election, Stress)} (1.5, 0.5);
    \draw[poporange, very thick, -{Stealth[length=3mm]}] (1.5, -0.5) -- node[below, font=\small, popdark, align=center]{Modification du milieu\\(Humus, O$_2$, CO$_2$, Erosion)} (-1.5, -0.5);
    % Central title
    \node[draw=popdark, thick, rounded corners, fill=popcream, minimum width=10cm, minimum height=1cm] (eco) at (0, -2.5) {};
    \node[popdark, font=\bfseries\Large] at (0, -2.5) {ECOSYSTÈME = Biotope + Biocénose + INTERACTIONS};
    % Energy/Matter flows
    \draw[poppink, thick, dashed, -{Stealth[length=2mm]}] (-4, -2) -- (-4, -3) node[below, font=\tiny, poppink] {Apports énergie (Soleil) \& Matière (Pluie, Vent)};
    \draw[poppink, thick, dashed, -{Stealth[length=2mm]}] (4, -2) -- (4, -3) node[below, font=\tiny, poppink] {Pertes énergie (Chaleur) \& Matière (Ruissellement, Migration)};
\end{tikzpicture}
\legende{Schéma de la structure d'un écosystème : le biotope (milieu physique) et la biocénose (êtres vivants) échangent matière et énergie (flèches orange). L'écosystème est un système ouvert (flèches roses) recevant l'énergie solaire et la matière (pluie, vent) et perdant de la chaleur et des éléments par ruissellement ou migration.}
\end{popfigure}

\begin{aretenir}[Propriété fondamentale : L'écosystème est un système OUVERT]
Un écosystème n'est pas isolé. Il échange en permanence \textbf{matière} et \textbf{énergie} avec l'extérieur.
\begin{itemize}
    \item \textbf{Entrées d'énergie :} Rayonnement solaire (source quasi exclusive pour la photosynthèse).
    \item \textbf{Sorties d'énergie :} Chaleur dissipée (2ème principe de la thermodynamique).
    \item \textbf{Entrées de matière :} Pluie (eau, $\text{CO}_2$, poussières), vent (graines, pollens, sels), roches mères (altération $\to$ sels minéraux).
    \item \textbf{Sorties de matière :} Ruissellement (eau, sels lessivés), évaporation, migration d'animaux, exportation par l'homme (récoltes, bois).
\end{itemize}
C'est parce qu'il est \textbf{ouvert} qu'un écosystème peut s'auto-entretenir (si les apports $\ge$ pertes) ou se dégrader (si pertes $>$ apports, ex. : déforestation $\to$ lessivage massif des sels).
\end{aretenir}

\courssub{Les composants biotiques : Producteurs, Consommateurs, Décomposeurs}

Dans la biocénose, les êtres vivants ne sont pas tous au même niveau. Ils s'organisent en \textbf{réseaux trophiques} (chaînes alimentaires imbriquées) selon leur mode de nutrition.

\begin{definition}[Groupes trophiques]
\begin{itemize}
    \item \textbf{Producteurs (Autotrophes) :} Êtres vivants capables de synthétiser leur matière organique à partir de matière minérale ($\text{CO}_2$, $\text{H}_2\text{O}$, sels minéraux) en utilisant l'énergie lumineuse. \emph{Ex. :} Arbres, herbes, algues, cyanobactéries. \textbf{Base de toute chaîne alimentaire.}
    \item \textbf{Consommateurs (Hétérotrophes) :} Êtres vivants qui prélèvent leur matière organique et leur énergie en consommant d'autres êtres vivants.
    \begin{itemize}
        \item \textbf{Consommateurs primaires (Herbivores) :} Mangent les producteurs. \emph{Ex. :} Souris, gazelles, chenilles, larves de moustiques, escargots.
        \item \textbf{Consommateurs secondaires (Carnivores/Prédateurs) :} Mangent les consommateurs primaires. \emph{Ex. :} Grenouilles, lézards, petits rapaces, araignées.
        \item \textbf{Consommateurs tertiaires (Super-prédateurs) :} Mangent les consommateurs secondaires. \emph{Ex. :} Panthère, aigle martial, python, homme (parfois).
    \end{itemize}
    \item \textbf{Décomposeurs (Saprotrophes) :} Êtres vivants microscopiques (champignons, bactéries) qui dégradent la matière organique \textbf{morte} (cadavres, déchets, humus) en matière minérale (minéralisation). Ils \textbf{ferment le cycle de la matière}.
    \item \textbf{Détritivores :} Animaux (vers de terre, termites, cloportes, larves de diptères) qui fragmentent et ingèrent la matière organique morte, facilitant le travail des décomposeurs.
\end{itemize}
\end{definition}

\begin{methode}[Identifier les composants biotiques et abiotiques d'un milieu (photo ou description)]
\textbf{Objectif :} Dresser la liste structurée des composants d'un écosystème à partir d'un document.
\begin{enumerate}
    \item \textbf{Lister les facteurs abiotiques visibles ou mentionnés :} Climat (pluie, soleil, vent), Sol (type, humidité), Eau (présence, type), Lumière, Température.
    \item \textbf{Lister les êtres vivants (biocénose) :} Noter chaque espèce citée ou visible sur l'image.
    \item \textbf{Classer chaque espèce par groupe trophique :}
    \begin{itemize}
        \item Chlorophyllien (feuilles, tiges vertes) $\to$ \textbf{Producteur}.
        \item Animal mangeant plante $\to$ \textbf{Consommateur primaire}.
        \item Animal mangeant animal $\to$ \textbf{Consommateur secondaire/tertiaire}.
        \item Champignon/Bactérie sur déchet/cadavre $\to$ \textbf{Décomposeur}.
        \item Animal sur déchet/cadavre (vers, termites) $\to$ \textbf{Détritivore}.
    \end{itemize}
    \item \textbf{Repérer au moins 2 interactions :} Relation proie-prédateur, relation plante-sol (racines), décomposition.
    \item \textbf{Conclure :} « Ce milieu est un écosystème car il associe un biotope ([citer 2 facteurs]) et une biocénose ([citer 2 groupes]) en interactions ([citer 1 interaction]). »
\end{enumerate}
\end{methode}

\begin{exoresolu}[Analyse d'un document : La lisière de forêt à Dschang]
\textbf{Énoncé :} Un élève observe une lisière de forêt à Dschang (Ouest Cameroun). Il note : « Sol rouge latéritique, humide, riche en humus. Température 22 °C. Pluie fréquente. Arbres : \emph{Musanga cecropioides} (parasolier), fougères arborescentes. Au sol : fourmis, termites, vers de terre, champignons sur bois mort. Dans les branches : singes \emph{Cercopithecus}, oiseaux (touraco), chenilles sur feuilles. Un python dort enroulé. »
\textbf{Question :} Caractériser cet écosystème en identifiant le biotope, la biocénose (classée par groupes trophiques) et deux interactions.

\textbf{Solution détaillée :}
\begin{enumerate}
    \item \textbf{Biotope (Facteurs abiotiques) :} Climat tropical humide (pluie fréquente, T° 22 °C), Sol latéritique rouge, humide, riche en humus, Lumière filtrée par la canopée.
    \item \textbf{Biocénose classifiée :}
    \begin{itemize}
        \item \textbf{Producteurs :} \emph{Musanga cecropioides} (arbre pionnier), fougères arborescentes, autres arbres de la canopée (feuilles vertes).
        \item \textbf{Consommateurs primaires (Herbivores) :} Chenilles (feuilles), Singes \emph{Cercopithecus} (fruits, feuilles, fleurs), Termites (bois mort, cellulose), Vers de terre (matière organique du sol -- détritivores).
        \item \textbf{Consommateurs secondaires :} Oiseaux insectivores (mangent chenilles), Fourmis prédatrices (mangent autres insectes/termites).
        \item \textbf{Consommateurs tertiaires :} Python (mange singes, oiseaux, rongeurs), Rapaces (non vus mais probables).
        \item \textbf{Décomposeurs :} Champignons (sur bois mort), Bactéries (dans le sol/humus, non visibles mais essentielles).
    \end{itemize}
    \item \textbf{Interactions (exemples) :}
    \begin{itemize}
        \item \emph{Prédation :} Python $\to$ mange Singe / Oiseau.
        \item \emph{Herbivorie :} Chenille $\to$ mange Feuille de \emph{Musanga}.
        \item \emph{Décomposition/Minéralisation :} Champignons + Bactéries $\to$ décomposent Bois mort / Humus $\to$ libèrent Sels minéraux $\to$ absorbés par Racines de \emph{Musanga}.
    \end{itemize}
    \item \textbf{Conclusion :} Cet ensemble constitue un écosystème forestier tropical humide : le biotope (climat, sol) supporte une biocénose diversifiée (producteurs, consommateurs 1/2/3, décomposeurs, détritivores) reliée par des flux de matière et d'énergie.
\end{enumerate}
\end{exoresolu}

\courssub{Les grands écosystèmes du Cameroun : Forêt et Savane}

Le Cameroun est souvent surnommé l'\textbf{« Afrique en miniature »} car il condense sur son triangle national (475 442 km²) la quasi-totalité des grands écosystèmes africains. Cette diversité exceptionnelle est dictée principalement par le \textbf{gradient climatique Nord-Sud}, notamment la \textbf{pluviométrie} (quantité de pluie annuelle) et la \textbf{saisonnalité}.

\begin{popfigure}
\begin{tikzpicture}[scale=0.85]
    % Carte simplifiée du Cameroun (contour approximatif)
    \def\cameroun{(-5,0) -- (-4.5,1.5) -- (-3,3) -- (-1,4) -- (1,4.5) -- (3,4) -- (4.5,2.5) -- (5,0.5) -- (4.5,-1) -- (3,-2.5) -- (1,-3.5) -- (-1,-4) -- (-3,-3.5) -- (-4.5,-2) -- cycle}
    % Fond de carte
    \draw[popline, thick] \cameroun;
    \clip \cameroun;
    % ZONE FORÊT (Sud) - Vert
    \fill[popgreen!30] (-5,-4) rectangle (5,0.5); 
    \fill[popgreen!40] (-5,-4) -- (-5,0.5) -- (0,2.5) -- (5,0.5) -- (5,-4) -- cycle; 
    % ZONE SAVANE (Nord) - Orange/Jaune
    \fill[poporange!30] (-5,0.5) rectangle (5,4.5); 
    \fill[poporange!40] (-5,0.5) -- (0,2.5) -- (5,0.5) -- (5,4.5) -- (-5,4.5) -- cycle; 
    % Ligne de transition (Adamaoua / 6-7°N)
    \draw[popdark, dashed, thick] (-5,0.5) .. controls (-2,1.5) and (2,1.5) .. (5,0.5) node[right, font=\small, popdark] {Frontière écologique $\approx$ 1500 mm/an};
    % Villes repères
    \node[fill=popdark, circle, inner sep=1.5pt, label={[font=\tiny, popdark]below:Yaoundé}] at (-1.5, 0.5) {};
    \node[fill=popdark, circle, inner sep=1.5pt, label={[font=\tiny, popdark]above:Douala}] at (-2.5, 1.5) {};
    \node[fill=popdark, circle, inner sep=1.5pt, label={[font=\tiny, popdark]above:Ngaoundéré}] at (0.5, 2.5) {};
    \node[fill=popdark, circle, inner sep=1.5pt, label={[font=\tiny, popdark]above:Maroua}] at (2.5, 3.5) {};
    \node[fill=popdark, circle, inner sep=1.5pt, label={[font=\tiny, popdark]left:Bertoua}] at (-3.5, 1) {};
    \node[fill=popdark, circle, inner sep=1.5pt, label={[font=\tiny, popdark]below:Kribi}] at (-2, -1) {};
    % Légendes zones
    \node[popgreen!70!black, font=\bfseries\large, align=center] at (-3, -2) {FORÊT DENSE\\ HUMIDE\\ (Sud, Est, Littoral)};
    \node[poporange!70!black, font=\bfseries\large, align=center] at (2.5, 3.5) {SAVANE\\ (Nord, Adamaoua,\\ Extrême-Nord)};
    % Pluviométrie indicative
    \node[popblue, font=\small, align=center, fill=popcream, rounded corners, inner sep=2pt] at (-3.5, 3) {Pluviométrie\\ > 2000 mm/an\\ (ex: Debundscha 10 000 mm)};
    \node[popblue, font=\small, align=center, fill=popcream, rounded corners, inner sep=2pt] at (3.5, -1.5) {Pluviométrie\\ 1000 -- 1500 mm/an\\ (ex: Ngaoundéré 1500 mm)};
    \node[popblue, font=\small, align=center, fill=popcream, rounded corners, inner sep=2pt] at (4, 4) {Pluviométrie\\ < 1000 mm/an\\ (ex: Maroua 700 mm)};
    % Fleche Nord
    \node at (0, 5) {\textbf{N}};
    \draw[->, thick] (0, 4.7) -- (0, 4.2);
\end{tikzpicture}
\legende{Carte simplifiée du Cameroun montrant la répartition des deux grands écosystèmes : la forêt dense humide au Sud et à l'Est (pluviométrie élevée $> 1500$ mm/an) et la savane au Nord et sur les hauts plateaux de l'Adamaoua (pluviométrie plus faible $< 1500$ mm/an, saison sèche marquée). La frontière est progressive (zone de transition forêt-savane). Villes repères : Yaoundé, Douala (forêt), Ngaoundéré (transition), Maroua (savane sahélienne).}
\end{popfigure}

\begin{definition}[Forêt dense humide et Savane : définitions climatiques]
\begin{itemize}
    \item \textbf{Forêt dense humide (Sud, Est, Littoral) :} Écosystème caractérisé par une \textbf{pluviométrie élevée} ($> 1500$ mm/an, souvent $> 2000$ mm), une \textbf{température moyenne élevée et stable} ($\approx 24-26$ °C), une \textbf{humidité relative forte} ($> 80\%$), une \textbf{saison sèche courte} (1 à 3 mois) ou absente (côte). Le biotope favorise une biomasse et une biodiversité maximales.
    \item \textbf{Savane (Nord, Adamaoua, Extrême-Nord) :} Écosystème caractérisé par une \textbf{pluviométrie plus faible et très saisonnière} ($500$ à $1500$ mm/an), une \textbf{saison sèche longue et marquée} (4 à 7 mois, harmattan), des \textbf{températures plus contrastées} (fortes amplitudes journalières/annuelles). Le biotope sélectionne des espèces adaptées à la sécheresse (xérophytes) et au feu.
\end{itemize}
\end{definition}

\begin{savaistu}[Le Cameroun : « Afrique en miniature » et hotspot de biodiversité]
Grâce à sa position géographique unique — charnière entre la \textbf{forêt congolaise} (au Sud) et la \textbf{savane soudanienne/sahélienne} (au Nord) — et à son relief varié (plaine côtière, plateau de l'Adamaoua $\approx 1100$ m, montagnes volcaniques : Mont Cameroun 4095 m, Monts Mandara), le Cameroun abrite :
\begin{itemize}
    \item Plus de \textbf{9000 espèces de plantes vasculaires} (dont de nombreux endémiques comme l'\emph{Ancistrocladus korupensis}).
    \item Environ \textbf{900 espèces d'oiseaux} (2ème d'Afrique après la RDC).
    \item Plus de \textbf{300 espèces de mammifères} (gorilles, chimpanzés, éléphants de forêt, girafes, lions, kob de Buffon).
    \item Des écosystèmes uniques : forêts de montagne (Kilum-Ijim), mangroves (Rio del Rey, Wouri), prairies d'altitude (Adamaoua), lacs de cratère (Nyos, Monoun, Barombi Mbo).
\end{itemize}
Cette richesse fait du Cameroun un \textbf{hotspot de biodiversité mondial} (région prioritaire pour la conservation). La protéger n'est pas seulement un devoir patriotique, c'est une responsabilité planétaire.
\end{savaistu}

\courssub{Exemples d'écosystèmes accessibles pour l'élève de 3ème}

Tu n'as pas besoin d'aller au Parc National de Waza ou de Korup pour étudier un écosystème. Ton environnement immédiat en regorge :

\begin{itemize}
    \item \textbf{La mare ou le point d'eau du quartier/village :} Écosystème aquatique doux, idéal pour observer producteurs (nénuphars, algues), consommateurs (têtards, insectes, grenouilles) et décomposeurs (vase). \emph{Attention : risque paludisme (larves anophèles), ne pas prélever sans protection.}
    \item \textbf{Le champ en jachère ou la friche de l'école :} Écosystème herbacé/buissonnant en succession secondaire. Producteurs : herbes (imperata), arbustes (chromolaena). Consommateurs : sauterelles, papillons, lézards, oiseaux granivores. Décomposeurs : termitières visibles, champignons après pluie.
    \item \textbf{Le jardin botanique / la cour arborée :} Écosystème forestier simplifié, planté par l'homme. Permet d'étudier la strate herbacée, arbustive, arborée, l'humus, la faune du sol.
    \item \textbf{Le parc national ou la réserve de faune le plus proche (Waza, Bénoué, Bouba Ndjida, Korup, Campo-Ma'an, Dja, Mbam-Djerem) :} Écosystèmes de référence, « témoins » peu perturbés. Sorties pédagogiques possibles (clubs nature, ONG partenaires).
\end{itemize}

\begin{experience}[Activité pratique : Inventaire biodiversité de la cour de récréation (ou jardin de l'école)]
\textbf{Objectif :} Appliquer la méthode d'identification des composants d'un écosystème anthropisé.
\textbf{Durée :} 1h30 (2 séances de 45 min).
\textbf{Matériel :} Quadrats de $1 \text{ m}^2$ (cadre bois/PVC), fiches d'inventaire, loupe, guide d'identification simplifié (arbres, herbes, insectes communs du Cameroun), appareil photo (téléphone).
\textbf{Protocole :}
\begin{enumerate}
    \item Définir 3 stations : \textbf{Zone herbeuse tondue}, \textbf{Pied d'arbre (ombre, litière)}, \textbf{Zone cimentée/bordure (milieu extrême)}.
    \item Poser le quadrat, relever : % couvert végétal, nombre d'espèces végétales différentes, liste animaux vus/piégés (piège Barber = gobelet enterré au ras du sol avec eau + savon 24h avant).
    \item Mesurer : T° sol, T° air, humidité sol (doigt/test tactile), luminosité (appli luxmètre).
\end{enumerate}
\textbf{Résultats attendus (type) :}
\begin{center}
\begin{tabularx}{\linewidth}{|l|X|X|X|}\hline
\textbf{Paramètre} & \textbf{Zone herbeuse (Soleil)} & \textbf{Pied d'arbre (Ombre/Litière)} & \textbf{Zone cimentée} \\ \hline
\textbf{Biotope} & Sec, chaud, lumineux, sol compact & Humide, frais, sombre, sol meuble (humus) & Très sec, très chaud, minéral, sans sol \\ \hline
\textbf{Producteurs} & Herbes (Cynodon, Imperata), Trèfle & Mousses, Fougères, Plantules d'arbres & Quasi nul (quelques algues vertes fentes) \\ \hline
\textbf{Consommateurs} & Fourmis, Sauterelles, Papillons, Lézards & Mille-pattes, Cloportes, Araignées, Petits coléoptères & Fourmis (nourriture humaine), Cafards \\ \hline
\textbf{Décomposeurs / Détritivores} & Peu visibles (bactéries sol) & \textbf{Abondants :} Champignons, Bactéries, Vers de terre, Termites & Absents \\ \hline
\textbf{Richesse spécifique} & Moyenne (5-10 sp. vég.) & \textbf{Élevée} (10-20 sp. vég. + faune sol) & Très faible \\ \hline
\end{tabularx}
\end{center}
\textbf{Interprétation :} La structure du biotope (humidité, lumière, sol) \textbf{commande} la composition de la biocénose. L'arbre crée un micro-biotope favorable (humus, ombre, humidité) qui abrite une biodiversité bien supérieure à la pelouse tondue ou au ciment. \textbf{L'homme (tonte, ciment) simplifie l'écosystème et appauvrit la biodiversité.}
\end{experience}

\begin{attention}[Erreurs fréquentes au BEPC - Vigilance !]
\begin{enumerate}
    \item \textbf{Confondre « Écosystème » et « Biotope » :} Dire « L'écosystème, c'est le sol et le climat. » $\to$ \textbf{Faux.} L'écosystème \textbf{inclut} les êtres vivants et leurs relations.
    \item \textbf{Oublier les décomposeurs :} Dans une chaîne alimentaire, ne jamais terminer par le super-prédateur. Il faut écrire : $\dots \to$ Super-prédateur $\xrightarrow{\text{mort}}$\textbf{Décomposeurs} $\to$ Sels minéraux $\to$ Producteurs. \textbf{Le cycle se ferme par les décomposeurs.}
    \item \textbf{Dire que la savane est « pauvre » :} La savane a une \textbf{biodiversité différente} (adaptée à la sécheresse/feu), pas « nulle ». Elle abrite des mégafaunes uniques (éléphants, girafes, lions, kob) absentes de la forêt dense.
    \item \textbf{Localiser la forêt au Nord :} La forêt est au \textbf{Sud} (et Est/Littoral). La savane est au \textbf{Nord} (et Adamaoua). Le gradient pluviométrique diminue du Sud vers le Nord.
    \item \textbf{Ne pas citer d'unités :} Pluviométrie en \textbf{mm/an}, Température en \textbf{°C}, Surface en \textbf{km²} ou \textbf{ha}.
    \item \textbf{Confondre Décomposeurs et Détritivores :} Les champignons et bactéries sont décomposeurs (microscopiques, chimiques). Vers de terre et termites sont détritivores (macroscopiques, mécaniques).
\end{enumerate}
\end{attention}

\courssub{Biodiversité et interdépendance dans la forêt et la savane}

\courssub{Structure et biodiversité de la forêt dense humide}

La forêt dense humide est l'écosystème le plus complexe et le plus riche en espèces. Sa structure verticale en \textbf{strates} crée une multitude de niches écologiques.

\begin{popfigure}
\begin{tikzpicture}[scale=0.7, every node/.style={font=\small}]
    % Strates
    \fill[popgreen!10] (-5,0) rectangle (5,0.5); % Sol
    \node[popgreen!70!black, font=\bfseries] at (-4.5, 0.2) {Sol / Humus};
    % Strate herbacée / moussue
    \fill[popgreen!20] (-5,0.5) rectangle (5,2);
    \foreach \x in {-4,-2,0,2,4} \draw[popgreen!70!black, thick] (\x,0.5) -- (\x+0.2,1.5) -- (\x-0.2,1.5) -- cycle; % Fougères
    \node[popgreen!70!black, font=\bfseries] at (-4.5, 1.2) {Strate herbacée (0-2m)};
    % Strate arbustive
    \fill[popgreen!30] (-5,2) rectangle (5,10);
    \foreach \x in {-3.5,-1,1.5,4} {
        \draw[popgreen!60!black, thick] (\x,2) -- (\x+0.5,6) -- (\x-0.5,6) -- cycle; % Arbres moyens
        \draw[popgreen!60!black, thick] (\x,6) ellipse ({0.8} and {1.5}); % Houppier
    }
    \node[popgreen!70!black, font=\bfseries] at (-4.5, 6) {Strate arbustive / moyenne (2-20m)};
    % Canopée / Emergents
    \foreach \x in {-4, 0, 3.5} {
        \draw[popgreen!50!black, very thick] (\x,10) -- (\x,18); % Tronc
        \draw[popgreen!50!black, very thick] (\x,18) ellipse ({2.5} and {2}); % Grande canopée
    }
    % Emergent
    \draw[popgreen!40!black, very thick] (4,10) -- (4,22);
    \draw[popgreen!40!black, very thick] (4,22) ellipse ({1.5} and {1.5});
    \node[popgreen!70!black, font=\bfseries] at (-4.5, 15) {Canopée (20-40m)};
    \node[popgreen!70!black, font=\bfseries] at (4.5, 20) {Émergents (>40m)};
    % Animaux indicateurs
    \node[poporange, font=\tiny] at (-2, 0.8) {Termites, Vers, Champignons};
    \node[poporange, font=\tiny] at (-2, 3) {Fourmis, Coléoptères, Petits mammifères};
    \node[poporange, font=\tiny] at (-2, 8) {Singes, Oiseaux, Serpents, Chauves-souris};
    \node[poporange, font=\tiny] at (-2, 16) {Aigles, Hornbills, Éperviers};
    % Flèches lumière/eau
    \draw[popgold, very thick, ->] (5.5, 20) -- (5.5, 0) node[midway, right, font=\tiny, popgold] {Lumière \searrow};
    \draw[popblue, very thick, ->] (-5.5, 0) -- (-5.5, 20) node[midway, left, font=\tiny, popblue] {Humidité \nearrow};
\end{tikzpicture}
\legende{Structure verticale type de la forêt dense humide camerounaise. La stratification crée des micro-habitats distincts (lumidité, hygrométrie, lumière) exploités par des communautés spécifiques. La canopée capte 95\% de la lumière.}
\end{popfigure}

\begin{aretenir}[Biodiversité forestière : records et endémisme]
\begin{itemize}
    \item \textbf{Richesse spécifique exceptionnelle :} 1 ha de forêt du Dja ou de Korup peut contenir $> 300$ espèces d'arbres (contre $< 20$ en forêt tempérée).
    \item \textbf{Endémisme fort :} Espèces uniques au Cameroun (ex. : \emph{Prunus africana} - Pygeum, \emph{Ancistrocladus korupensis} - anticancéreux).
    \item \textbf{Interdépendances spécialisées :} Pollinisation spécifique (figuier/guêpe), dispersion des graines par grands mammifères (éléphants, gorilles, chimpanzés) ou gros oiseaux (caloës, touracos). Disparition du disperseur $\to$ disparition de l'arbre.
\end{itemize}
\end{aretenir}

\courssub{Structure et biodiversité de la savane}

La savane est une formation herbacée continue (strate graminéenne) parsemée d'arbres et d'arbustes (strate ligneuse), soumise à une alternance saisonnière stricte.

\begin{popfigure}
\begin{tikzpicture}[scale=0.8, every node/.style={font=\small}]
    % Sol
    \fill[poporange!20] (-5,0) rectangle (5,0.5);
    \node[poporange!70!black, font=\bfseries] at (-4.5, 0.2) {Sol nu / Herbes sèches (saison sèche)};
    % Strate herbacée (haute en saison des pluies)
    \fill[popgreen!30!poporange!30] (-5,0.5) rectangle (5,2.5);
    \foreach \x in {-4.5,-3.5,...,4.5} \draw[popgreen!60!black, thick] (\x,0.5) -- (\x,2.2); % Herbes
    \node[popgreen!70!black, font=\bfseries] at (-4.5, 1.5) {Strate herbacée (0-2m) : Poacées (Andropogon, Hyparrhenia)};
    % Arbres épars (acacias, baobabs, karités)
    \foreach \x/\h/\w in {-3/5/1.5, 1/6/2, 4/4/1} {
        \draw[poporange!70!popdark, thick] (\x,0.5) -- (\x,0.5+\h); % Tronc
        \draw[popgreen!50!black, thick] (\x,0.5+\h) ellipse ({\w} and {1}); % Houppier plat
    }
    \node[poporange!70!popdark, font=\bfseries] at (-4.5, 4) {Strate ligneuse discontinue : Acacias, Baobabs, Karités, Nérés};
    % Faune typique
    \node[poporange, font=\small] at (3, 0.8) {Kob, Buffon, Gazelles};
    \node[poporange, font=\small] at (3, 1.5) {Phacochères, Damans};
    \node[poporange, font=\small] at (3, 3) {Lions, Guépards, Hyènes};
    \node[poporange, font=\small] at (3, 5) {Circaètes, Vautours};
    % Feu
    \draw[poporange, thick, decorate, decoration={snake, amplitude=1mm, segment length=3mm}] (-4, 3) -- (-2, 4) node[above, font=\tiny, poporange] {Feu de brousse};
    % Racines profondes
    \draw[poporange!70!popdark!50!black, dashed, thick] (-3, 0.5) -- (-3.5, -1.5) node[below, font=\tiny, poporange!70!popdark] {Racines pivotantes > 10m};
    \draw[poporange!70!popdark!50!black, dashed, thick] (1, 0.5) -- (1.5, -1.5);
\end{tikzpicture}
\legende{Structure de la savane soudanaise (Nord/Adamaoua). Strate herbacée dominante (graminées pérennes C4), strate ligneuse claire (arbres xérophytes à écorce épaisse, racines profondes). Adaptations au feu (écorce liégeuse, bourgeons protégés) et à la sécheresse (feuilles caduques, photosynthèse C4).}
\end{popfigure}

\begin{aretenir}[Adaptations clés de la savane]
\begin{itemize}
    \item \textbf{Végétation (Xérophytes / Pyrophytes) :} Feuilles caduques (perte en saison sèche), cuticule épaisse, photosynthèse C4 (efficace forte lumière/faible eau), racines pivotantes très profondes, écorce épaisse (protection feu), bourgeons souterrains ou protégés.
    \item \textbf{Faune :} Migrations saisonnières (kobs, éléphants), estivation/hibernation (escargots, amphibiens), régime alimentaire flexible, activité crépusculaire/nocturne.
    \item \textbf{Rôle du feu :} Facteur écologique naturel (foudre) ou anthropique. \textbf{Effets positifs :} minéralisation rapide (cendres), stimulation repoussé, maintien paysage ouvert, lutte contre tiques. \textbf{Effets négatifs (feux tardifs/intenses) :} destruction graines/jeunes plants, érosion, perte azote, mortalité faune.
\end{itemize}
\end{aretenir}

\courssub{Interdépendance des êtres vivants : Chaînes et Réseaux trophiques}

L'énergie solaire captée par les producteurs circule le long des \textbf{chaînes alimentaires}. La matière, elle, \textbf{cycle} grâce aux décomposeurs.

\begin{definition}[Chaîne alimentaire, Réseau trophique, Niveau trophique]
\begin{itemize}
    \item \textbf{Chaîne alimentaire :} Suite linéaire d'êtres vivants où chacun mange le précédent. Représentée par des flèches $\to$ signifiant « est mangé par » ou « donne de l'énergie/matière à ». Sens de la flèche = sens du transfert d'énergie/matière.
    \item \textbf{Niveau trophique :} Rang occupé dans la chaîne. Niveau 1 = Producteurs, Niveau 2 = Consommateurs primaires, Niveau 3 = Secondaires, etc.
    \item \textbf{Réseau trophique :} Ensemble des chaînes alimentaires interconnectées d'un écosystème (la réalité du terrain).
    \item \textbf{Pyramide des biomasses / de l'énergie :} La quantité d'énergie (ou biomasse) diminue drastiquement d'un niveau à l'autre (règle des \textbf{10\%} environ transférés, 90\% perdus en chaleur, respiration, déchets non mangés).
\end{itemize}
\end{definition}

\begin{exemplebox}[Réseau trophique simplifié : Savane de Bénoué (Saison des pluies)]
\begin{center}
\begin{tikzpicture}[scale=0.85, every node/.style={font=\small, align=center}]
    % Niveau 1 Producteurs
    \node[draw=popgreen, thick, fill=popgreen!10, rounded corners, align=center] (herb) at (0,0){Herbes \\ (Andropogon, Hyparrhenia)};
    \node[draw=popgreen, thick, fill=popgreen!10, rounded corners, align=center] (arb) at (-3,-1.5){Arbres/Arbustes \\ (Acacia, Baobab, Karité)};
    % Niveau 2 Consommateurs primaires
    \node[draw=poporange, thick, fill=poporange!10, rounded corners, align=center] (kob) at (-4, -3.5){Kob de Buffon \\ (Herbivore)};
    \node[draw=poporange, thick, fill=poporange!10, rounded corners, align=center] (phaco) at (0, -3.5){Phacochère \\ (Herbivore/Détritivore)};
    \node[draw=poporange, thick, fill=poporange!10, rounded corners, align=center] (saut) at (4, -3.5){Sauterelles \\ (Herbivores)};
    % Niveau 3 Secondaires
    \node[draw=poppink, thick, fill=poppink!10, rounded corners, align=center] (lion) at (-4, -5.5){Lion / Hyène \\ (Carnivores)};
    \node[draw=poppink, thick, fill=poppink!10, rounded corners, align=center] (circ) at (0, -5.5){Circaète \\ (Ophiophage)};
    \node[draw=poppink, thick, fill=poppink!10, rounded corners, align=center] (lez) at (4, -5.5){Lézards / Oiseaux \\ (Insectivores)};
    % Niveau 4 Tertiaires / Décomposeurs
    \node[draw=poppurple, thick, fill=poppurple!10, rounded corners, align=center] (dec) at (0, -7.5){Décomposeurs \\ (Bactéries, Champignons) \\ + Détritivores (Termites, Vers)};
    % Flèches
    \foreach \from/\to in {herb/kob, herb/phaco, herb/saut, arb/kob, arb/phaco, arb/saut, kob/lion, phaco/lion, saut/lez, saut/circ, lez/circ, lion/dec, circ/dec, lez/dec, kob/dec, phaco/dec, saut/dec, herb/dec, arb/dec}
        \draw[popdark, thick, ->] (\from) -- (\to);
    % Soleil
    \node[popgold, font=\bfseries\large] at (0, 1.2) {(soleil)\ Énergie Solaire};
    \draw[popgold, thick, ->] (0, 0.8) -- (herb);
    \draw[popgold, thick, ->] (-0.5, 0.8) -- (arb);
    % Légende niveaux
    \node[popgreen, font=\bfseries] at (-6, 0) {Niv. 1};
    \node[poporange, font=\bfseries] at (-6, -3.5) {Niv. 2};
    \node[poppink, font=\bfseries] at (-6, -5.5) {Niv. 3};
    \node[poppurple, font=\bfseries] at (-6, -7.5) {Recyclage};
\end{tikzpicture}
\end{center}
\legende{Réseau trophique de la savane de Bénoué. Les flèches indiquent le sens du transfert de matière et d'énergie. Notez la convergence vers les décomposeurs (recyclage) et la multiplicité des proies pour un prédateur (stabilité du réseau).}
\end{exemplebox}

\begin{propriete}[Règle des 10\% et efficacité de transfert]
Seule une fraction (environ \textbf{10\%} en moyenne, de 1\% à 20\%) de l'énergie (ou biomasse) d'un niveau trophique est transférée au niveau supérieur.
\begin{itemize}
    \item \textbf{Conséquence 1 :} Les chaînes sont courtes (rarement $> 4-5$ niveaux). Il faut une énorme biomasse de producteurs pour soutenir peu de super-prédateurs.
    \item \textbf{Conséquence 2 :} La biomasse des producteurs $>$ biomasse consommateurs primaires $>$ biomasse consommateurs secondaires... (Pyramide de biomasse généralement droite pour les écosystèmes terrestres).
    \item \textbf{Conséquence 3 :} Polluants persistants (métaux lourds, pesticides organochlorés) se \textbf{bioaccumulent} et se \textbf{biomagnifient} le long de la chaîne (concentration $\times 10$ à chaque niveau). Danger pour l'homme (sommet de chaîne).
\end{itemize}
\end{propriete}

\courssub{Les activités humaines destructrices des écosystèmes}

L'homme, par ses activités, modifie brutalement le biotope et décime la biocénose, rompant l'équilibre dynamique.

\begin{center}
\begin{tabularx}{\linewidth}{|l|X|X|}\hline
\rowcolor{popblueL}
\textbf{Activité} & \textbf{Mécanisme de destruction (Biotope $\to$ Biocénose)} & \textbf{Conséquences majeures au Cameroun} \\ \hline
\textbf{Déforestation} (Bois d'œuvre, agriculture itinérante sur brûlis, plantations palmier/cacao/hévéa, mines) & Destruction du biotope (sol nu $\to$ érosion, lessivage, hausse T°, baisse hygrométrie). Fragmentation (îlots forestiers isolés $\to$ effet de lisière, isolement populations). & Perte biodiversité (Gorilles, Chimpanzés, Éléphants de forêt), érosion hydrique (fleuves boueux : Sanaga, Nyong), perturbation régime des pluies locale, émission CO$_2$. \\ \hline
\textbf{Feux de brousse} (Chasse, pâturage, miel, malveillance, feux précoces/tardifs) & Destruction biocénose (mortalité directe faune lente, destruction graines/jeunes plants). Minéralisation brutale (perte azote volatilisé). Sol nu $\to$ érosion précoce. & Régression forêt (avancée savane), disparition espèces sensibles au feu, appauvrissement sols (perte matière organique), risques incendies habitations. \\ \hline
\textbf{Braconnage / Surpêche} (Viande de brousse, ivoire, écailles pangolin, filets maillants) & Prélèvement non durable : cibles espèces à reproduction lente (grands mammifères, gros poissons). Rupture chaînes trophiques (cascade trophique). & « Forêts vides » (syndrome forêt vide), extinction locale (Rhinocéros noir, Girafe de Kordofan), déséquilibre (pullulation proies $\to$ surpâturage). \\ \hline
\textbf{Urbanisation / Artificialisation} (Villes, routes, barrages, carrières) & Imperméabilisation sols, pollution (eau, air, bruit, lumière), fragmentation habitats, introduction espèces invasives (Jacinthe d'eau, Chromolaena). & Perte terres agricoles, crues urbaines (Yaoundé, Douala), maladies vectorielles, conflits homme-faune. \\ \hline
\textbf{Agriculture intensive / Pesticides} (Coton Nord, Banane/Ananas Sud, Intrants chimiques) & Pollution sols/eaux (nitrates, phosphates, résidus pesticides). Destruction auxiliaires (pollinisateurs, prédateurs ravageurs). Salinisation (irrigation mal gérée). & Eutrophisation lacs/rivières, résistance ravageurs, perte fertilité sols, risques sanitaires (cancers, perturbateurs endocriniens). \\ \hline
\end{tabularx}
\end{center}

\begin{savaistu}[Cas concret : Le Lac Ossa (Littoral) et la Jacinthe d'eau]
Le Lac Ossa (réserve de faune, 4000 ha) subit une invasion massive de \emph{Eichhornia crassipes} (Jacinthe d'eau) depuis 2016. Causes : Eutrophisation (apports azotés/phosphatés agriculture/élevages amont) + absence de prédateurs naturels.
\textbf{Conséquences :} Asphyxie du lac (baisse O$_2$ $\to$ mortalité poissons), blocage navigation/pêche, perte biodiversité (lamantins, tortues), prolifération moustiques.
\textbf{Lutte :} Lâchers de \emph{Neochetina} spp. (coléoptères brouteurs spécifiques, lutte biologique) + arrachage mécanique + sensibilisation riverains. Exemple d'interaction homme-écosystème et de solution fondée sur la nature.
\end{savaistu}

\courssub{Restauration et conservation de la biodiversité}

Face à la dégradation, deux stratégies complémentaires s'imposent : \textbf{protéger} ce qui reste et \textbf{restaurer} ce qui est perdu.

\begin{definition}[Aires protégées et statuts au Cameroun (Loi 94/01 forestière, Loi 2024/016 faune)]
\begin{itemize}
    \item \textbf{Protection intégrale (Zones interdites d'exploitation) :} \textbf{Parcs Nationaux} (Waza, Korup, Bouba Ndjida, Bénoué, Campo-Ma'an, Lobéké, Mbam-Djerem, Nki, Boumba Bek), \textbf{Réserves de Faune} (Dja, Santchou, Bassa), \textbf{Sanctuaires} (Mengamé - gorilles). Objectif : conservation écosystèmes naturels, recherche, éducation, écotourisme.
    \item \textbf{Gestion durable / Usage multiple :} \textbf{Zones d'Intérêt Cynégétique (ZIC)} (ZICGC gérées par communautés, ZICAD gérées par l'État), \textbf{Forêts de production} (UFA - Unités Forestières d'Aménagement) avec plans d'aménagement (PEA), \textbf{Forêts communautaires}.
    \item \textbf{Corridors écologiques :} Liaisons entre aires protégées (ex. corridor TRIDOM : Tri-National Dja-Odzala-Minkébé ; corridor Bénoué-Bouba Ndjida) pour flux géniques et migrations.
\end{itemize}
\end{definition}

\begin{aretenir}[Stratégies de conservation : In situ / Ex situ]
\begin{itemize}
    \item \textbf{In situ (sur place) :} Création/extension aires protégées, corridors, lutte anti-braconnage (écogardes, drones, SMART), gestion communautaire (COVAREF), restauration habitats (reboisement, lutte espèces invasives), paiements pour services environnementaux (PSE / REDD+).
    \item \textbf{Ex situ (hors lieu naturel) :} Banques de graines (IRAD, Herbier National), jardins botaniques (Limbe, Yaoundé), élevage conservatoire (ex. tortues, crocodiles), cryoconservation (ADN, gamètes), zoos éducatifs (Mvog-Betsi).
    \item \textbf{Approche écosystémique :} Intégrer les populations locales (droits d'usage, partage bénéfices), valorisation produits forestiers non ligneux (PFNL : Gnetum/okok, Irvingia/andok, Ricinodendron/ndjanssang), écotourisme communautaire.
\end{itemize}
\end{aretenir}

\courssub{TP N°10 : Pratique du reboisement}

\begin{experience}[TP N°10 : Pratique du reboisement (Restauration d'un écosystème dégradé)]
\textbf{Objectif :} Maîtriser les techniques de production de plants en pépinière et de mise en terre pour restaurer un écosystème forestier ou agroforestier.
\textbf{Contexte :} Lutte contre la déforestation, restauration des berges (protection cours d'eau), brise-vent, bois-énergie, agroforesterie (cacao/ombrage).
\textbf{Durée :} 2 séances (1h30 pépinière + 1h30 plantation) + suivi sur l'année.
\textbf{Matériel :}
\begin{itemize}
    \item \textbf{Pépinière :} Sacs pépinière (polyéthylène noir 15x25 cm ou 20x30 cm), substrat (terre de surface + fumure décomposée + sable 2:1:1), graines/pré-germés/plantules (essences locales : \emph{Acacia mangium}, \emph{Leucaena leucocephala}, \emph{Terminalia superba}, \emph{Irvingia gabonensis}, \emph{Ricinodendron heudelotii}, \emph{Persea americana} - avocatier), arrosoirs, ombrière (palmes/toile 50\% ombre), fongicide (optionnel).
    \item \textbf{Plantation :} Pioches, pelles, machettes, cordeau, piquets, seaux, eau, paillis (herbes sèches, feuilles), engrais organique (compost/fumier).
\end{itemize}

\textbf{Phase 1 : Travail en pépinière (Production des plants)}
\begin{enumerate}
    \item \textbf{Préparation substrat :} Tamiser, mélanger, humidifier. Remplir sacs (tasser légèrement, laisser 2 cm bord).
    \item \textbf{Semis / Repiquage :}
    \begin{itemize}
        \item \emph{Grosses graines (Irvingia, Ricinodendron) :} Trempage 24-48h, semis direct 1 graine/sac, profondeur 2-3 cm.
        \item \emph{Petites graines (Acacia, Leucaena, Terminalia) :} Semis en godets/caisses $\to$ repiquage au stade 2-4 feuilles vraies (1 plant/sac). \emph{Attention : ne pas casser la racine pivotante.}
    \end{itemize}
    \item \textbf{Entretien pépinière (3-6 mois selon essence) :} Arrosage quotidien (matin/soir), désherbage manuel, éclaircissage (garder le plus vigoureux), traitement fongicide préventif (damping-off), \textbf{endurcissement} (réduction arrosage, levée ombrière 2-3 semaines avant plantation pour habituer au plein soleil).
\end{enumerate}

\textbf{Phase 2 : Mise en terre (Plantation)}
\begin{enumerate}
    \item \textbf{Choix du site et préparation :} Zone dégradée, sol nu, érosif. Nettoyage (coupe herbes/arbustes, conservation arbres utiles). \textbf{Traçage} : alignement (cordeau) ou quinconce. Espacement : $3 \times 3$ m (bois d'œuvre), $4 \times 4$ m (agroforesterie), $2 \times 2$ m (bois-énergie/brise-vent).
    \item \textbf{Fosses de plantation :} Dimensions \textbf{40 $\times$ 40 $\times$ 40 cm} (standard) ou $50^3$ cm (sol pauvre). Séparer terre de surface (noire, fertile) et terre de fond (rouge, stérile).
    \item \textbf{Amendement :} Mélanger terre de surface + 5-10 kg compost/fumier bien décomposé + 100-200 g engrais NPK (optionnel, fond de fosse). Remettre mélange au fond.
    \item \textbf{Plantation :} Ôter sac plastique \textbf{délicatement} (couper le fond + fente latérale). Placer motte au centre, collet au niveau du sol. Reboucher avec terre de surface tassée par couches. Former \textbf{cuvette d'arrosage} (bourrelet circulaire 50 cm diam).
    \item \textbf{Arrosage de mise en place :} 10-20 L d'eau par plant (même s'il a plu). \textbf{Paillage} : 5-10 cm paille/feuilles autour du pied (éloigné du collet) $\to$ conserve humidité, limite herbes, enrichit sol.
    \item \textbf{Protection :} Tuteur (tige droite), grillage/épines si broutage (chèvres, vaches), pare-feu (coupe-feu 5 m large autour parcelle).
\end{enumerate}

\textbf{Phase 3 : Suivi et entretien (Indispensable la 1ère année)}
\begin{itemize}
    \item \textbf{Arrosage :} 1 fois/semaine saison sèche (si pas pluie) 1ère année.
    \item \textbf{Désherbage :} 3-4 fois/an (manuel, 1 m rayon pied). \emph{Interdit : herbicide au pied jeune.}
    \item \textbf{Remplacement (Reprise) :} Compter plants morts à 1 mois, remplacer immédiatement (même essence).
    \item \textbf{Élagage / Taille de formation :} À partir an 2-3 (bois d'œuvre) pour tige unique, pas de fourche basse.
\end{itemize}

\textbf{Interprétation / Bilan pédagogique :}
\begin{itemize}
    \item Le reboisement \textbf{recrée le biotope} (sol, humidité, ombre) $\to$ permet retour \textbf{biocénose} (faune, flore, micro-organismes).
    \item Choix \textbf{essences locales} (adaptées climat/sol, utiles populations) $>$ essences exotiques à croissance rapide (Eucalyptus, Pin) qui appauvrissent sol/eau et hébergent peu de biodiversité locale.
    \item Taux de réussite visé $> 80\%$ (exigence BEPC/Projet). Échec fréquent = mauvais endurcissement, plantation trop tardive (fin saison pluies), absence entretien, broutage.
\end{itemize}
\end{experience}

\begin{exercice}[Entraînement BEPC - Situation d'intégration]
\textbf{Document :} Extrait d'un rapport d'un élève de 3ème après une sortie au Parc National de la Bénoué (région du Nord, savane soudanaise).
\begin{quote}
« Nous sommes arrivés à 8h. Le ciel était dégagé, le thermomètre indiquait 28 °C. Le sol était rouge, dur, craquelé par endroits. L'herbe était jaune, sèche, haute de 1,5 m par endroits. Nous avons vu des troupeaux de kob de Buffon (antilopes) broutant l'herbe. Plus loin, un groupe de phacochères fouissait le sol avec leurs groins. Un circaète Jean-le-Blanc (rapace) tournoyait haut dans le ciel. Au sol, de grosses termitières dures comme de la pierre ponctuaient la plaine. Notre guide nous a dit qu'en mars, des feux contrôlés sont allumés par les gardes pour renouveler l'herbe. »
\end{quote}
\textbf{Questions :}
\begin{enumerate}
    \item Identifie dans le texte \textbf{trois facteurs abiotiques} du biotope et \textbf{trois espèces} de la biocénose.
    \item Classe les espèces citées en \textbf{Producteurs, Consommateurs primaires, Consommateurs secondaires/tertiaires, Décomposeurs/Détritivores}. Justifie chaque classement par une phrase du texte.
    \item Le guide parle de « feux contrôlés en mars ». Mars est-il en saison sèche ou des pluies au Nord ? Ce feu est-il un facteur abiotique ou biotique ? Quelle en est la conséquence sur l'écosystème (avantage/inconvénient) ?
    \item En utilisant la notion d'écosystème \textbf{ouvert}, explique pourquoi la pluie qui tombe en juin-juillet est indispensable au maintien de cette savane.
\end{enumerate}
\end{exercice}

\begin{corrige}[Exercice n°1 - Entraînement BEPC]
\begin{enumerate}
    \item \textbf{Facteurs abiotiques (Biotope) :} Température (28 °C), Nature du sol (rouge, dur, craquelé), État de l'herbe (sèche, jaune -- indicateur climat/saison), Luminosité (ciel dégagé). \textbf{Espèces (Biocénose) :} Herbe (Poacées), Kob de Buffon, Phacochère, Circaète Jean-le-Blanc, Termites (constructeurs de termitières).
    \item \textbf{Classement trophique :}
    \begin{itemize}
        \item \textbf{Producteur :} Herbe (Poacées). \emph{Justification :} « broutant l'herbe » $\to$ base de l'alimentation, plante chlorophyllienne.
        \item \textbf{Consommateurs primaires (Herbivores) :} Kob de Buffon (\emph{« broutant l'herbe »}), Phacochère (\emph{« fouissait le sol »} pour racines/tubercules/herbes).
        \item \textbf{Consommateur secondaire/tertiaire (Carnivore) :} Circaète Jean-le-Blanc (rapace, mange reptiles/petits mammifères).
        \item \textbf{Détritivores / Décomposeurs :} Termites (\emph{« termitières »} $\to$ dégradent cellulose bois mort/herbes sèches, enrichissent sol). Bactéries/Champignons (non cités mais implicites dans termitières/sol).
    \end{itemize}
    \item \textbf{Saison / Nature du feu :} Mars = \textbf{Saison sèche} (Nord Cameroun : saison sèche $\approx$ nov--avril). Le feu est un \textbf{facteur abiotique} (phénomène physique/chimique de combustion, même s'il est allumé par l'homme). \textbf{Conséquences :} Avantage : élimine herbes sèches, stimule repousse verte (pousse nutritive pour herbivores), limite invasion arbustive, recycle minéraux (cendres K, Ca, Mg). Inconvénient : risque destruction nids/oiseaux au sol, perte azote (volatilisation N$_2$, NO$_x$), érosion si sol nu avant pluies, émission CO$_2$/particules.
    \item \textbf{Écosystème ouvert \& Pluie :} La savane est un système ouvert. La pluie (juin-juillet) apporte \textbf{l'eau} (facteur limitant n°1 en saison sèche) et \textbf{les sels minéraux} (lessivage atmosphère, altération roches). Sans cet apport externe de matière (eau, $\text{NO}_3^-$, $\text{K}^+$, $\text{PO}_4^{3-}$), les producteurs (herbes) ne peuvent pas photosynthétiser $\to$ pas de matière organique $\to$ effondrement de la biocénose (herbivores meurent/partent $\to$ prédateurs partent). La pluie recharge aussi la nappe phréatique pour la saison sèche.
\end{enumerate}
\end{corrige}

\begin{exercice}[Synthèse de l'unité : Gestion durable d'un terroir villageois]
\textbf{Situation :} Dans un village de la zone de transition forêt-savane (région du Centre, près de Mbandjock), les habitants constatent : baisse des rendements champs (sol lessivé), tarissement source d'eau potable en saison sèche, disparition gibier, conflits éléphants/champs. Le chef de village demande un plan d'action basé sur les connaissances SVT.
\textbf{Travail demandé :} En t'appuyant sur les notions du chapitre, propose un plan d'action structuré en 3 axes : \textbf{1. Protection / Conservation}, \textbf{2. Restauration}, \textbf{3. Valorisation durable}. Pour chaque action, cite le principe écologique mobilisé (ex. : corridor, strate, cycle matière, pyramide énergie, espèce parapluie).
\end{exercice}

\begin{corrige}[Exercice n°2 - Synthèse]
\textbf{Axe 1 : Protection / Conservation}
\begin{itemize}
    \item Créer/officialiser une \textbf{forêt communautaire} ou \textbf{zone de chasse villageoise (ZICGC)} sur les reliquats forestiers restants. \emph{Principe :} Protection \emph{in situ} biotope/biocénose, connectivité (corridor) vers parc/réserve proche.
    \item Interdire feux tardifs, instaurer \textbf{feux précoces contrôlés} (déc-janv) par comités villageois formés. \emph{Principe :} Maîtrise facteur abiotique feu, stimulation herbacée, protection strates ligneuses.
    \item Surveillance anti-braconnage (brigades villageoises + écogardes). \emph{Principe :} Maintien pyramide trophique (espèces parapluies : éléphant, grand ongulés).
\end{itemize}
\textbf{Axe 2 : Restauration}
\begin{itemize}
    \item \textbf{Reboisement berges rivière/source} (essences ripicoles : \emph{Mitragyna}, \emph{Raphia}, \emph{Ficus}) sur 20-50 m largeur. \emph{Principe :} Restauration biotope (filtration eau, ombrage $\downarrow$ T°, stabilisation berges, corridor faune).
    \item \textbf{Agroforesterie dans champs cacao/manioc :} Intégrer arbres fertilitaires (\emph{Leucaena, Gliricidia, Acacia}) et fruitiers (\emph{Irvingia, Ricinodendron, Dacryodes/safoutier}). \emph{Principe :} Recyclage matière (litières $\to$ humus $\to$ sels), fixation azote, strates multiples $\uparrow$ biodiversité auxiliaires, ombrage cacao.
    \item \textbf{Couloirs de migration éléphants :} Laisser bandes non cultivées (1-2 km large) reliant blocs forestiers. \emph{Principe :} Connectivité paysagère, flux géniques, réduction conflits homme-faune.
\end{itemize}
\textbf{Axe 3 : Valorisation durable}
\begin{itemize}
    \item Domestication/valorisation \textbf{PFNL} (Okok/Gnetum, Andok/Irvingia, Njanssang/Ricinodendron, Moabi) $\to$ revenus sans coupe bois. \emph{Principe :} Valorisation biodiversité, incitation conservation.
    \item \textbf{Apiculture} (ruches kenyanes en lisière forêt/haies). \emph{Principe :} Service pollinisation (rendements cultures $\uparrow$), revenu miel/cire, dissuasion braconnage/feux (gardiens ruches).
    \item \textbf{Écotourisme communautaire} (observation éléphants, oiseaux, grottes, culture). \emph{Principe :} Valorisation économique écosystème intact $\to$ financement conservation.
\end{itemize}
\end{corrige}

\begin{aretenir}[Synthèse finale de la leçon : Ce qu'il faut maîtriser pour le BEPC]
\begin{enumerate}
    \item \textbf{Définition rigoureuse :} \textbf{Écosystème = Biotope (abiotique) + Biocénose (biotique) + Interactions.} Système \textbf{ouvert} (flux énergie/matière).
    \item \textbf{Composants biotiques :} \textbf{Producteurs} (base), \textbf{Consommateurs} (1°, 2°, 3°), \textbf{Décomposeurs} (bactéries/champignons = minéralisation), \textbf{Détritivores} (vers/termites = fragmentation). \textbf{Jamais oublier décomposeurs dans une chaîne !}
    \item \textbf{Cameroun :} \textbf{Forêt dense humide} au \textbf{Sud} (Pluie $> 1500$ mm, T° stable, strates complexes, biodiversité max, endémisme) $\leftrightarrow$ \textbf{Savane} au \textbf{Nord} (Pluie $< 1500$ mm, Saison sèche 4-7 mois, strate herbacée C4, arbres xérophytes, feu facteur écologique).
    \item \textbf{Interdépendance :} Chaînes/Réseaux trophiques, Règle des 10\% (pyramides), Bioaccumulation polluants, Cycles biogéochimiques (C, N, Eau) fermés par décomposeurs.
    \item \textbf{Menaces :} Déforestation, Feux destructifs, Braconnage, Artificialisation, Pollution/Intrants. Conséquences : Érosion, Perte biodiversité, Changement microclimat, Conflits, Appauvrissement populations.
    \item \textbf{Solutions :} \textbf{Aires protégées} (Parcs, Réserves, ZIC, Corridors), \textbf{Restauration} (Reboisement essences locales, Agroforesterie, Protection berges), \textbf{Gestion durable} (PFNL, Apiculture, Écotourisme, PSE/REDD+), \textbf{Législation} (Loi forestière, Loi faune, CITES).
    \item \textbf{Méthode BEPC :} Document $\to$ 1. Biotope (3 facteurs + unités), 2. Biocénose (espèces + groupes trophiques + justifications texte), 3. Interactions (proie-prédateur, décomposition, compétition), 4. Conclusion « C'est un écosystème car... ». \textbf{Unités obligatoires !}
\end{enumerate}
\end{aretenir}

\coursec{Biodiversité de la forêt et de la savane}

Dans la section précédente, nous avons découvert ce qu'est un \cle{écosystème} : un ensemble formé par un milieu de vie (le \cle{biotope}) et les êtres vivants qui l'occupent (la \cle{biocénose}). Nous avons aussi situé les grands écosystèmes du Cameroun : la \cle{forêt dense} au Sud et au Centre, la \cle{savane} au Nord et à l'Adamaoua. Mais ces deux milieux sont-ils peuplés de la même façon ? Abritent-ils les mêmes espèces, en même nombre ? Pour répondre à ces questions, il faut introduire une notion fondamentale : la \cle{biodiversité}.

\courssub{La notion de biodiversité}

\textbf{Observation.} Un élève qui se promène dans la forêt du Dja (région du Sud) note en une heure : des papillons, des singes, des oiseaux, des champignons, des arbres immenses, des lianes, des insectes... Dans la savane de Waza (Extrême-Nord), un autre élève note : des herbes hautes, quelques arbres isolés, des antilopes, des autruches, des lions. Les deux listes sont différentes : chaque milieu possède son propre « inventaire » d'espèces. C'est cet inventaire que les scientifiques appellent la biodiversité.

\begin{definition}[Biodiversité]
La \cle{biodiversité} est la \textbf{diversité des espèces vivantes} présentes dans un milieu donné. Elle se mesure notamment par le \textbf{nombre d'espèces différentes} (richesse spécifique) que l'on peut recenser dans cet écosystème : plantes, animaux, champignons et microorganismes.
\end{definition}

\begin{aretenir}[L'essentiel]
Plus un écosystème abrite d'espèces différentes, plus sa biodiversité est \textbf{riche}. On admet qu'une biodiversité riche assure l'\textbf{équilibre} et la \textbf{stabilité} de l'écosystème : chaque espèce y joue un rôle (nourriture, pollinisation, décomposition...), et la disparition de l'une d'elles fragilise tout l'ensemble.
\end{aretenir}

\begin{exemplebox}[Deux écosystèmes camerounais remarquables]
\begin{itemize}
\item Le \textbf{parc national de Waza} (Extrême-Nord) abrite plus de \textbf{300 espèces d'oiseaux}, ainsi que des éléphants, des lions, des girafes et des antilopes.
\item La \textbf{réserve de faune du Dja} (Sud) est classée au \textbf{patrimoine mondial de l'UNESCO} depuis 1987 : c'est l'une des forêts les plus riches d'Afrique, avec des gorilles, des chimpanzés et des arbres géants.
\end{itemize}
\end{exemplebox}

\courssub{La biodiversité de la forêt dense}

\textbf{Observation.} En forêt dense, on remarque que les plantes ne poussent pas toutes à la même hauteur : les unes rasent le sol, d'autres forment des buissons, d'autres encore montent jusqu'à 40 ou 50 mètres. La forêt est donc organisée en « couches » superposées : c'est la \cle{stratification}.

\begin{definition}[Stratification de la forêt]
La forêt dense est organisée en \textbf{étages verticaux} appelés \cle{strates}. Du bas vers le haut, on distingue :
\begin{enumerate}
\item la \textbf{strate herbacée} : herbes, fougères et jeunes pousses vivant dans la pénombre du sol forestier ;
\item la \textbf{strate arbustive} : arbustes et jeunes arbres (quelques mètres de haut) ;
\item la \textbf{canopée} : toit formé par les cimes des grands arbres soudées les unes aux autres (environ 20 à 40 m) ; c'est là que vivent la plupart des animaux de la forêt ;
\item les \textbf{émergents} : quelques arbres géants qui dépassent la canopée (jusqu'à 50--60 m).
\end{enumerate}
\end{definition}

\begin{popfigure}
\begin{center}
\begin{tikzpicture}[scale=0.92]
% sol
\fill[popgreenL!60] (0,0) rectangle (11,0.5);
\draw[popdark] (0,0.5) -- (11,0.5);
% strate herbacée
\foreach \x in {0.4,1.2,2.0,4.6,5.4,9.6,10.4}{
  \draw[popgreen,thick] (\x,0.5) -- (\x,1.1) (\x,0.5) -- (\x-0.18,0.95) (\x,0.5) -- (\x+0.18,0.95);}
% arbustes
\draw[popgreen!70!black,thick] (3.2,0.5) -- (3.2,2.0);
\fill[popgreen!60!black] (3.2,2.2) circle (0.55);
\draw[popgreen!70!black,thick] (8.6,0.5) -- (8.6,1.9);
\fill[popgreen!60!black] (8.6,2.1) circle (0.5);
% arbres canopée
\draw[popgold!60!black,line width=2pt] (1.8,0.5) -- (1.8,4.3);
\fill[popgreen!75!black] (1.8,4.6) circle (1.15);
\draw[popgold!60!black,line width=2pt] (5.2,0.5) -- (5.2,4.4);
\fill[popgreen!75!black] (5.2,4.7) circle (1.25);
\draw[popgold!60!black,line width=2pt] (7.2,0.5) -- (7.2,4.3);
\fill[popgreen!75!black] (7.2,4.6) circle (1.1);
% émergent
\draw[popgold!60!black,line width=2.5pt] (9.9,0.5) -- (9.9,6.0);
\fill[popgreen!85!black] (9.9,6.3) circle (0.9);
% liane
\draw[poppurple,thick,decorate,decoration={coil,segment length=6pt,amplitude=2.5pt}] (5.2,0.6) -- (5.2,4.2);
% singe dans canopée
\fill[popdark] (6.35,4.35) circle (0.14);
\draw[popdark,thick] (6.35,4.35) -- (6.6,4.15) (6.35,4.35) -- (6.15,4.6);
% oiseau sur émergent
\draw[poporange,thick] (10.6,7.0) .. controls (10.75,7.2) and (10.9,7.2) .. (11.05,7.0);
% délimitations d'étages
\draw[dashed,popmuted] (0,1.3) -- (11,1.3);
\draw[dashed,popmuted] (0,3.0) -- (11,3.0);
\draw[dashed,popmuted] (0,5.6) -- (11,5.6);
% légendes latérales
\node[anchor=west,font=\small] at (0.1,0.9) {\textbf{1. Strate herbacée} : fougères, herbes};
\node[anchor=west,font=\small] at (4.0,2.15) {\textbf{2. Strate arbustive} : arbustes, jeunes arbres};
\node[anchor=west,font=\small] at (0.1,3.35) {\textbf{3. Canopée} : cimes soudées (singes, oiseaux, insectes)};
\node[anchor=west,font=\small] at (6.3,6.9) {\textbf{4. Émergents} : arbres géants (iroko)};
\node[anchor=west,font=\small,poppurple] at (5.5,1.6) {liane};
\end{tikzpicture}
\end{center}
\legende{Schéma de la stratification verticale de la forêt dense. 1 : strate herbacée (fougères, herbes) ; 2 : strate arbustive (arbustes, jeunes arbres) ; 3 : canopée, toit des cimes où vivent singes, oiseaux et insectes ; 4 : émergents, arbres géants dépassant la canopée. Une liane (en violet) grimpe le long du tronc pour atteindre la lumière.}
\end{popfigure}

\textbf{Espèces caractéristiques de la forêt dense camerounaise.}
\begin{itemize}
\item \textbf{Végétaux} : l'\cle{iroko} et l'\cle{ébène} (arbres géants au bois précieux), les palmiers, les fougères arborescentes, les nombreuses lianes.
\item \textbf{Animaux} : les \cle{gorilles} et les \cle{chimpanzés} (grands singes), les \cle{éléphants de forêt} (plus petits que ceux de savane), les perroquets gris, les pangolins, les serpents et une immense variété d'insectes.
\end{itemize}

\begin{attention}[Ne pas confondre les étages]
La canopée n'est pas « le sommet d'un seul arbre » : c'est le \textbf{toit continu} formé par les cimes de \textbf{tous} les grands arbres. La plupart de la vie de la forêt (fruits, fleurs, singes, oiseaux) se concentre dans la canopée, et non au sol.
\end{attention}

\courssub{La biodiversité de la savane}

\textbf{Observation.} En savane, le paysage est très différent : une immense étendue d'herbes, parsemée çà et là d'arbres isolés. Pourquoi si peu d'arbres ? Parce que la savane connaît une \textbf{longue saison sèche} (plusieurs mois sans pluie) : seules les plantes capables de résister à la sécheresse y survivent.

\begin{definition}[La savane]
La \cle{savane} est un écosystème dominé par la \textbf{strate herbacée} (graminées hautes), avec des \textbf{arbres épars} (dispersés, isolés). Elle est adaptée à un climat marqué par une \textbf{saison sèche longue} et une saison des pluies courte.
\end{definition}

\textbf{Espèces caractéristiques de la savane camerounaise.}
\begin{itemize}
\item \textbf{Végétaux} : les graminées (herbes hautes), le \cle{karité}, le \cle{baobab}, le \cle{néré}.
\item \textbf{Animaux} : les grands \textbf{herbivores} (\cle{antilopes}, \cle{buffles}, éléphants de savane, girafes) et les \textbf{carnivores} qui les chassent (\cle{lions}, hyènes), ainsi que de très nombreux oiseaux (autruches, calaos, rapaces).
\end{itemize}

\begin{attention}[La savane n'est pas une forêt dégradée]
On entend parfois que la savane serait « une forêt qui a été détruite ». C'est \textbf{faux} : la savane est un \textbf{écosystème naturel à part entière}, adapté depuis très longtemps à une saison sèche longue. Ses plantes et ses animaux possèdent des adaptations propres (racines profondes, réserves d'eau, migrations) qui n'ont rien à voir avec une dégradation.
\end{attention}

\begin{savaistu}[Le baobab, une véritable citerne vivante]
Le \textbf{baobab} peut stocker \textbf{plusieurs milliers de litres d'eau} dans son tronc gonflé et spongieux. Grâce à cette réserve, il survit sans problème aux longs mois de sécheresse de la savane. Son écorce épaisse le protège aussi des feux de brousse. Certains baobabs vivent plus de mille ans !
\end{savaistu}

\courssub{Les adaptations des êtres vivants à leur milieu}

Chaque écosystème impose des contraintes : en forêt, la difficulté est le \textbf{manque de lumière} au sol (la canopée intercepte les rayons du soleil) ; en savane, c'est le \textbf{manque d'eau} pendant la saison sèche. Les êtres vivants possèdent des \cle{adaptations}, c'est-à-dire des caractéristiques qui leur permettent de survivre dans leur milieu.

\begin{exemplebox}[Quelques adaptations remarquables]
\begin{itemize}
\item \textbf{En forêt} : les \textbf{lianes} s'enroulent autour des troncs et grimpent \textbf{vers la lumière} sans avoir à fabriquer un tronc épais ; les plantes du sous-bois ont de \textbf{larges feuilles} pour capter le peu de lumière disponible.
\item \textbf{En savane} : les arbres comme le karité ont des \textbf{racines très profondes} qui puisent l'eau loin sous le sol ; leurs \textbf{feuilles sont réduites} (petites, parfois épineuses) pour limiter l'évaporation de l'eau ; le baobab \textbf{stocke l'eau} dans son tronc ; de nombreux animaux \textbf{migrent} vers les points d'eau pendant la saison sèche.
\end{itemize}
\end{exemplebox}

\begin{aretenir}[Idée clé]
La forme et le mode de vie d'une espèce reflètent les contraintes de son écosystème : \textbf{chercher la lumière} en forêt dense, \textbf{économiser l'eau} en savane.
\end{aretenir}

\courssub{Activité documentaire : comparer la biodiversité de deux milieux}

Pour comparer scientifiquement deux écosystèmes, on réalise des \textbf{relevés d'espèces} : on note, sur des surfaces comparables, toutes les espèces observées et leurs effectifs. Voici les résultats (fictifs mais réalistes) de deux relevés scolaires.

\begin{center}
\rowcolors{1}{popcream}{white}
\begin{tabularx}{\linewidth}{|l|X|X|}
\hline
\rowcolor{popblueL} \textbf{Groupe d'êtres vivants} & \textbf{Forêt du Dja (nombre d'espèces)} & \textbf{Savane de Waza (nombre d'espèces)} \\
\hline
Arbres et arbustes & 45 & 12 \\
Herbes & 18 & 30 \\
Mammifères & 32 & 25 \\
Oiseaux & 60 & 85 \\
Insectes & 120 & 70 \\
\hline
\rowcolor{popgreenL} \textbf{Total} & \textbf{275} & \textbf{222} \\
\hline
\end{tabularx}
\end{center}

\begin{popfigure}
\begin{center}
\begin{tikzpicture}
\begin{axis}[
  ybar, bar width=14pt,
  width=0.9\linewidth, height=6.5cm,
  symbolic x coords={Arbres, Herbes, Mammifères, Oiseaux, Insectes},
  xtick=data,
  ylabel={Nombre d'espèces relevées},
  ymin=0, ymax=140,
  legend style={at={(0.5,1.02)},anchor=south,legend columns=2,font=\small},
  nodes near coords, every node near coord/.append style={font=\scriptsize},
]
\addplot[fill=popgreen!70!black] coordinates {(Arbres,45) (Herbes,18) (Mammifères,32) (Oiseaux,60) (Insectes,120)};
\addplot[fill=popgold] coordinates {(Arbres,12) (Herbes,30) (Mammifères,25) (Oiseaux,85) (Insectes,70)};
\legend{Forêt du Dja, Savane de Waza}
\end{axis}
\end{tikzpicture}
\end{center}
\legende{Histogramme comparatif du nombre d'espèces relevées par groupe d'êtres vivants dans la forêt du Dja et la savane de Waza (données fictives réalistes). La forêt domine pour les arbres et les insectes ; la savane domine pour les herbes et les oiseaux.}
\end{popfigure}

\begin{methode}[Exploiter un tableau de relevés d'espèces]
\begin{enumerate}
\item \textbf{Lire} attentivement les colonnes et identifier ce qui est compté (ici : le nombre d'espèces par groupe).
\item \textbf{Calculer} les totaux par milieu en additionnant les effectifs de chaque colonne.
\item \textbf{Comparer} groupe par groupe : repérer où chaque milieu est le plus riche.
\item \textbf{Conclure} en reliant les différences observées aux conditions du milieu (lumière, eau, climat).
\end{enumerate}
\end{methode}

\begin{exoresolu}[Comparaison forêt du Dja / savane de Waza]
\textbf{Énoncé.} À partir du tableau de relevés ci-dessus : 1) calcule le nombre total d'espèces relevées dans chaque milieu ; 2) indique quel groupe d'êtres vivants est le plus riche dans chaque milieu ; 3) propose une explication.
\tcblower
\textbf{Solution détaillée.}
\begin{enumerate}
\item \textbf{Totaux.}
\begin{align*}
\text{Forêt du Dja} &: 45 + 18 + 32 + 60 + 120 = 275 \text{ espèces}\\
\text{Savane de Waza} &: 12 + 30 + 25 + 85 + 70 = 222 \text{ espèces}
\end{align*}
La forêt du Dja est donc, au total, plus riche en espèces que la savane de Waza ($275 > 222$).
\item \textbf{Groupes dominants.} Dans la forêt, ce sont les \textbf{insectes} (120 espèces) puis les oiseaux (60) ; dans la savane, ce sont les \textbf{oiseaux} (85 espèces) puis les insectes (70). La savane compte davantage d'herbes (30 contre 18) et la forêt beaucoup plus d'arbres (45 contre 12).
\item \textbf{Explication.} La forêt dense, humide et chaude toute l'année, offre de nombreux étages de vie (strates), donc de nombreuses « places » pour des espèces différentes : d'où sa grande richesse en arbres et en insectes. La savane, soumise à une longue saison sèche, favorise les herbes résistantes à la sécheresse et attire de grands herbivores suivis de nombreux oiseaux. Chaque milieu possède donc une biodiversité adaptée à ses conditions.
\end{enumerate}
\end{exoresolu}

\begin{attention}[Pièges fréquents]
\begin{itemize}
\item Ne confonds pas le \textbf{nombre d'espèces} et le \textbf{nombre d'individus} : 100 fourmis d'une seule espèce, c'est une biodiversité faible ; 10 insectes d'espèces différentes, c'est plus riche.
\item Ne conclus jamais qu'un milieu est « meilleur » qu'un autre : forêt et savane ont chacune une biodiversité \textbf{adaptée} à leur climat.
\end{itemize}
\end{attention}

\begin{aretenir}[Bilan de la section]
\begin{itemize}
\item La \textbf{biodiversité} est la diversité des espèces vivantes d'un milieu ; elle assure l'\textbf{équilibre} et la \textbf{stabilité} de l'écosystème.
\item La \textbf{forêt dense} est stratifiée (strate herbacée, strate arbustive, canopée, émergents) et très riche en espèces (iroko, ébène, gorilles, chimpanzés, éléphants de forêt).
\item La \textbf{savane} est dominée par les herbes, avec des arbres épars (karité, baobab, néré), de grands herbivores (antilopes, buffles) et des carnivores (lions).
\item Les êtres vivants présentent des \textbf{adaptations} à leur milieu : lianes grimpant vers la lumière en forêt ; racines profondes, feuilles réduites et réserves d'eau en savane.
\end{itemize}
\end{aretenir}

Maintenant que nous savons \emph{qui} vit dans la forêt et la savane, une question se pose : comment toutes ces espèces vivent-elles \emph{ensemble} ? C'est ce que nous verrons dans la section suivante, consacrée à l'\textbf{interdépendance des êtres vivants} dans un écosystème.

\coursec{Interdépendance des êtres vivants dans un écosystème}

Dans la section précédente, nous avons découvert l'extraordinaire richesse de la biodiversité de la forêt et de la savane camerounaises : des milliers d'espèces végétales et animales y cohabitent. Mais ces espèces ne vivent pas les unes à côté des autres par hasard. Pourquoi trouve-t-on des lions là où il y a des antilopes ? Pourquoi certaines forêts se régénèrent-elles mal lorsque les éléphants disparaissent ? Parce que tous les êtres vivants d'un écosystème sont \cle{liés les uns aux autres} : c'est ce que nous appelons l'\cle{interdépendance}. Dans cette section, nous allons découvrir ces liens, les classer, puis démontrer, à partir de faits observés, que la disparition d'une seule espèce peut déséquilibrer tout l'écosystème.

\courssub{Les relations alimentaires : chaînes et réseaux alimentaires}

\textbf{Observation.} Dans la savane du parc national de Waza (Extrême-Nord du Cameroun), un élève observateur note : les sauterelles mangent les herbes ; les antilopes broutent aussi les herbes ; les lions chassent les antilopes ; les vautours se nourrissent des restes des carcasses abandonnées par les lions. Comment organiser ces informations ?

\begin{definition}[Chaîne alimentaire]
Une \cle{chaîne alimentaire} est une suite d'êtres vivants dans laquelle chacun se nourrit de celui qui le précède. On la représente par des flèches orientées dans le \textbf{sens du transfert de matière et d'énergie}, c'est-à-dire de l'être mangé vers le mangeur.
\end{definition}

\begin{exemplebox}[Une chaîne alimentaire de la savane]
\[\text{Herbe} \longrightarrow \text{Antilope} \longrightarrow \text{Lion}\]
L'herbe est mangée par l'antilope ; l'antilope est mangée par le lion. La matière et l'énergie passent de l'herbe vers l'antilope, puis de l'antilope vers le lion.
\end{exemplebox}

\begin{methode}[Construire une chaîne alimentaire à partir d'observations]
\begin{enumerate}
\item Observer et noter \textbf{qui mange qui} (observations directes, restes de repas, contenus stomacaux, crottes).
\item Identifier le premier maillon : c'est toujours un \textbf{végétal} (producteur).
\item Placer ensuite les animaux dans l'ordre où ils se mangent.
\item Flécher dans le sens du \textbf{transfert de matière} : la flèche part de l'être \emph{mangé} et pointe vers l'être \emph{mangeur}. Lire la flèche « est mangé par ».
\item Vérifier : la chaîne commence par un végétal et chaque flèche doit pouvoir se lire « est mangé par ».
\end{enumerate}
\end{methode}

\begin{attention}[Le piège classique : le sens des flèches]
L'erreur la plus fréquente est d'inverser le sens des flèches, en fléchant « du mangeur vers le mangé ». Retiens bien : la flèche indique le \textbf{sens de circulation de la matière et de l'énergie}, donc elle va de la \textbf{proie vers le prédateur}. On écrit herbe $\rightarrow$ antilope $\rightarrow$ lion, et jamais lion $\rightarrow$ antilope $\rightarrow$ herbe. Astuce : la flèche « sort » de celui qui est mangé.
\end{attention}

Dans la réalité, un animal ne mange pas une seule espèce : le lion mange aussi des zèbres et des buffles, l'antilope mange plusieurs herbes, et le vautour récupère les restes de tous. Les chaînes alimentaires sont donc \textbf{entrecroisées}.

\begin{definition}[Réseau alimentaire]
Un \cle{réseau alimentaire} est l'ensemble des chaînes alimentaires d'un écosystème, reliées entre elles par des maillons communs.
\end{definition}

\begin{popfigure}
\begin{tikzpicture}[scale=1, every node/.style={font=\small}]
\% Producteurs
\node[draw, fill=popgreenL, rounded corners, minimum width=2.6cm, minimum height=0.8cm] (herbes) at (0,0) {\textbf{Herbes}};
\% Consommateurs primaires
\node[draw, fill=popgoldL, rounded corners, minimum width=2.4cm, minimum height=0.8cm] (saut) at (-3.2,2) {Sauterelles};
\node[draw, fill=popgoldL, rounded corners, minimum width=2.4cm, minimum height=0.8cm] (anti) at (3.2,2) {Antilopes};
\% Consommateurs secondaires
\node[draw, fill=poporangeL, rounded corners, minimum width=2.4cm, minimum height=0.8cm] (oiseau) at (-3.2,4) {Petits oiseaux};
\node[draw, fill=poporangeL, rounded corners, minimum width=2.4cm, minimum height=0.8cm] (lion) at (3.2,4) {\textbf{Lions}};
\% Charognards
\node[draw, fill=poppurpleL, rounded corners, minimum width=2.4cm, minimum height=0.8cm] (vaut) at (0,5.6) {Vautours};
\% Décomposeurs
\node[draw, fill=poppinkL, rounded corners, minimum width=3.4cm, minimum height=0.8cm] (decomp) at (0,-2.2) {Décomposeurs (bactéries, champignons)};
\% Flèches alimentaires
\draw[->, very thick, popgreen!70!black] (herbes) -- (saut);
\draw[->, very thick, popgreen!70!black] (herbes) -- (anti);
\draw[->, very thick, poporange] (saut) -- (oiseau);
\draw[->, very thick, poporange] (anti) -- (lion);
\draw[->, very thick, poppurple] (lion) -- (vaut);
\draw[->, very thick, poppurple] (oiseau) -- (vaut);
\% Retour vers décomposeurs
\draw[->, dashed, thick, poppink] (herbes.south) .. controls (-2,-1.2) .. (decomp.west);
\draw[->, dashed, thick, poppink] (vaut.north) .. controls (4,6.5) and (5,-1) .. (decomp.east);
\draw[->, dashed, thick, poppink] (lion.south) .. controls (2.2,1) .. (decomp.east);
\% Retour minéral
\draw[->, dashed, thick, popblue] (decomp.west) .. controls (-5.5,-1) and (-5.5,0.5) .. (herbes.west) node[midway, left, font=\footnotesize, text=popblue]{sels minéraux};
\end{tikzpicture}
\legende{Réseau alimentaire simplifié de la savane. Les flèches pleines indiquent le transfert de matière (qui est mangé par qui) ; les flèches en pointillés montrent le retour de la matière organique vers les décomposeurs, puis le recyclage en sels minéraux utilisés par les herbes.}
\end{popfigure}

\courssub{Les rôles des êtres vivants dans l'écosystème}

Chaque maillon d'une chaîne alimentaire joue un rôle précis, défini par sa façon de se nourrir.

\begin{definition}[Producteurs, consommateurs, décomposeurs]
\begin{itemize}
\item Les \cle{producteurs} sont les \textbf{végétaux verts} (herbes, arbustes, arbres). Grâce à la chlorophylle et à la lumière du Soleil, ils fabriquent eux-mêmes leur matière organique à partir de l'eau, du CO$_2$ et des sels minéraux du sol : ce sont des êtres \emph{autotrophes}.
\item Les \cle{consommateurs} sont les animaux qui mangent d'autres êtres vivants. On distingue :
      \begin{itemize}
      \item les \textbf{consommateurs primaires} (herbivores : sauterelle, antilope, zèbre), qui mangent les producteurs ;
      \item les \textbf{consommateurs secondaires} (carnivores : lion, petit oiseau insectivore, mangouste), qui mangent les consommateurs primaires ;
      \item les \textbf{consommateurs tertiaires} (superprédateurs : aigle, léopard), qui mangent les consommateurs secondaires.
      \end{itemize}
\item Les \cle{décomposeurs} (bactéries, champignons, mais aussi termites et vers de terre) se nourrissent des restes et des déchets de tous les êtres vivants morts. Ils transforment la matière organique en \cle{humus}, puis en sels minéraux qui enrichissent le sol et sont réutilisés par les producteurs.
\end{itemize}
\end{definition}

\begin{exemplebox}[Les termites, décomposeurs champions de la savane]
Les termites de la savane camerounaise décomposent l'herbe sèche, le bois mort et les bouses : ils recyclent ainsi d'énormes quantités de matière organique et enrichissent le sol en humus. De plus, ils servent eux-mêmes de nourriture à de nombreux animaux : fourmiliers, oiseaux, lézards, et même aux hommes dans certaines régions où les termites ailés sont consommés ! Les termites occupent donc \emph{deux places} dans le réseau : décomposeurs \emph{et} proies. Cela montre bien que les rôles sont interconnectés.
\end{exemplebox}

\begin{aretenir}[Le cycle de la matière dans l'écosystème]
Les producteurs fabriquent la matière organique ; les consommateurs la transfèrent le long des chaînes alimentaires ; les décomposeurs la recyclent en humus et en sels minéraux, qui reviennent aux producteurs. Rien ne se perd : la matière circule en \textbf{cycle} dans l'écosystème.
\end{aretenir}

\courssub{Les autres relations entre les êtres vivants}

Tous les liens ne sont pas alimentaires. Les espèces s'entraident aussi, ou s'exploitent mutuellement.

\begin{definition}[Relations d'entraide et d'exploitation]
\begin{itemize}
\item La \cle{pollinisation} : en butinant les fleurs pour leur nectar, les insectes (abeilles, papillons) transportent le pollen d'une fleur à l'autre, ce qui permet la reproduction des plantes. La plante nourrit l'insecte, l'insecte reproduit la plante : les deux y gagnent (mutualisme).
\item La \cle{dissémination des graines} : les animaux qui mangent des fruits (éléphants, singes, oiseaux, chauves-souris) rejettent les graines loin de l'arbre parent, dans leurs excréments, ce qui permet à de nouvelles plantes de s'installer ailleurs.
\item La \cle{symbiose} (au sens large, souvent mutualisme) est une association durable entre deux espèces qui y trouvent toutes les deux un avantage (exemple : les oiseaux pique-bœufs qui débarrassent les buffles de leurs parasites tout en se nourrissant).
\item Le \cle{parasitisme} est une relation où un être vivant (le parasite) vit aux dépens d'un autre (l'hôte) qu'il affaiblit sans le tuer immédiatement (exemples : les tiques sur les antilopes, les vers intestinaux).
\end{itemize}
\end{definition}

\courssub{Mise en évidence de l'interdépendance : études de cas}

\textbf{Étude de cas 1 : les éléphants et la régénération de la forêt.} Dans les forêts du Sud-Cameroun, certains arbres à gros fruits (comme l'espèce \emph{Irvingia gabonensis} ou le \emph{Balanites wilsoniana}) ne peuvent voir leurs graines germer efficacement que si elles ont traversé le tube digestif d'un éléphant : le passage dans l'intestin ramollit l'enveloppe de la graine (scarification), et les excréments déposés loin de l'arbre parent lui offrent un engrais naturel et un site de germination dégagé. Là où le braconnage a fait disparaître les éléphants, on observe que les graines de ces arbres restent sous l'arbre parent, germent mal, et que les jeunes plants se raréfient : la \textbf{régénération forestière est compromise}. La disparition d'une seule espèce animale menace donc la survie de plusieurs espèces végétales.

\textbf{Étude de cas 2 : suppression des prédateurs et explosion des herbivores.} Dans une réserve où les lions ont été éliminés par le braconnage, les effectifs d'antilopes ont été comptés chaque année. Les résultats sont reportés sur la courbe ci-dessous.

\begin{popfigure}
\begin{center}
\begin{tikzpicture}
\begin{axis}[
width=0.85\linewidth, height=6.5cm,
xlabel={Temps (années)}, ylabel={Effectif d'antilopes (milliers d'individus)},
xmin=0, xmax=10, ymin=0, ymax=14,
xtick={0,1,2,3,4,5,6,7,8,9,10},
ytick={0,2,4,6,8,10,12,14},
grid=both, grid style={popline!40},
legend style={at={(0.03,0.97)}, anchor=north west, font=\footnotesize}
]
\addplot[popcurve, mark=*, mark size=1.5pt] coordinates {
(0,4) (1,4.1) (2,4.2) (3,6.5) (4,9) (5,11.5) (6,12.8) (7,11) (8,8.5) (9,6.5) (10,5.5)
};
\addplot[dashed, popmuted, domain=0:10] {4.1};
\draw[poporange, dashed] (axis cs:2,0) -- (axis cs:2,14);
\node[font=\footnotesize, text=poporange, anchor=west] at (axis cs:2.1,13) {élimination des lions};
\end{axis}
\end{tikzpicture}
\end{center}
\legende{Évolution de l'effectif d'une population d'antilopes après l'élimination des lions (année 2). La ligne horizontale en pointillés indique l'effectif d'équilibre initial.}
\end{popfigure}

\textbf{Interprétation.} Avant l'année 2, l'effectif des antilopes est stable autour de \num{4000} individus : les lions prélèvent chaque année à peu près autant d'antilopes qu'il en naît. Après l'élimination des lions, les antilopes n'ont plus de prédateur : leur effectif \textbf{augmente fortement} jusqu'à environ \num{13000} individus (année 6). Mais ce surpeuplement épuise les ressources en herbes : la nourriture manque, les antilopes affaiblies meurent ou se reproduisent moins, et l'effectif \textbf{diminue} ensuite. L'écosystème a été profondément déséquilibré.

\begin{propriete}[Interdépendance des êtres vivants]
Les êtres vivants d'un écosystème sont \cle{interdépendants} : chacun dépend des autres pour se nourrir, se reproduire ou se développer. La disparition ou la diminution d'une seule espèce (un maillon du réseau) entraîne des modifications en cascade sur les effectifs des autres espèces et \textbf{déséquilibre l'ensemble de l'écosystème}.
\end{propriete}

\begin{exoresolu}[Construire et exploiter un réseau alimentaire]
Dans une savane, on observe les faits suivants : la sauterelle mange l'herbe ; le lézard mange la sauterelle ; la mangouste mange le lézard ; l'aigle mange la mangouste et la sauterelle ; l'herbe est aussi broutée par le zèbre, lui-même chassé par le lion.
\begin{enumerate}
\item Construis deux chaînes alimentaires à quatre maillons.
\item Indique le rôle de chaque maillon de la première chaîne.
\item Que risque-t-il d'arriver aux lézards et aux sauterelles si les mangoustes disparaissent ?
\end{enumerate}
\tcblower
\textbf{Solution.}
\begin{enumerate}
\item \textbf{Chaîne 1 :} Herbe $\rightarrow$ Sauterelle $\rightarrow$ Lézard $\rightarrow$ Mangouste. \\
      \textbf{Chaîne 2 :} Herbe $\rightarrow$ Sauterelle $\rightarrow$ Mangouste $\rightarrow$ Aigle. \\
      (On note que l'Aigle mange aussi directement la Sauterelle, ce qui crée une chaîne plus courte Herbe $\rightarrow$ Sauterelle $\rightarrow$ Aigle, mais la chaîne à 4 maillons via la Mangouste est bien présente).
\item Herbe : \textbf{Producteur} (autotrophe) ; Sauterelle : \textbf{Consommateur primaire} (herbivore) ; Lézard : \textbf{Consommateur secondaire} (carnivore mangeant un herbivore) ; Mangouste : \textbf{Consommateur tertiaire} (carnivore mangeant un consommateur secondaire).
\item Si les mangoustes disparaissent, les lézards, n'ayant plus de prédateur (la mangouste), voient leur effectif \textbf{augmenter} ; ils mangent alors davantage de sauterelles, dont l'effectif \textbf{diminue}. Cette diminution des sauterelles peut aussi priver les aigles d'une partie de leur nourriture (puisque l'aigle mange aussi les sauterelles). Toute la chaîne est perturbée : c'est l'interdépendance.
\end{enumerate}
\end{exoresolu}

\begin{attention}[Ne pas conclure trop vite]
Quand on supprime un maillon, l'effet n'est pas toujours immédiat ni simple : dans un \textbf{réseau}, un animal peut compenser en mangeant une autre proie (exemple : l'aigle mange encore les sauterelles si les mangoustes disparaissent). C'est pourquoi on raisonne toujours sur l'ensemble du réseau alimentaire, et jamais sur une seule chaîne isolée.
\end{attention}

\begin{aretenir}[Conclusion de la section]
Dans un écosystème, les êtres vivants sont unis par des relations alimentaires (chaînes et réseaux alimentaires) et par d'autres relations (pollinisation, dissémination, symbiose, parasitisme). Les producteurs fabriquent la matière, les consommateurs la transfèrent, les décomposeurs la recyclent en humus. Les études de cas (éléphants et forêt, lions et antilopes) le prouvent : \textbf{toute perturbation d'un maillon se répercute sur tout l'écosystème}. C'est pourquoi protéger un écosystème, c'est protéger \emph{toutes} ses espèces, même les plus discrètes.
\end{aretenir}

\coursec{Les activités humaines destructrices des écosystèmes}

Dans la section précédente, nous avons découvert que les êtres vivants d'un écosystème sont \cle{interdépendants} : chaque espèce, du plus petit insecte au plus grand éléphant, joue un rôle dans les chaînes alimentaires et dans l'équilibre du milieu. Cet équilibre, construit sur de très longues périodes, est fragile. Or, depuis plusieurs décennies, les activités de l'Homme le perturbent de plus en plus fortement. Au Cameroun, qui possède l'une des plus riches biodiversités d'Afrique, les forêts reculent, le gibier se raréfie, certaines sources tarissent. Qui n'a jamais vu, en saison sèche, un feu de brousse dévorer une colline entière ? Qui n'a jamais entendu un ancien du village dire : « avant, la forêt arrivait jusqu'ici » ? Cette section a pour but d'identifier les principales \cle{activités humaines destructrices} des écosystèmes, de comprendre leurs mécanismes et de mesurer leurs conséquences, afin de préparer la section suivante consacrée à la restauration et à la conservation.

\courssub{Le feu de brousse}

\begin{definition}[Feu de brousse]
Un \cle{feu de brousse} est un incendie de végétation (herbes, arbustes, jeunes arbres) allumé volontairement ou accidentellement en milieu naturel, le plus souvent en saison sèche, dans la savane ou en lisière de forêt.
\end{definition}

\textbf{Pourquoi allume-t-on des feux de brousse ?} Les motivations sont nombreuses : les agriculteurs brûlent les champs pour éliminer les résidus de récolte et fertiliser momentanément le sol avec les cendres ; les éleveurs brûlent les vieilles herbes pour provoquer la repousse de jeunes pousses tendres pour le bétail ; les chasseurs brûlent la brousse pour débusquer le gibier ; certains feux partent d'un mégot de cigarette ou d'un feu de cuisine mal surveillé.

\textbf{Que détruit réellement un feu de brousse ?} On pourrait croire qu'il ne brûle que les herbes sèches visibles. En réalité, les dégâts sont beaucoup plus profonds :

\begin{itemize}
\item la \cle{végétation} est détruite : herbes, arbustes, jeunes plants d'arbres, graines enfouies dans le sol (qui auraient permis la régénération naturelle) ;
\item les \cle{animaux} sont tués ou chassés : insectes, reptiles, rongeurs, oiseaux nicheurs au sol, œufs et larves ; les grands animaux fuient et perdent leur habitat ;
\item l'\cle{humus} est détruit : cette couche superficielle du sol, riche en matière organique en décomposition et en micro-organismes, est brûlée ; or c'est elle qui nourrit les plantes ;
\item le \cle{sol nu} est ensuite exposé à l'\cle{érosion} : sans végétation pour retenir la terre, les premières pluies emportent la couche fertile vers les rivières.
\end{itemize}

\begin{popfigure}
\begin{center}
\begin{tikzpicture}[
  cause/.style={rectangle, rounded corners=3pt, draw=popblue, fill=popblueL, text width=3.4cm, align=center, minimum height=1.1cm, font=\small},
  effet/.style={rectangle, rounded corners=3pt, draw=poporange, fill=poporangeL, text width=3.4cm, align=center, minimum height=1.1cm, font=\small},
  fleche/.style={-{Stealth[length=3mm]}, very thick, popdark}]
\node[cause] (feu) at (0,0) {\textbf{Feu de brousse} (destruction de la végétation)};
\node[effet] (humus) at (5.2,0) {\textbf{Perte de l'humus} et des micro-organismes du sol};
\node[effet] (erosion) at (10.4,0) {\textbf{Érosion} du sol nu par les pluies};
\node[effet, draw=poppink, fill=poppinkL] (pauvre) at (5.2,-2.6) {\textbf{Appauvrissement du sol} : baisse des rendements, désertification};
\draw[fleche] (feu) -- (humus);
\draw[fleche] (humus) -- (erosion);
\draw[fleche] (erosion.south) |- (pauvre.east);
\draw[fleche] (humus.south) -- (pauvre.north);
\end{tikzpicture}
\end{center}
\legende{Chaîne de causes à effets montrant comment un feu de brousse conduit à l'appauvrissement du sol : destruction de la végétation, perte de l'humus, érosion, puis baisse de fertilité.}
\end{popfigure}

\begin{attention}[Un feu de brousse ne brûle pas que les herbes !]
Erreur fréquente : croire qu'un feu de brousse est « sans conséquence » parce que l'herbe repousse après. En réalité, le feu tue aussi les \cle{graines} enfouies, les \cle{jeunes plants} d'arbres, les \cle{insectes} et les \cle{micro-organismes} du sol qui fabriquent l'humus. Répétés chaque année, les feux empêchent la savane de se régénérer et appauvrissent durablement le sol.
\end{attention}

\courssub{La déforestation}

\begin{definition}[Déforestation]
La \cle{déforestation} est la destruction des forêts par l'Homme, par coupe des arbres ou par le feu, pour utiliser le bois ou libérer des terres pour l'agriculture et l'habitat.
\end{definition}

\textbf{Quelles sont les causes de la déforestation au Cameroun ?}

\begin{itemize}
\item l'\cle{agriculture itinérante sur brûlis} : le paysan coupe et brûle une parcelle de forêt, la cultive deux ou trois ans, puis l'abandonne quand le sol s'épuise et en défriche une nouvelle ;
\item l'\cle{exploitation forestière} : coupe industrielle des grands arbres (ayous, sapelli, iroko) pour l'exportation du bois ;
\item la coupe de \cle{bois de chauffe} et de bois à charbon : dans beaucoup de villages et de villes, le bois reste la principale source d'énergie pour cuisiner ;
\item l'extension des \cle{plantations} industrielles (palmiers à huile, hévéas, bananiers) et des zones d'habitation.
\end{itemize}

\textbf{Quelles sont les conséquences de la déforestation ?}

\begin{itemize}
\item \cle{perte des habitats} : des milliers d'espèces animales et végétales perdent leur milieu de vie et disparaissent ;
\item \cle{érosion des sols} : sans arbres pour retenir la terre, les pluies lessivent les sols, qui deviennent stériles ; les rivières sont ensablées ;
\item \cle{diminution des pluies} : la forêt évapore beaucoup d'eau qui alimente les nuages ; sa destruction réduit les précipitations locales et favorise l'assèchement des sources ;
\item perturbation du \cle{climat} : les forêts stockent du carbone ; leur destruction aggrave le réchauffement climatique.
\end{itemize}

\begin{exemplebox}[Une perte qui se compte en dizaines de milliers d'hectares]
On estime que le Cameroun perd chaque année des \cle{dizaines de milliers d'hectares} de forêt (de l'ordre de \SI{100000}{ha} par an selon certaines périodes, soit environ la superficie de 140 000 terrains de football !). À cette vitesse, une forêt entière peut disparaître en quelques décennies si rien n'est fait. C'est pourquoi le reboisement (TP n°10) est si important.
\end{exemplebox}

\courssub{Exploitation d'un document : l'évolution de la surface forestière du Cameroun}

\begin{popfigure}
\begin{center}
\begin{tikzpicture}
\begin{axis}[
  width=0.85\linewidth, height=6.2cm,
  xlabel={Année}, ylabel={Superficie forestière (millions d'ha)},
  xmin=1968, xmax=2022, ymin=16, ymax=26,
  xtick={1970,1980,1990,2000,2010,2020},
  ytick={16,18,20,22,24,26},
  grid=major, grid style={popline!50},
  legend style={at={(0.03,0.05)}, anchor=south west, font=\small}]
\addplot[popcurve, mark=*] coordinates {
  (1970,24.5) (1980,23.2) (1990,22.0) (2000,20.8) (2010,19.8) (2020,19.0)};
\legend{Superficie forestière (valeurs approchées)}
\end{axis}
\end{tikzpicture}
\end{center}
\legende{Évolution approximative de la superficie forestière du Cameroun entre 1970 et 2020. La courbe décroissante traduit la déforestation.}
\end{popfigure}

\begin{methode}[Lire et interpréter une courbe d'évolution]
Pour exploiter une courbe comme celle-ci, procède en trois étapes :
\begin{enumerate}
\item \textbf{Repérer la tendance} : observer le sens général de la courbe (ici elle \emph{décroît} : la forêt recule) ;
\item \textbf{Calculer une variation} : lire deux valeurs sur les axes, calculer la différence, puis le pourcentage de variation : $\text{pourcentage de perte} = \dfrac{\text{valeur initiale} - \text{valeur finale}}{\text{valeur initiale}} \times 100$ ;
\item \textbf{Conclure} : relier le résultat à une cause (ici les activités humaines) et à une conséquence (perte de biodiversité, érosion).
\end{enumerate}
\end{methode}

\begin{exoresolu}[Calcul du pourcentage de forêt perdue entre 1970 et 2020]
Énoncé : d'après le document, la superficie forestière du Cameroun était d'environ 24,5 millions d'hectares en 1970 et d'environ 19,0 millions d'hectares en 2020. Calculer la perte de surface forestière en millions d'hectares, puis le pourcentage de forêt perdue sur cette période. Conclure.
\tcblower
Solution détaillée :
\begin{enumerate}
\item Perte de surface : $24{,}5 - 19{,}0 = 5{,}5$ millions d'hectares perdus en 50 ans.
\item Pourcentage de perte : $\dfrac{24{,}5 - 19{,}0}{24{,}5} \times 100 = \dfrac{5{,}5}{24{,}5} \times 100 \approx 22{,}4\,\%$.
\item Conclusion : en un demi-siècle, le Cameroun a perdu environ \textbf{un cinquième de sa forêt} (plus de 22 \%), soit en moyenne plus de \SI{100000}{ha} par an. Cette perte s'explique principalement par l'agriculture sur brûlis, l'exploitation forestière et la coupe de bois de chauffe. Elle entraîne la disparition des habitats, l'érosion des sols et la raréfaction des espèces.
\end{enumerate}
\end{exoresolu}

\courssub{Le braconnage}

\begin{definition}[Braconnage]
Le \cle{braconnage} est la chasse ou la capture illégale d'animaux sauvages, en dehors des périodes autorisées, dans des zones protégées, ou sur des espèces protégées par la loi.
\end{definition}

\textbf{Pourquoi braconne-t-on ?} Certains braconnent pour se nourrir (viande de brousse), mais beaucoup braconnent pour le \cle{commerce} : les éléphants sont tués pour leurs défenses en \cle{ivoire}, les pangolins pour leurs écailles, les perroquets gris pour être vendus comme animaux de compagnie, les grands singes pour leur viande.

\textbf{Quelles sont les conséquences du braconnage ?}

\begin{itemize}
\item \cle{disparition des espèces} : tués plus vite qu'ils ne se reproduisent, les animaux deviennent de plus en plus rares ; certaines espèces risquent l'extinction totale ;
\item \cle{rupture des chaînes alimentaires} : la disparition d'une espèce déséquilibre tout l'écosystème ; par exemple, si les éléphants disparaissent, les graines de nombreux arbres (qu'ils dispersaient dans leurs crottes) ne germent plus, et la forêt elle-même se modifie ;
\item appauvrissement de la \cle{biodiversité}, patrimoine irremplaçable du pays.
\end{itemize}

\begin{savaistu}[Le pangolin, champion malheureux du trafic]
Le \cle{pangolin}, ce petit mammifère inoffensif couvert d'écailles qui se roule en boule pour se défendre, est considéré comme \textbf{l'un des mammifères les plus trafiqués au monde}. Au Cameroun, il est très recherché pour sa viande et pour ses écailles, exportées illégalement vers l'Asie. Pourtant, le pangolin est utile : il mange des millions de fourmis et de termites chaque année, protégeant ainsi les cultures et les habitations !
\end{savaistu}

\courssub{Les autres menaces sur les écosystèmes}

Au-delà des feux, de la déforestation et du braconnage, d'autres activités humaines dégradent les écosystèmes :

\begin{itemize}
\item la \cle{pollution} : déchets plastiques jetés dans la nature, produits chimiques agricoles (engrais, pesticides) et huiles de vidange déversés dans les rivières, qui empoisonnent les animaux et contaminent l'eau ;
\item l'\cle{urbanisation} : l'extension des villes (Yaoundé, Douala) détruit les forêts et les zones humides, et fragmente les habitats ;
\item la \cle{surexploitation des ressources} : pêche excessive dans les rivières et les lagunes, prélèvement excessif de plantes médicinales, exploitation minière qui détruit le sol et pollue les cours d'eau.
\end{itemize}

\courssub{Des changements visibles dans notre vie quotidienne}

Ces destructions ne sont pas des problèmes lointains : elles se voient autour de nous. Les anciens constatent la \cle{raréfaction du gibier} : là où l'on chassait facilement il y a trente ans, il faut désormais marcher des heures sans rien trouver. Dans les zones déboisées, on observe l'\cle{assèchement des sources} et des ruisseaux, car sans forêt, l'eau des pluies ruisselle au lieu de s'infiltrer dans le sol. Les paysans constatent la \cle{baisse des rendements} : les sols érodés et privés d'humus produisent de moins en moins de maïs, de manioc ou d'arachides, ce qui pousse à défricher de nouvelles parcelles — un véritable cercle vicieux.

\begin{aretenir}[L'essentiel de la section]
\begin{itemize}
\item Les principales activités humaines destructrices sont : les \cle{feux de brousse}, la \cle{déforestation} (brûlis, exploitation forestière, bois de chauffe), le \cle{braconnage}, la pollution, l'urbanisation et la surexploitation.
\item Un feu de brousse détruit la végétation, les animaux et l'humus ; le sol nu est ensuite emporté par l'érosion.
\item La déforestation provoque la perte des habitats, l'érosion et la diminution des pluies ; le Cameroun a perdu plus de 20 \% de sa forêt en 50 ans.
\item Le braconnage fait disparaître des espèces (éléphants, pangolins, perroquets gris) et rompt les chaînes alimentaires.
\item Ces destructions se traduisent concrètement par la raréfaction du gibier, l'assèchement des sources et la baisse des rendements agricoles.
\end{itemize}
\end{aretenir}

Heureusement, ces dégradations ne sont pas une fatalité : il est possible d'agir pour restaurer les milieux et protéger les espèces. C'est précisément l'objet de la prochaine section, consacrée à la \cle{restauration} et à la \cle{conservation} de la biodiversité, que tu mettras en pratique lors du TP n°10 sur le reboisement.

\coursec{Restauration et conservation de la biodiversité}

Nous venons de voir que les activités humaines — feux de brousse, déforestation, braconnage — dégradent gravement la forêt et la savane camerounaises. Faut-il alors se résigner à voir disparaître nos écosystèmes ? Non. L'être humain qui détruit peut aussi \emph{réparer} et \emph{protéger}. C'est tout l'objet de cette section : comprendre comment on \cle{conserve} la biodiversité encore intacte et comment on \cle{restaure} les milieux déjà dégradés. Tu découvriras que ces actions reposent sur des principes scientifiques simples, mais exigent rigueur et persévérance.

\courssub{Principes de la conservation de la biodiversité}

\courssub{Qu'est-ce que conserver ?}

\begin{definition}[Conservation]
La \cle{conservation} de la biodiversité est l'ensemble des actions visant à \textbf{protéger les êtres vivants et leurs habitats} afin d'éviter la disparition des espèces et la dégradation des écosystèmes.
\end{definition}

L'idée directrice est simple : une espèce ne survit que si son \emph{habitat} (le milieu où elle vit, se nourrit et se reproduit) est intact. On ne peut pas sauver le gorille sans sauver la forêt dense où il vit ; on ne peut pas sauver l'éléphant sans préserver les grandes étendues de savane qu'il parcourt. La conservation repose donc sur \textbf{deux principes complémentaires} :

\begin{itemize}
\item \textbf{Protéger les habitats} : interdire ou limiter la coupe des arbres, les cultures et les feux dans certaines zones ;
\item \textbf{Protéger les espèces menacées} : interdire leur chasse, leur capture et leur commerce.
\end{itemize}

\begin{attention}[Conserver n'est pas figer]
Conserver un écosystème ne signifie pas l'empêcher de vivre : les êtres vivants continuent d'y naître, de s'y nourrir et d'y mourir naturellement. On protège le milieu contre les \emph{destructions humaines}, pas contre les équilibres naturels (prédation, vieillissement des arbres).
\end{attention}

\courssub{Les moyens de conservation : les aires protégées}

\begin{definition}[Aire protégée]
Une \cle{aire protégée} est un espace (terrestre ou aquatique) délimité par la loi, où les activités humaines destructrices (chasse, coupe du bois, agriculture) sont interdites ou strictement contrôlées afin de préserver la nature. Les principales aires protégées sont les \textbf{parcs nationaux} et les \textbf{réserves} (réserves de faune, réserves forestières).
\end{definition}

Le Cameroun possède un réseau remarquable d'aires protégées, réparties dans les grands écosystèmes du pays :

\begin{center}
\begin{tabularx}{\linewidth}{|l|X|X|}
\hline
\rowcolor{popblueL} \textbf{Aire protégée} & \textbf{Écosystème} & \textbf{Espèces emblématiques protégées} \\
\hline
Parc national de Waza (Extrême-Nord) & Savane soudanienne & Éléphants, lions, girafes, antilopes, oiseaux aquatiques \\
\hline
Parc national de la Bénoué (Nord) & Savane & Hippopotames, buffles, lions, élands de Derby \\
\hline
Réserve de faune du Dja (Sud) & Forêt dense humide & Gorilles, chimpanzés, éléphants de forêt \\
\hline
Parc national de Korup (Sud-Ouest) & Forêt dense humide & Chimpanzés, léopards, arbres et plantes très diversifiés \\
\hline
\end{tabularx}
\end{center}

\begin{exemplebox}[La réserve de faune du Dja]
Créée en 1950 et classée \textbf{patrimoine mondial de l'UNESCO} depuis 1987, la réserve de faune du Dja (région du Sud) couvre environ \SI{526000}{ha} de forêt dense presque intacte, entourée par le fleuve Dja qui forme une barrière naturelle. Elle abrite l'une des plus grandes populations de \textbf{gorilles} et de \textbf{chimpanzés} d'Afrique centrale, ainsi que des éléphants de forêt et des centaines d'espèces d'oiseaux. C'est un exemple de conservation \emph{réussie} grâce à la protection de tout un habitat.
\end{exemplebox}

\begin{popfigure}
\begin{tikzpicture}[scale=1.0]
% Contour simplifié du Cameroun
\draw[thick, popdark, fill=popcream] (0,0) -- (0.3,1.6) -- (0.1,3.2) -- (0.6,4.6) -- (1.0,6.4) -- (1.6,7.6) -- (2.4,8.2) -- (3.2,7.8) -- (3.8,6.6) -- (4.6,5.8) -- (5.2,4.6) -- (5.6,3.4) -- (5.2,2.2) -- (4.4,1.2) -- (3.0,0.4) -- (1.4,0.0) -- cycle;
% Océan
\node[popblue] at (-0.9,0.6) {\small Océan};
\node[popblue] at (-0.9,0.2) {\small Atlantique};
% Aires protégées
\fill[popgreen!70] (1.9,7.3) circle (0.45);
\node[popdark, font=\small\bfseries] at (1.9,7.3) {1};
\fill[popgreen!70] (3.4,5.6) circle (0.5);
\node[popdark, font=\small\bfseries] at (3.4,5.6) {2};
\fill[popgreen!70] (4.3,2.4) circle (0.55);
\node[popdark, font=\small\bfseries] at (4.3,2.4) {3};
\fill[popgreen!70] (0.9,1.1) circle (0.4);
\node[popdark, font=\small\bfseries] at (0.9,1.1) {4};
% Limites des grands écosystèmes
\draw[dashed, popmuted] (0.35,4.4) -- (4.7,5.0);
\draw[dashed, popmuted] (0.15,2.6) -- (5.35,2.9);
\node[popmuted, font=\footnotesize\itshape] at (4.6,6.6) {Savane soudanienne};
\node[popmuted, font=\footnotesize\itshape] at (4.7,3.9) {Savane / transition};
\node[popmuted, font=\footnotesize\itshape] at (3.0,1.4) {Forêt dense humide};
\legende{Carte simplifiée des principales aires protégées du Cameroun : \textbf{1} = parc national de Waza ; \textbf{2} = parc national de la Bénoué ; \textbf{3} = réserve de faune du Dja ; \textbf{4} = parc national de Korup. Les traits pointillés délimitent schématiquement les grands écosystèmes.}
\end{tikzpicture}
\end{popfigure}

Outre les aires protégées, la conservation s'appuie sur :

\begin{itemize}
\item \textbf{Des lois contre le braconnage} : la loi camerounaise interdit la chasse des espèces protégées (gorilles, chimpanzés, éléphants, panthères...). Le braconnier encourt des amendes et des peines de prison. Des éco-gardes patrouillent dans les parcs ;
\item \textbf{La sensibilisation des populations} : émissions de radio, campagnes d'affichage, éducation environnementale à l'école, pour faire comprendre que la faune et la forêt sont un patrimoine à transmettre ;
\item \textbf{La répression du commerce illégal} : saisie de l'ivoire, des peaux et des animaux vivants vendus sur les marchés.
\end{itemize}

\begin{savaistu}[Une journée pour planter des arbres]
Au Cameroun, il existe une \textbf{journée nationale de l'arbre}, célébrée chaque année pour encourager le reboisement. Ce jour-là, élèves, enseignants, administrations et villageois plantent des arbres dans les écoles, les villes et les villages. C'est une occasion concrète pour toi de participer à la restauration de ton environnement !
\end{savaistu}

\courssub{La restauration des écosystèmes dégradés}

Conserver, c'est protéger ce qui existe encore. Mais que faire quand un milieu est \emph{déjà détruit} : une colline chauve, une savane brûlée chaque année, une forêt transformée en champ abandonné ? Il faut alors \emph{restaurer}.

\begin{definition}[Restauration et reboisement]
La \cle{restauration} est l'ensemble des actions visant à \textbf{réparer un écosystème dégradé} pour qu'il retrouve peu à peu sa biodiversité et ses fonctions (sol fertile, ombrage, habitat des animaux). Le \cle{reboisement} est la principale technique de restauration : il consiste à \textbf{planter des arbres} sur un terrain dénudé ou dégradé.
\end{definition}

\courssub{Les techniques de restauration}

\textbf{a) Le reboisement.} On plante de jeunes plants (issus de pépinières) en choisissant de préférence des \emph{espèces locales} adaptées au milieu. Les arbres plantés retiennent le sol par leurs racines, freinent l'érosion, créent de l'ombre et de l'humus, ce qui permet ensuite à d'autres plantes et animaux de revenir.

\textbf{b) La régénération naturelle.} Quand la dégradation n'est pas trop avancée, il suffit parfois de \emph{protéger} le terrain (contre les feux, le pâturage excessif, la coupe) pour que la végétation repousse \textbf{toute seule} à partir des graines et des souches présentes dans le sol. Cette technique, moins coûteuse, montre la capacité naturelle des écosystèmes à se reconstituer si on les laisse tranquilles.

\textbf{c) La lutte contre les feux de brousse.} Les feux tardifs (en pleine saison sèche) détruisent jeunes arbres, graines et animaux. Deux moyens simples permettent de s'en protéger :

\begin{itemize}
\item Les \textbf{pare-feu} : bandes de terrain nettoyées (herbes et broussailles enlevées) qui coupent la progression du feu autour des villages, des champs et des zones reboisées ;
\item Les \textbf{feux précoces} : brûler l'herbe \emph{en début de saison sèche}, quand elle est encore un peu verte : le feu est alors moins violent, n'atteint pas la cime des arbres et épargne la faune qui a le temps de fuir.
\end{itemize}

\textbf{d) L'agroforesterie (complément utile).} L'\cle{agroforesterie} consiste à \textbf{associer des arbres et des cultures} sur une même parcelle : par exemple, conserver ou planter des arbres dans une cacaoyère ou un champ de manioc. Les arbres protègent le sol du soleil et de la pluie battante, enrichissent la terre avec leurs feuilles mortes et fournissent fruits et bois. C'est une manière de cultiver \emph{sans détruire} l'écosystème.

\courssub{Expérience : le suivi des jeunes plants, condition de la réussite}

Reboiser, est-ce seulement mettre des plants en terre ? Pour le vérifier, des élèves de 3\textsuperscript{ème} ont mené l'expérience suivante lors d'une campagne de reboisement de la colline de leur quartier.

\begin{experience}[Survie des plants selon les soins apportés]
\textbf{Protocole.} 400 jeunes plants d'espèces locales sont répartis en 4 lots de 100 plants, plantés le même jour sur le même terrain :
\begin{itemize}
\item \textbf{Lot A} : plants simplement mis en terre, sans aucun soin ;
\item \textbf{Lot B} : plants arrosés régulièrement pendant la saison sèche ;
\item \textbf{Lot C} : plants protégés par un grillage contre les animaux et les feux ;
\item \textbf{Lot D} : plants arrosés \emph{et} protégés.
\end{itemize}
\textbf{Résultats.} Nombre de plants encore vivants après 2 ans :
\begin{center}
\begin{tabular}{|l|c|c|}
\hline
\rowcolor{popblueL} \textbf{Lot} & \textbf{Plants vivants / 100} & \textbf{Taux de survie} \\
\hline
A : sans soin & 25 & 25 \% \\
B : arrosage seul & 55 & 55 \% \\
C : protection seule & 60 & 60 \% \\
D : arrosage + protection & 85 & 85 \% \\
\hline
\end{tabular}
\end{center}
\textbf{Interprétation.} Le lot A montre que la plantation seule donne des résultats médiocres : les 3/4 des plants meurent (sécheresse, piétinement, feux). L'arrosage (B) et la protection (C) améliorent chacun la survie, et leur combinaison (D) donne les meilleurs résultats. \textbf{Conclusion :} la réussite d'un reboisement dépend autant du \emph{suivi} des plants (arrosage, protection) que de la plantation elle-même.
\end{experience}

\begin{popfigure}
\begin{tikzpicture}
\begin{axis}[
ybar, bar width=0.9cm,
width=0.9\linewidth, height=6.2cm,
ylabel={Taux de survie (\%)},
xlabel={Lot de plants},
symbolic x coords={A, B, C, D},
xtick=data,
ymin=0, ymax=100,
ytick={0,20,40,60,80,100},
nodes near coords,
nodes near coords align={vertical},
every node near coord/.append style={font=\small\bfseries},
]
\addplot[fill=popgreen!70, draw=popdark] coordinates {(A,25) (B,55) (C,60) (D,85)};
\end{axis}
\end{tikzpicture}
\legende{Taux de survie des plants après 2 ans selon les soins apportés : A = sans soin ; B = arrosage seul ; C = protection seule ; D = arrosage + protection. Le suivi multiplie par plus de 3 le taux de survie.}
\end{popfigure}

\begin{attention}[Reboiser ne suffit pas !]
Erreur fréquente : croire qu'un reboisement est réussi dès que les plants sont en terre. En réalité, un jeune plant est fragile : il faut le \textbf{protéger et le suivre pendant plusieurs années} (arrosage en saison sèche, désherbage, pare-feu, protection contre les animaux). Sans suivi, la plupart des plants meurent et l'argent investi est perdu. Au BEPC, précise toujours le \emph{suivi} quand tu proposes le reboisement comme solution.
\end{attention}

\courssub{Le rôle du citoyen et de l'élève}

La conservation n'est pas l'affaire des seuls éco-gardes et ministères : chaque citoyen, et chaque élève, peut agir à son échelle :

\begin{itemize}
\item \textbf{Planter des arbres} : dans la cour de l'école, de la maison, le jour de la journée nationale de l'arbre ;
\item \textbf{Refuser la viande de brousse d'espèces protégées} (gorille, chimpanzé, pangolin...) : sans acheteurs, le braconnage cesse d'être rentable ;
\item \textbf{Éviter les feux de brousse incontrôlés} et participer à la construction des pare-feu du village ;
\item \textbf{Ne pas jeter les déchets} (sachets plastiques) dans la nature ;
\item \textbf{Sensibiliser} son entourage : expliquer à sa famille et à ses voisins pourquoi la forêt et la savane sont indispensables.
\end{itemize}

\begin{methode}[Proposer et justifier une solution de conservation]
Face à un problème environnemental décrit dans un exercice, procède en 4 étapes :
\begin{enumerate}
\item \textbf{Identifier le problème} et sa cause (ex. : disparition des éléphants $\rightarrow$ braconnage ; colline érodée $\rightarrow$ déforestation) ;
\item \textbf{Choisir la solution adaptée} : protéger (aire protégée, loi, sensibilisation) si le milieu est menacé, restaurer (reboisement, pare-feu) s'il est déjà dégradé ;
\item \textbf{Préciser les modalités pratiques} : qui agit, comment, avec quel suivi ;
\item \textbf{Justifier} : expliquer \emph{pourquoi} cette solution répond à la cause du problème (ex. : le reboisement restaure le couvert végétal, ce qui freine l'érosion).
\end{enumerate}
\end{methode}

\begin{exoresolu}[Sauver la colline de Nkolbisson]
\textbf{Énoncé.} Près de Yaoundé, une colline autrefois boisée a été entièrement défrichée pour les cultures. Aujourd'hui, les pluies emportent la terre fertile, les rendements s'effondrent et la source au pied de la colline tarit en saison sèche. 1) Identifie la cause du problème. 2) Propose deux solutions de restauration et justifie-les. 3) Précise les conditions de réussite du reboisement.
\tcblower
\textbf{Solution.}
1) \textbf{Cause} : la déforestation de la colline. Sans arbres, les racines ne retiennent plus le sol (d'où l'érosion) et l'eau de pluie ruisselle au lieu de s'infiltrer (d'où le tarissement de la source).

2) \textbf{Solutions} :
\begin{itemize}
\item \emph{Reboiser la colline} avec des espèces locales : les racines fixeront le sol, les feuilles mortes formeront de l'humus et l'eau s'infiltrera à nouveau, alimentant la source ;
\item \emph{Pratiquer l'agroforesterie} en bas de la colline (arbres associés aux cultures) et laisser la régénération naturelle agir sur les zones protégées : on cultive tout en préservant le sol.
\end{itemize}

3) \textbf{Conditions de réussite} : choisir des espèces adaptées au milieu, protéger les jeunes plants (grillage, pare-feu contre les feux de brousse), les arroser pendant la saison sèche et assurer ce suivi \emph{pendant plusieurs années}, car l'expérience montre que sans soins, environ 75 \% des plants meurent.
\end{exoresolu}

\begin{exercice}[À toi de jouer]
Dans la région de l'Extrême-Nord, les habitants d'un village constatent que les feux de brousse tardifs détruisent chaque année les jeunes arbres et font fuir le gibier. Propose deux mesures concrètes pour limiter ces dégâts et justifie chacune d'elles. (Correction : voir la boîte « corrigé » du cahier d'exercices de l'unité.)
\end{exercice}

\begin{aretenir}[L'essentiel de la section]
\begin{itemize}
\item La \cle{conservation} protège les habitats et les espèces menacées : aires protégées (Waza, Bénoué, Dja, Korup), lois contre le braconnage, sensibilisation ;
\item La \cle{restauration} répare les milieux dégradés : reboisement, régénération naturelle, pare-feu et feux précoces, agroforesterie ;
\item Un reboisement ne réussit qu'avec un \textbf{suivi prolongé} des plants (arrosage, protection) : 85 \% de survie avec soins contre 25 \% sans soins ;
\item Chaque citoyen peut agir : planter des arbres, refuser la viande de brousse d'espèces protégées, éviter les feux, sensibiliser.
\end{itemize}
\end{aretenir}

\coursec{TP N°10 : Pratique du reboisement}

Après avoir identifié les menaces qui pèsent sur nos écosystèmes et les stratégies pour les restaurer, il est temps de passer à l'action concrète. La restauration de la biodiversité ne se décrète pas, elle se construit, plant par plant. Ce Travail Pratique est l'aboutissement naturel de notre séquence : vous allez mettre en œuvre, sur le terrain de l'établissement ou dans un espace communautaire voisin, les gestes techniques qui transforment une parcelle dégradée en un écosystème naissant. Au Cameroun, où la déforestation pour l'agriculture sur brûlis, le bois de chauffe ou l'urbanisation grignote chaque année des milliers d'hectares, savoir reboiser efficacement est une compétence citoyenne majeure.

\courssub{Objectif et choix des essences : planter l'arbre adapté au milieu}

\begin{definition}[Objectif du TP]
Réaliser concrètement l'ensemble des opérations techniques conduisant à l'installation durable de jeunes arbres (plants) sur une parcelle, depuis la production ou l'acquisition des plants en pépinière jusqu'au suivi de leur reprise, en respectant les exigences écologiques des essences choisies.
\end{definition}

Le succès d'un reboisement ne tient pas au hasard. Il repose d'abord sur un choix judicieux des espèces. Nous distinguons trois grandes catégories d'essences locales adaptées à nos contextes (forêt dense humide du Sud, savane de l'Adamaoua, zone soudano-sahélienne du Nord) :

\begin{itemize}
    \item \textbf{Essences fruitières :} \textit{Mangifera indica} (Manguier), \textit{Persea americana} (Avocatier), \textit{Citrus sinensis} (Oranger), \textit{Psidium guajava} (Goyavier), \textit{Dacryodes edulis} (Safoutier ou Prunier d'Afrique). Elles assurent la sécurité alimentaire et génèrent des revenus.
    \item \textbf{Essences forestières nobles (à croissance lente mais durables) :} \textit{Milicia excelsa} (Iroko), \textit{Entandrophragma cylindricum} (Sapelli), \textit{Terminalia superba} (Fraké), \textit{Triplochiton scleroxylon} (Ayous). Elles reconstituent la strate dominante de la forêt et stockent le carbone sur le long terme.
    \item \textbf{Essences à croissance rapide (pionnières, fertilisantes) :} \textit{Acacia mangium}, \textit{Leucaena leucocephala}, \textit{Gliricidia sepium}, \textit{Calliandra calothyrsus}. Ces légumineuses enrichissent le sol en azote grâce à des bactéries (\textit{Rhizobium}) vivant en association avec leurs racines, améliorent rapidement la structure du sol par leur litière et offrent du bois de feu à court terme (3 à 5 ans).
\end{itemize}

\begin{attention}[Choix de l'essence et station]
Ne jamais planter un jeune plant d'essence forestière de grande taille (ex : Iroko) en plein champ découvert sans protection : le jeune plant souffre du soleil brûlant et du dessèchement. À l'inverse, une essence pionnière qui aime la lumière (ex : \textit{Acacia}) ne s'épanouira pas sous un couvert dense. Adaptez toujours l'essence à la \cle{station}, c'est-à-dire au type de sol, à la pluviométrie, à l'exposition et à la végétation déjà présente.
\end{attention}

\begin{exemplebox}[Choix pour la cour de l'école (Yaoundé, zone de forêt)]
Pour une cour d'école à Yaoundé (pluviométrie environ \SI{1600}{mm/an}, sol ferrallitique), on privilégiera un mélange :
\begin{itemize}
    \item \textbf{Bordures et ombrage :} Manguier, Safoutier (fruitiers, ombre dense, valorisation par les élèves).
    \item \textbf{Brise-vent et haie vive :} \textit{Gliricidia sepium} (croissance très rapide, bouturable, enrichit le sol en azote).
    \item \textbf{Boisement d'agrément :} \textit{Terminalia mantaly} (ornemental, croissance modérée) ou jeunes plants d'Iroko si une strate arbustive protectrice existe déjà.
\end{itemize}
\end{exemplebox}

\courssub{Les quatre étapes du reboisement : protocole rigoureux}

\begin{methode}[Protocole du reboisement en 4 étapes]
\textbf{Étape 1 -- Pépinière : produire ou acquérir des plants de qualité.}
\begin{enumerate}
    \item \textbf{Semis :} en godets (sacs plastiques noirs perforés de 15 $\times$ 25 cm) remplis d'un substrat composé de 1/3 de terre de surface, 1/3 de sable et 1/3 de compost ou fumier bien décomposé. Semer 2 à 3 graines par godet, puis ne garder qu'un seul plant vigoureux après la levée.
    \item \textbf{Bouturage (pour \textit{Gliricidia}, \textit{Leucaena}) :} prélever des boutures de 20 à 30 cm, tremper la base dans l'eau, planter en pépinière ombragée (ombrière à 50\%).
    \item \textbf{Conduite de la pépinière (3 à 6 mois) :} arrosage quotidien (matin et soir), désherbage, protection contre les ravageurs (chenilles, termites), puis \textbf{endurcissement} : on enlève progressivement l'ombrière pendant les 2 semaines qui précèdent la plantation afin d'habituer les plants au plein soleil.
    \item \textbf{Critères d'un bon plant :} hauteur de 30 à 50 cm selon l'essence, collet (zone de transition entre racine et tige) de diamètre supérieur ou égal à 5--8 mm, système racinaire dense et non enroulé (pas de « chignon » au fond du godet), feuillage sain, exempt de maladies.
\end{enumerate}

\textbf{Étape 2 -- Préparation du terrain : accueillir la racine.}
\begin{enumerate}
    \item \textbf{Nettoyage :} débroussaillage manuel (coupe à ras, pas de labour profond qui détruit la structure du sol et la matière organique). Conserver les débris végétaux au sol.
    \item \textbf{Alignement et piquetage :} tracer les lignes (cordeau, piquets). Respecter les \textbf{distances de plantation} adaptées à l'essence et au but recherché (voir tableau ci-dessous).
    \item \textbf{Trous de plantation :} dimensions standard \textbf{40 cm $\times$ 40 cm $\times$ 40 cm} (longueur $\times$ largeur $\times$ profondeur). Séparer la terre de surface (premiers 15 à 20 cm, riche en humus) de la terre profonde (moins fertile).
    \item \textbf{Amendement :} mélanger la terre de surface avec \textbf{5 à 10 kg de compost ou de fumier bien décomposé} par trou. Remettre ce mélange au fond du trou (environ 15 cm d'épaisseur) pour créer un lit fertile.
\end{enumerate}

\textbf{Étape 3 -- Mise en terre : le geste technique décisif.}
\begin{enumerate}
    \item \textbf{Période impérative :} \textbf{début de la saison des pluies} (mars-avril au Sud, mai-juin au Nord). Le sol doit être humide en profondeur.
    \item \textbf{Extraction du plant :} couper le fond du godet, fendre le côté verticalement. Retirer le godet \textbf{sans casser la motte}. Ne jamais tirer sur la tige.
    \item \textbf{Plantation :} placer la motte au centre du trou. Le \textbf{collet doit affleurer le niveau du sol} : s'il est enterré, il pourrit ; s'il est surélevé, les racines supérieures se dessèchent.
    \item \textbf{Rebouchage :} reboucher avec le mélange terre-compost en tassant \textbf{progressivement par couches} pour chasser l'air sans trop compacter le sol.
    \item \textbf{Cuvette d'arrosage :} former une cuvette circulaire (diamètre environ 60 cm, bordure de 10 à 15 cm de haut) autour du pied pour retenir l'eau.
    \item \textbf{Arrosage de plantation :} \textbf{20 à 30 litres d'eau par plant} (même s'il pleut) pour assurer un bon contact entre les racines et le sol.
    \item \textbf{Tuteurage (si vent fort ou plant grand) :} planter un tuteur (bambou, bois) \textit{avant} la mise en terre pour ne pas blesser les racines. Attacher en « 8 » avec un lien souple, non blessant.
    \item \textbf{Paillage :} couvrir la cuvette de paille, de feuilles sèches ou d'herbe coupée (épaisseur 5 à 10 cm) pour limiter l'évaporation et l'érosion.
\end{enumerate}

\textbf{Étape 4 -- Protection et suivi : la durée fait la forêt.}
\begin{enumerate}
    \item \textbf{Arrosage de soutien :} 10 à 15 L par plant et par semaine en l'absence de pluie pendant les 6 à 12 premiers mois (phase critique d'enracinement).
    \item \textbf{Protection physique :} grillage ou haie vive contre la divagation des animaux (chèvres, moutons, bovins). Pare-feu (bande débroussaillée de 3 à 5 m) autour de la parcelle avant la saison sèche.
    \item \textbf{Entretien :} sarclage manuel (rayon 50 cm) tous les 2 à 3 mois pendant les 2 premières années pour éliminer la concurrence des herbes. Remplacement des plants morts au début de la 2\up{e} saison des pluies.
    \item \textbf{Suivi quantitatif :} comptage des survivants à 3 mois, 6 mois et 1 an. Calcul du \textbf{taux de reprise} :
    \[ \text{Taux de reprise} = \frac{\text{Nombre de plants vivants}}{\text{Nombre de plants mis en terre}} \times 100 \]
    Objectif : taux supérieur à 80\%.
\end{enumerate}
\end{methode}

\begin{exemplebox}[Distances de plantation recommandées]
L'espacement conditionne l'architecture future du peuplement et la facilité d'entretien.
\begin{center}
\begin{tabularx}{\linewidth}{|X|l|l|}\hline
\rowcolor{popblueL}\textbf{Type d'essence / Objectif} & \textbf{Espacement (m $\times$ m)} & \textbf{Densité (plants/ha)} \\ \hline
Fruitiers (Manguier, Avocatier, Safoutier) & $10 \times 10$ à $12 \times 12$ & 70 à 100 \\ \hline
Essences forestières nobles (Iroko, Sapelli, Ayous) & $4 \times 4$ à $5 \times 5$ & 400 à 625 \\ \hline
Essences à croissance rapide, bois de feu, haies vives & $2 \times 2$ à $3 \times 3$ & 1100 à 2500 \\ \hline
Agroforesterie (arbres alignés dans les cultures) & 10 m entre lignes, 2 m sur la ligne & 500 \\ \hline
\end{tabularx}
\end{center}
\textbf{Règle pratique pour l'élève :} « 3 à 5 mètres entre les plants » est un bon compromis pour un boisement scolaire mixte (fruitiers + forestiers) permettant le passage et l'entretien.
\end{exemplebox}

\courssub{Illustration technique : les étapes et le plant bien mis en terre}

\begin{popfigure}
\begin{tikzpicture}[
    node distance=1.5cm and 2.5cm,
    box/.style={rectangle, draw=popdark, thick, fill=popcream, rounded corners, minimum width=3.2cm, minimum height=1.8cm, align=center, font=\small},
    arrow/.style={-Stealth, thick, popblue},
    labelbox/.style={rectangle, draw=poporange, fill=poporangeL, rounded corners, font=\tiny, align=center, minimum width=2.5cm}
]
% Nodes
\node[box, align=center] (pep){\textbf{ÉTAPE 1 : PÉPINIÈRE}\\[0.2cm] Semis / Bouturage\\Godets + Substrat\\Conduite 3-6 mois\\Endurcissement};
\node[box, right=of pep, align=center] (prep){\textbf{ÉTAPE 2 : PRÉPARATION}\\[0.2cm] Débroussaillage\\Alignement / Piquetage\\Trous 40$\times$40$\times$40 cm\\Amendement (Compost)};
\node[box, right=of prep, align=center] (plant){\textbf{ÉTAPE 3 : PLANTATION}\\[0.2cm] Début saison pluies\\Mise en terre (Collet=Sol)\\Cuvette + Arrosage 20-30 L\\Tuteur + Paillage};
\node[box, right=of plant, align=center] (suivi){\textbf{ÉTAPE 4 : SUIVI}\\[0.2cm] Arrosage soutien\\Protection (Animaux/Feu)\\Sarclage / Remplacement\\Comptage (Taux reprise)};

% Arrows
\draw[arrow] (pep.east) -- node[above, font=\small\bfseries] {Plants prêts} (prep.west);
\draw[arrow] (prep.east) -- node[above, font=\small\bfseries] {Terrain prêt} (plant.west);
\draw[arrow] (plant.east) -- node[above, font=\small\bfseries] {Plant installé} (suivi.west);

% Key details labels
\node[labelbox, above=0.3cm of pep.north, anchor=south] {Substrat : 1/3 terre + 1/3 sable + 1/3 compost};
\node[labelbox, above=0.3cm of prep.north, anchor=south] {Séparer terre de surface / terre profonde};
\node[labelbox, above=0.3cm of plant.north, anchor=south] {Collet au niveau du sol !};
\node[labelbox, above=0.3cm of suivi.north, anchor=south] {Objectif : taux de reprise $>$ 80\%};
\end{tikzpicture}
\legende{Schéma synthétique des 4 étapes du reboisement : de la production du plant en pépinière au suivi de la reprise sur le terrain. Chaque étape conditionne la suivante.}
\end{popfigure}

\begin{popfigure}
\begin{tikzpicture}[
    scale=0.9,
    every node/.style={font=\small},
    soil/.style={fill=brown!40!black, pattern=north east lines, pattern color=brown!60!black},
    compost/.style={fill=green!30!brown, pattern=crosshatch dots, pattern color=green!50!brown},
    mulch/.style={fill=yellow!60!brown, pattern=horizontal lines, pattern color=yellow!50!brown},
]
% Sol en place
\fill[soil] (-3,-3) rectangle (3,-1);
% Trou rebouché (mélange)
\fill[compost] (-1.5,-3) rectangle (1.5,-1);
% Cuvette
\draw[thick, brown!70!black] (-2,-1) -- (-2,0.5) -- (2,0.5) -- (2,-1);
\fill[mulch] (-2,-1) rectangle (2,-0.5); % Paillis surface
% Plant
\draw[line width=3pt, brown!60!black] (0,-1) -- (0,2.5); % Tige
% Branches
\draw[line width=2pt, brown!60!black] (0,1.5) -- (-0.8,2.2);
\draw[line width=2pt, brown!60!black] (0,1.5) -- (0.8,2.2);
\draw[line width=1.5pt, brown!60!black] (0,2.0) -- (-0.5,2.6);
\draw[line width=1.5pt, brown!60!black] (0,2.0) -- (0.5,2.6);
% Feuilles
\foreach \x/\y in {-0.8/2.2, 0.8/2.2, -0.5/2.6, 0.5/2.6, 0/2.5} {
    \fill[green!60!black, rotate=30] (\x,\y) ellipse (0.3 and 0.1);
}
% Racines (schématisées)
\draw[line width=1.5pt, brown!50!black] (0,-1) .. controls (-1,-1.5) and (-1.2,-2) .. (-1.2,-2.5);
\draw[line width=1.5pt, brown!50!black] (0,-1) .. controls (1,-1.5) and (1.2,-2) .. (1.2,-2.5);
\draw[line width=1.5pt, brown!50!black] (0,-1) .. controls (0,-1.8) and (0,-2.3) .. (0,-2.7);
% Collet
\draw[thick, red] (-0.3,-1) -- (0.3,-1);
\node[red, anchor=north] at (0,-1.1) {\textbf{Collet}};
% Tuteur
\draw[line width=4pt, orange!70!black] (1.2,-2.5) -- (1.2,2.8);
\draw[thick, poppink] (0,1.8) .. controls (0.6,1.8) .. (1.2,1.8); % Attache en 8
\node[orange!70!black, anchor=west] at (1.3,2.5) {\textbf{Tuteur}};
% Cuvette label
\node[anchor=east, align=right] at (-2.3,0) {\textbf{Cuvette}\\\textbf{d'arrosage}};
% Paillis label
\node[anchor=west, align=left] at (2.3,-0.7) {\textbf{Paillis}\\(paille, feuilles)};
% Flèche arrosage
\draw[popblue, very thick, -Stealth] (0,3) -- (0,2.7);
\node[popblue, anchor=south] at (0,3) {\textbf{Arrosage 20-30 L}};
% Légendes coté
\begin{scope}[xshift=4.5cm]
\node[anchor=west] at (0,2.5) {\textbf{1.} Tige droite, extrémité intacte};
\node[anchor=west] at (0,2.0) {\textbf{2.} Collet au niveau du sol};
\node[anchor=west] at (0,1.5) {\textbf{3.} Cuvette + paillis};
\node[anchor=west] at (0,1.0) {\textbf{4.} Tuteur + attache en 8};
\node[anchor=west] at (0,0.5) {\textbf{5.} Racines étalées, non enroulées};
\node[anchor=west] at (0,0.0) {\textbf{6.} Mélange terre + compost};
\node[anchor=west] at (0,-0.5) {\textbf{7.} Sol naturel non compacté};
\end{scope}
\end{tikzpicture}
\legende{Plant correctement mis en terre : respect du collet (2), cuvette d'arrosage et paillis (3), tuteurage souple (4), racines libérées dans un substrat amendé (5 et 6). C'est la garantie d'une reprise racinaire rapide.}
\end{popfigure}

\courssub{Compte rendu de TP : structurer sa démarche scientifique}

Le compte rendu n'est pas une simple liste de courses. C'est la trace écrite de votre démarche d'investigation. Il doit permettre à un tiers de reproduire le travail et d'en juger la validité.

\begin{methode}[Rédiger un compte rendu de TP de reboisement]
\begin{enumerate}
    \item \textbf{Titre précis :} « TP N°10 : Reboisement de la parcelle X de l'établissement [Nom] — Essences : [Liste] — Date : [JJ/MM/AAAA] ».
    \item \textbf{Objectif :} formuler la question scientifique ou technique. Ex : « Quel taux de reprise obtient-on pour \textit{Mangifera indica} et \textit{Gliricidia sepium} plantés en début de saison des pluies sur sol ferrallitique amendé ? »
    \item \textbf{Matériel et espèces :} liste exhaustive (godets, pioches, cordeau, mètre, arrosoirs, compost, plants identifiés par étiquettes).
    \item \textbf{Protocole suivi (démarche) :} décrire les 4 étapes \textbf{avec vos mesures réelles} (ex : « Trous de 40 $\times$ 40 $\times$ 40 cm réalisés à la pioche ; 5 kg de compost par trou ; espacement 4 m $\times$ 4 m ; arrosage de 25 L par plant »). Mentionner les aléas (pluie le jour J, sol dur, godets déchirés).
    \item \textbf{Observations et résultats (données brutes) :} tableau de suivi.
    \begin{center}
    \begin{tabularx}{\linewidth}{|l|X|X|X|X|}\hline
    \rowcolor{popblueL}\textbf{Essence} & \textbf{Plants mis en terre ($N_0$)} & \textbf{Vivants à 1 mois ($N_1$)} & \textbf{Vivants à 3 mois ($N_3$)} & \textbf{Taux de reprise à 3 mois (\%)} \\ \hline
    Manguier & 20 & 19 & 18 & 90 \\ \hline
    \textit{Gliricidia} & 20 & 20 & 19 & 95 \\ \hline
    Iroko & 10 & 8 & 6 & 60 \\ \hline
    \end{tabularx}
    \end{center}
    Exemple de calcul pour le Manguier : $TR = \dfrac{18}{20} \times 100 = 90\%$.
    \item \textbf{Interprétation et discussion :} expliquer les résultats. Pourquoi l'Iroko a-t-il un taux plus faible ? (Ex : jeune plant sensible au plein soleil, installé sans ombrage transitoire, ou godets trop petits provoquant un enroulement des racines.)
    \item \textbf{Conclusion justifiée :} répondre à l'objectif. « Le reboisement est réussi pour les essences pionnières et fruitières (taux supérieur à 90\%). L'Iroko nécessite une couverture préalable (par exemple sous \textit{Gliricidia}) pour atteindre 80\%. L'apport de compost et l'arrosage de plantation ont été déterminants. »
    \item \textbf{Perspectives et recommandations :} actions pour l'an prochain (remplacement des plants morts, taille de formation, installation d'ombrières pour les essences forestières).
\end{enumerate}
\end{methode}

\begin{savaistu}[L'arbre à l'école : un capital vivant]
Au Cameroun, les programmes d'éducation environnementale encouragent chaque établissement à planter des arbres chaque année. Un manguier adulte peut produire de 200 à 500 kg de fruits par an. Un \textit{Gliricidia} taillé régulièrement fournit du paillis riche en azote pour le potager scolaire et du fourrage pour l'élevage de l'école. Le reboisement scolaire n'est pas un exercice académique : c'est un \textbf{investissement productif} qui nourrit les élèves et rafraîchit les salles de classe (l'ombrage peut abaisser la température de 2 à 4 \degres C). Votre TP d'aujourd'hui est le capital de l'école de demain.
\end{savaistu}

\begin{attention}[Pièges classiques à éviter dans le compte rendu]
\begin{itemize}
    \item \textbf{Confondre « arrosage » et « pluie » :} noter « il a plu » ne remplace pas la mesure du volume d'eau apporté au pied (cuvette). Quantifiez : « Arrosage de 20 L effectué le 15/04 car 10 jours sans pluie ».
    \item \textbf{Oublier le collet :} ne pas mentionner le positionnement du collet (au niveau du sol, enterré ou surélevé) invalide l'analyse des mortalités.
    \item \textbf{Données sans unité :} « trous de 40 » (40 quoi ? cm ? m ?). Toujours écrire : 40 cm, 20 L, 5 kg.
    \item \textbf{Conclusion sans chiffre :} « ça a bien marché » est insuffisant. Écrivez plutôt : « Taux de reprise global de 85\% (Manguier 90\%, Iroko 60\%) ».
    \item \textbf{Photo sans légende :} une photo de vous tenant une pioche n'est pas une illustration scientifique. Photographiez le trou amendé, le plant dans la cuvette paillée, l'alignement, et légendez : « Fig. 1 : plant de \textit{Gliricidia} 1 mois après plantation, cuvette intacte, reprise visible (nouvelles feuilles) ».
\end{itemize}
\end{attention}

\begin{exoresolu}[Analyse d'un échec de reboisement]
\textbf{Énoncé :} un club SVT a planté 50 plants d'Iroko (\textit{Milicia excelsa}) en janvier (saison sèche) sur une parcelle nue, sans arrosage, sans compost, espacés de 2 m. Trois mois plus tard, 48 plants sont secs. Analysez les causes techniques de cet échec et proposez un protocole correctif pour la campagne suivante.
\tcblower
\textbf{Solution détaillée :}
\begin{enumerate}
    \item \textbf{Analyse des causes (facteurs limitants) :}
    \begin{itemize}
        \item \textbf{Période (cause majeure) :} janvier correspond à la saison sèche. L'absence d'eau dans le sol provoque le dessèchement rapide du système racinaire, pas encore ancré. \textit{Règle violée : planter en début de saison des pluies.}
        \item \textbf{Absence d'arrosage de soutien :} aucun apport d'eau (20 à 30 L par plant à la plantation, puis 10 à 15 L par semaine) pour compenser l'évaporation et la transpiration des feuilles.
        \item \textbf{Station inadaptée :} le jeune plant d'Iroko supporte mal le plein soleil brûlant et la chaleur du sol nu ; il a besoin d'un ombrage latéral pendant ses premières années.
        \item \textbf{Absence d'amendement :} sol nu, pauvre et compacté, d'où une mauvaise reprise racinaire.
        \item \textbf{Espacement trop serré (2 m) :} concurrence racinaire immédiate pour l'eau, qui est rare.
    \end{itemize}
    \item \textbf{Protocole correctif (campagne suivante, début de la saison des pluies) :}
    \begin{enumerate}
        \item \textbf{Année 1 :} planter une essence de couverture pionnière, \textit{Gliricidia sepium} ou \textit{Acacia mangium}, à 3 m $\times$ 3 m, en début de saison des pluies, avec amendement (5 kg de compost par trou), arrosage de plantation et paillage.
        \item \textbf{Années 1 et 2 :} sarcler, protéger contre le feu et les animaux. Le \textit{Gliricidia} atteint 3 à 4 m, crée un ombrage latéral et enrichit le sol (azote, litière).
        \item \textbf{Année 2 ou 3 :} planter les jeunes Iroko (issus de pépinière de 6 à 8 mois, collet de 8 à 10 mm) \textit{sous} le couvert du \textit{Gliricidia}, à un espacement de 5 m $\times$ 5 m.
        \item \textbf{Suivi :} éclaircir progressivement le \textit{Gliricidia} (taille) au fur et à mesure de la croissance de l'Iroko.
    \end{enumerate}
    \item \textbf{Conclusion :} l'échec n'est pas dû à une « mauvaise qualité » de l'Iroko, mais au non-respect de ses besoins écologiques (ombrage juvénile) et du calendrier des pluies. La réussite passe par une \textbf{plantation en mélange successif} : essence pionnière protectrice d'abord, essence noble ensuite.
\end{enumerate}
\end{exoresolu}

\begin{aretenir}[Savoir-faire évalué : le compte rendu de TP]
Au BEPC ou en contrôle continu, vous devez être capable de :
\begin{itemize}
    \item \textbf{Structurer} un compte rendu : titre, objectif, matériel, protocole (étapes numérotées), résultats (tableau chiffré), interprétation, conclusion.
    \item \textbf{Justifier} chaque geste technique par son rôle (ex : « paillage \textit{pour} limiter l'évaporation et l'érosion » ; « collet au niveau du sol \textit{pour} éviter la pourriture du collet ou le dessèchement des racines superficielles »).
    \item \textbf{Calculer et interpréter} un taux de reprise : $TR = \dfrac{N_{\text{vivants}}}{N_{\text{plantés}}} \times 100$, et discuter les causes d'un taux inférieur à 80\%.
    \item \textbf{Proposer des mesures correctives} argumentées (changement d'essence, de période, de préparation du sol, de protection).
\end{itemize}
\textbf{Mots-clés à maîtriser :} pépinière, godet, substrat, endurcissement, collet, cuvette d'arrosage, paillis, tuteurage, taux de reprise, essence pionnière, essence forestière noble, station, amendement.
\end{aretenir}

\begin{exercice}[Mise en situation : votre premier chantier]
Vous êtes responsable du reboisement d'une parcelle de 0,5 ha (5000 m$^2$) dans l'enceinte du collège, en zone soudano-sahélienne (Maroua, pluviométrie environ 800 mm/an, saison des pluies de juin à septembre). Le sol est sableux et pauvre en matière organique. Vous disposez de 200 plants d'\textit{Acacia nilotica} (épineux, enrichit le sol en azote, fourrage, bois dur), 50 plants de \textit{Moringa oleifera} (feuilles et fruits comestibles, croissance très rapide) et 50 plants d'\textit{Adansonia digitata} (Baobab, croissance lente, longévité exceptionnelle).
\begin{enumerate}
    \item Proposez un plan de plantation (schéma simplifié avec distances) optimisant l'espace et les interactions (ex : haie vive périphérique, alignements intérieurs).
    \item Calculez le nombre de trous à creuser et la quantité totale de compost nécessaire (5 kg par trou).
    \item Listez le matériel indispensable pour l'équipe de 10 élèves.
    \item Rédigez la consigne de sécurité n°1 à donner avant de commencer (manipulation des outils, épines de l'\textit{Acacia}).
    \item Prévoyez le calendrier de suivi pour la première année (plantation, 1\up{er} sarclage, 1\up{er} comptage, protection contre le feu).
\end{enumerate}
\end{exercice}

\begin{corrige}[Exercice : Mise en situation]
\begin{enumerate}
    \item \textbf{Plan :} sur le périmètre (environ 200 m linéaires), planter une haie vive d'\textit{Acacia nilotica} à 1 m d'espacement (200 plants) : protection contre les animaux et le feu, plus fourrage. À l'intérieur, aligner les 50 \textit{Moringa} tous les 5 m sur 2 lignes (ombrage rapide, feuilles comestibles pour la cantine). Répartir les 50 Baobabs en quinconce à 10 m $\times$ 10 m (futurs grands arbres, ombre durable).
    \item \textbf{Calculs :} nombre total de plants = $200 + 50 + 50 = 300$. Nombre de trous = 300. Quantité de compost = $300 \times 5 = 1500$ kg, soit \textbf{1,5 tonne}. Prévoir 2 tonnes pour compenser les pertes et le tassement.
    \item \textbf{Matériel :} 10 pioches, 5 pelles, 10 arrosoirs de 10 L, 2 cordeaux de 50 m, 2 décamètres, 10 paires de gants en cuir (épines), 300 étiquettes et feutres indélébiles, des sacs de compost, 1 brouette, 1 seau et du savon (hygiène).
    \item \textbf{Sécurité n°1 :} « Port des gants en cuir OBLIGATOIRE pour manipuler l'\textit{Acacia nilotica} (épines de 3 à 5 cm, risque de blessure et d'infection). Manipulation de la pioche et de la pelle : distance de sécurité de 2 m entre les travailleurs ; ne jamais lever l'outil au-dessus de l'épaule. »
    \item \textbf{Calendrier (Maroua) :}
    \begin{itemize}
        \item \textbf{1\up{er}--10 juin :} plantation (début des pluies bien établies).
        \item \textbf{15 juin -- 15 juillet :} arrosage de soutien 2 fois par semaine si l'arrêt des pluies dépasse 5 jours ; 1\up{er} sarclage (rayon 30 cm).
        \item \textbf{15 juillet :} 1\up{er} comptage de la reprise (objectif : taux supérieur à 85\%).
        \item \textbf{Août :} 2\up{e} sarclage ; réparation des cuvettes et du paillis.
        \item \textbf{Novembre (fin des pluies) :} 2\up{e} comptage ; installation du pare-feu (bande débroussaillée de 5 m) en périphérie.
        \item \textbf{Mars (saison sèche) :} arrosage d'urgence des Baobabs et \textit{Moringa} si stress visible (flétrissement) ; taille de formation de la haie d'\textit{Acacia} (hauteur 1,5 m).
    \end{itemize}
\end{enumerate}
\end{corrige}

Ce TP termine notre unité sur les écosystèmes. Vous avez maintenant les clés pour comprendre la forêt et la savane, analyser leurs dysfonctionnements et \textbf{agir} pour les restaurer. Le reboisement n'est pas une fin en soi : c'est le début d'une relation de long terme avec un écosystème que vous avez contribué à faire naître. Observez-le, mesurez-le, protégez-le : c'est tout l'esprit des SVTEEHB.

\newpage

\coursec{Activité d'intégration}

\courssub{Objectifs de l'activité}

À la fin de cette activité, tu seras capable de :
\begin{itemize}
\item mobiliser les notions d'\cle{écosystème}, de \cle{biodiversité} et d'\cle{interdépendance} pour analyser une situation réelle de dégradation du milieu ;
\item identifier les \cle{activités humaines destructrices} d'un écosystème (feux de brousse, déforestation, braconnage) et expliquer leurs conséquences ;
\item proposer et justifier un plan concret de \cle{restauration} et de \cle{conservation} de la biodiversité (reboisement) ;
\item rédiger une conclusion argumentée, claire et convaincante, destinée à un décideur (le chef du village).
\end{itemize}

\courssub{Situation}

\textbf{Sauvons la colline de notre village : un projet de reboisement.}

Le village de \textbf{Ngoumou}, situé dans la région du Centre, est dominé par une colline appelée « colline de Mbog ». Les anciens racontent qu'il y a 25 ans, cette colline était couverte d'une forêt dense où vivaient de nombreux animaux et où poussaient des arbres précieux. Aujourd'hui, la colline est presque nue : les \cle{feux de brousse} répétés (allumés chaque saison sèche pour défricher ou chasser) et la \cle{coupe de bois} de chauffe ont fait disparaître presque tous les arbres.

Les conséquences sont graves pour les habitants :
\begin{itemize}
\item la \textbf{source} située en contrebas de la colline, qui alimente le village en eau potable, \textbf{tarit chaque saison sèche} depuis quelques années ;
\item les pluies emportent la terre des champs de manioc et d'arachide situés sur les pentes (\cle{érosion}) : les rendements agricoles diminuent d'année en année.
\end{itemize}

Le comité de développement du village a réuni trois documents pour comprendre la situation.

\textbf{Document 1 : quelques espèces de la colline de Mbog, en 2000 et en 2024.}

\begin{center}
\begin{tabularx}{\linewidth}{|l|X|X|}
\hline
\rowcolor{popblueL} \textbf{Espèce} & \textbf{Effectif estimé en 2000} & \textbf{Effectif estimé en 2024} \\
\hline
Iroko (arbre) & 300 individus & 12 individus \\
\hline
Fraké (arbre) & 450 individus & 25 individus \\
\hline
Colobe (singe) & 80 individus & 0 (disparu) \\
\hline
Céphalophe (petite antilope) & 120 individus & 0 (disparu) \\
\hline
Calao (oiseau) & 200 individus & 15 individus \\
\hline
\end{tabularx}
\end{center}

\textbf{Document 2 : débit de la source de Ngoumou mesuré en fin de saison sèche (mars), de 2004 à 2024.}

\begin{popfigure}
\begin{center}
\begin{tikzpicture}
\begin{axis}[
width=0.85\linewidth, height=6cm,
xlabel={Année}, ylabel={Débit de la source (L/min)},
xmin=2004, xmax=2024, ymin=0, ymax=60,
xtick={2004,2008,2012,2016,2020,2024},
grid=both, grid style={popline!40},
]
\addplot[popcurve, mark=*, mark size=1.5pt] coordinates {
(2004,50) (2008,44) (2012,35) (2016,24) (2020,12) (2024,5)
};
\end{axis}
\end{tikzpicture}
\end{center}
\legende{Courbe du débit de la source en fin de saison sèche : il est passé de 50~L/min en 2004 à 5~L/min en 2024.}
\end{popfigure}

\textbf{Document 3 : résultats d'un essai de plantation réalisé par un village voisin (100 plants d'iroko mis en terre en 2019).}

\begin{center}
\begin{tabularx}{\linewidth}{|l|X|X|}
\hline
\rowcolor{popblueL} \textbf{Parcelle} & \textbf{Plants vivants après 1 an} & \textbf{Plants vivants après 5 ans} \\
\hline
Parcelle A : plantation \textbf{sans suivi} (pas d'arrosage, pas de protection contre les feux) & 40 sur 100 & 8 sur 100 \\
\hline
Parcelle B : plantation \textbf{avec suivi} (arrosage en saison sèche, pare-feu, protection contre les animaux) & 85 sur 100 & 70 sur 100 \\
\hline
\end{tabularx}
\end{center}

Le chef du village demande aux élèves de la classe de 3ème du collège de Ngoumou de l'aider à comprendre ce qui s'est passé et à préparer un \textbf{plan d'action} pour sauver la colline.

\courssub{Consignes}

\begin{enumerate}
\item À l'aide du document 1 et de la situation décrite, \textbf{identifie les causes} de la dégradation de la colline de Mbog et montre qu'il s'agit bien d'activités humaines destructrices de l'écosystème.
\item En t'appuyant sur la notion d'\textbf{interdépendance} entre les êtres vivants et leur milieu, \textbf{explique} pourquoi la disparition des arbres a entraîné :
\begin{enumerate}
\item le tarissement progressif de la source (document 2) ;
\item l'érosion des sols des champs sur les pentes.
\end{enumerate}
\item \textbf{Propose un plan de reboisement en 4 étapes}, en précisant pour chaque étape les actions à mener, le choix des essences locales à planter et le calendrier (saison favorable). Appuie tes propositions sur les résultats du document 3.
\item \textbf{Rédige une conclusion argumentée} (10 à 15 lignes) destinée au chef du village, pour le convaincre d'adopter ton plan d'action.
\end{enumerate}

\courssub{Ressources à mobiliser}

\begin{aretenir}[Ce que tu dois te rappeler]
\begin{itemize}
\item Un \cle{écosystème} est formé d'un milieu de vie (\cle{biotope}) et de l'ensemble des êtres vivants qui y habitent (\cle{biocénose}).
\item La \cle{biodiversité} est la variété des espèces vivantes dans un écosystème.
\item Dans un écosystème, les êtres vivants sont \cle{interdépendants} : ils dépendent les uns des autres (nourriture, abri, pollinisation\ldots) et du milieu. La disparition d'un élément entraîne des conséquences sur tout l'écosystème.
\item Les activités humaines qui détruisent les écosystèmes : \cle{feux de brousse}, \cle{déforestation} (coupe des arbres), \cle{braconnage}, pollution\ldots
\item Pour restaurer un écosystème dégradé : le \cle{reboisement} (plantation d'arbres), la protection contre les feux, la sensibilisation des populations. Pour conserver : créer des aires protégées, lutter contre le braconnage.
\item Les arbres retiennent l'eau de pluie, facilitent son infiltration dans le sol (ce qui alimente les sources) et fixent le sol par leurs racines (ce qui empêche l'érosion).
\end{itemize}
\end{aretenir}

\courssub{Démarche et corrigé commenté}

\begin{exoresolu}[Sauvons la colline de Mbog]
Énoncé : à partir des trois documents, identifier les causes de la dégradation de la colline, expliquer le tarissement de la source et l'érosion, proposer un plan de reboisement en 4 étapes et rédiger une conclusion pour le chef du village.
\tcblower
\textbf{Question 1 : identification des causes de la dégradation.}

Le document 1 montre que, entre 2000 et 2024, les arbres (iroko : de 300 à 12 individus ; fraké : de 450 à 25 individus) ont presque disparu, et avec eux les animaux (colobe et céphalophe disparus, calaos réduits à 15 individus). La biodiversité de la colline s'est donc effondrée.

D'après la situation, cette disparition est due à deux \textbf{activités humaines destructrices} :
\begin{itemize}
\item les \cle{feux de brousse} répétés, allumés chaque saison sèche, qui brûlent la végétation, les jeunes arbres et détruisent ou chassent les animaux ;
\item la \cle{coupe de bois} de chauffe (déforestation), qui élimine directement les arbres.
\end{itemize}
On peut ajouter que les animaux restants ont été chassés ou ont fui un milieu devenu incapable de les nourrir et de les abriter : c'est la conséquence de la destruction de leur habitat.

\textbf{Question 2 : explication par l'interdépendance.}

\emph{a) Tarissement de la source.} Les arbres jouent un rôle essentiel dans le cycle de l'eau : leur feuillage freine les gouttes de pluie, et le sol forestier, couvert de litière et traversé par les racines, laisse l'eau \textbf{s'infiltrer lentement} dans le sous-sol. Cette eau souterraine alimente ensuite la source de façon régulière, même en saison sèche. Quand les arbres disparaissent, l'eau de pluie ruisselle rapidement en surface au lieu de s'infiltrer : les réserves souterraines ne se rechargent plus et la source tarit. C'est exactement ce que confirme le document 2 : le débit de la source est passé de 50~L/min en 2004 à seulement 5~L/min en 2024, soit une \textbf{division par 10 en 20 ans} (une perte de 90~\% du débit), au rythme de la disparition de la forêt.

\emph{b) Érosion des sols.} Les racines des arbres \textbf{fixent le sol} comme un filet, et le couvert végétal protège la terre de l'impact direct des gouttes de pluie. Sans arbres, l'eau de pluie frappe le sol nu et l'emporte le long des pentes : c'est l'\cle{érosion}. La terre fertile des champs de manioc et d'arachide est ainsi emportée, ce qui explique la baisse des rendements.

Ces deux phénomènes illustrent l'\cle{interdépendance} : en détruisant un seul élément de l'écosystème (les arbres), l'homme a provoqué des conséquences en chaîne sur l'eau, le sol, les animaux et finalement sur lui-même.

\textbf{Question 3 : plan de reboisement en 4 étapes.}

\begin{enumerate}
\item \textbf{Étape 1 — Sensibilisation et protection (dès maintenant).} Réunir les habitants pour expliquer les causes du problème ; interdire les feux de brousse sur la colline et créer des \textbf{pare-feu} (bandes de terre débroussaillée) ; réglementer la coupe de bois.
\item \textbf{Étape 2 — Préparation de la pépinière (saison sèche, janvier--février).} Installer une petite pépinière villageoise et produire des plants d'\textbf{essences locales} : iroko et fraké (arbres de la forêt d'origine, document 1), auxquels on peut ajouter des arbres fruitiers (manguier, avocatier) qui donnent un revenu aux habitants et les encouragent à protéger les plantations.
\item \textbf{Étape 3 — Plantation (début de la saison des pluies, mars--avril).} Planter les jeunes plants sur la colline au début des pluies pour que l'eau ne manque pas, en commençant par les zones proches de la source et les pentes érodées.
\item \textbf{Étape 4 — Suivi et entretien (toute l'année, pendant au moins 5 ans).} Arroser les plants en saison sèche, les protéger des animaux et des feux, remplacer les plants morts. Le document 3 prouve que cette étape est \textbf{indispensable} : sans suivi, seulement 8 plants sur 100 survivent après 5 ans, contre 70 sur 100 avec suivi. C'est le suivi qui fait la différence entre l'échec et la réussite du reboisement.
\end{enumerate}

\textbf{Question 4 : conclusion argumentée pour le chef du village (exemple de rédaction).}

« Chef du village, la colline de Mbog est malade parce que les feux de brousse et la coupe de bois ont détruit sa forêt. Sans arbres, l'eau de pluie ne pénètre plus dans le sol : notre source a perdu 90~\% de son débit en 20 ans et tarit chaque saison sèche. Sans racines pour retenir la terre, nos champs sont emportés par l'érosion et nos récoltes diminuent. Tout est lié dans la nature : en perdant les arbres, nous avons perdu l'eau, le sol et les animaux. Mais nous pouvons réparer ces dégâts. Nous te proposons un plan en 4 étapes : protéger la colline contre les feux, préparer une pépinière d'essences locales, planter au début des pluies, et surtout suivre les plants pendant plusieurs années, car l'expérience du village voisin montre que le suivi fait passer la réussite de 8 à 70 plants sur 100. Accepte ce plan, et dans quelques années la colline redeviendra verte, la source coulera à nouveau et nos champs seront protégés. »
\end{exoresolu}

\begin{attention}[Erreurs fréquentes à éviter]
\begin{itemize}
\item Ne confonds pas \textbf{cause} (les feux de brousse, la coupe de bois) et \textbf{conséquence} (tarissement de la source, érosion, disparition des animaux) : au BEPC, il faut les distinguer clairement.
\item Ne dis pas « les arbres fabriquent l'eau » : les arbres \textbf{facilitent l'infiltration} de l'eau de pluie dans le sol, ce qui recharge les réserves souterraines qui alimentent la source.
\item Un reboisement ne se résume pas à « planter des arbres » : sans \textbf{suivi} (arrosage, protection contre les feux), la plupart des plants meurent, comme le montre le document 3.
\end{itemize}
\end{attention}

\courssub{Bilan de l'activité}

Cette activité t'a permis de mettre en évidence que :
\begin{itemize}
\item un écosystème est un \textbf{tout} : la destruction d'un élément (la forêt) entraîne des conséquences en chaîne sur l'eau, le sol, les animaux et les activités humaines — c'est le sens exact de l'\cle{interdépendance} ;
\item les \cle{feux de brousse} et la \cle{déforestation} sont des activités humaines qui dégradent gravement les écosystèmes du Cameroun, avec des conséquences concrètes pour les populations (manque d'eau, baisse des récoltes) ;
\item la dégradation n'est pas une fatalité : le \cle{reboisement}, s'il est bien organisé (essences locales, bon calendrier, suivi régulier), permet de \textbf{restaurer} un écosystème et de \textbf{conserver} sa biodiversité ;
\item les données chiffrées (courbe du débit, tableau de survie des plants) sont des outils précieux pour \textbf{argumenter} et convaincre : un bon citoyen sait s'appuyer sur des faits pour proposer des solutions à sa communauté.
\end{itemize}

Tu es maintenant prêt à devenir, dans ton propre village ou quartier, un acteur de la protection de l'environnement !

\newpage

\coursec{Méthodes et exercices résolus}

\coursec{Leçon 14 : Étude des écosystèmes : La forêt et la savane}

\courssub{Notion d'écosystème et grands écosystèmes du Cameroun}

\begin{definition}[Écosystème, Biotope, Biocénose]
    Un \textbf{écosystème} est un ensemble formé par une communauté d'êtres vivants (\textbf{biocénose}) et son milieu de vie (\textbf{biotope}), fonctionnant comme une unité dynamique par les échanges de matière et d'énergie.
    \begin{itemize}
        \item \textbf{Biotope (Milieu physique) :} Facteurs \textbf{abiotiques} : climat (température, pluviométrie, lumière), sol (nature, pH, humidité), topographie, eau.
        \item \textbf{Biocénose (Peuplement) :} Ensemble des \textbf{populations} d'espèces différentes (végétaux, animaux, microorganismes) vivant en interaction.
        \item \textbf{Interactions :} Relations trophiques (proie-prédateur), compétition, symbiose, relations avec le milieu (nutrition, reproduction).
    \end{itemize}
    \textbf{Échelle :} De la mare (micro-écosystème) à la forêt dense (méso-écosystème) jusqu'à la biosphère (macro-écosystème).
\end{definition}

\begin{savaistu}[Les grands domaines écologiques du Cameroun : un gradient climatique Nord-Sud]
    Le Cameroun, « Afrique en miniature », présente une succession latitudinale d'écosystèmes déterminée par le gradient pluviométrique et la durée de la saison sèche :
    \begin{center}
    \begin{tabularx}{\linewidth}{|l|X|X|X|}\hline
    \rowcolor{popblueL} \textbf{Domaine} & \textbf{Localisation} & \textbf{Climat (P / Saison sèche)} & \textbf{Végétation climax / Sol dominant} \\ \hline
    \textbf{Forêt dense humide} & Sud, Est, Littoral, Sud-Ouest & $P > 1800$ mm / $< 2$ mois & Forêt sempervirente, strates multiples, lianes, épiphytes / \textbf{Sol ferrallitique} lessivé, humifère \\ \hline
    \textbf{Forêt claire / Mosaïque} & Transition Centre (Yaoundé) & $P \approx 1500\text{--}1800$ mm / $3\text{--}4$ mois & Forêt décidue, galerie forestière, savanes enclaves / Sol ferrallitique appauvri \\ \hline
    \textbf{Savane guinéenne (boisée)} & Adamaoua, Centre-Nord & $P \approx 1200\text{--}1500$ mm / $4\text{--}5$ mois & Arbres clairs (Isoberlinia), strate herbacée haute / \textbf{Sol ferrugineux tropical} lessivé \\ \hline
    \textbf{Savane soudanaenne} & Nord, Nord-Ouest & $P \approx 800\text{--}1200$ mm / $5\text{--}7$ mois & Arbres épineux (Acacia), herbes xérophiles (Andropogon) / Sol ferrugineux, vertisols (bas-fonds) \\ \hline
    \textbf{Steppe sahélienne} & Extrême-Nord (Maroua, Lac Tchad) & $P < 800$ mm / $> 7$ mois & Steppe à épineux, herbacées annuelles / Sol peu évolué, salé (vertisols) \\ \hline
    \textbf{Montagne (Mont Cameroun, Bamenda)} & Altitude $> 1500$ m & Tempéré, $P$ très élevé & Forêt de montagne, landes à \emph{Erica}, tourbières / Sol andosol, humique \\ \hline
    \end{tabularx}
    \end{center}
    \textbf{Facteur limitant principal :} La \textbf{disponibilité en eau} (pluviométrie $\times$ durée saison sèche) conditionne le type de végétation et la productivité primaire.
\end{savaistu}

\courssub{Méthodes et exercices résolus}

\begin{methode}[Caractériser et comparer deux grands écosystèmes du Cameroun (Forêt et Savane)]
    \textbf{Objectif :} À partir d'un texte, d'un tableau de données climatiques ou de photographies, identifier le type d'écosystème et justifier par les facteurs abiotiques et la structure de la végétation.
    \begin{enumerate}
        \item \textbf{Relever les données climatiques :} Température moyenne annuelle, pluviométrie annuelle, nombre de mois secs (pluviométrie $< \SI{100}{mm}$), saisonnalité.
        \item \textbf{Analyser la structure de la végétation :} Strates (forêt : 4-5 strates, canopée fermée ; savane : strate herbacée continue, strate arbustive/arborée discontinue), types de feuilles (mésophylles persistantes vs microphylles caduques/xérophiles), écorce (lisse/fine vs épaisse/rugueuse).
        \item \textbf{Identifier le type de sol (si donné) :} Sol ferrallitique lessivé (forêt) vs sol ferrugineux tropical / vertisols (savane).
        \item \textbf{Conclure :} Nommer l'écosystème (Forêt dense humide, Forêt claire, Savane boisée, Savane herbeuse, Steppe) et citer les \textbf{deux} facteurs climatiques déterminants (Température, Pluviométrie/Saison sèche).
    \end{enumerate}
    \textbf{Repères chiffrés (Cameroun) :}
    \begin{itemize}
        \item Forêt dense humide (Sud) : $P > \SI{1500}{mm}$, $T \approx \SI{25}{\celsius}$, saison sèche $< 2$ mois.
        \item Savane guinéenne (Centre/Adamaoua) : $P \approx \SI{1200}{--}\SI{1500}{mm}$, saison sèche $3\text{--}5$ mois.
        \item Savane soudanaenne (Nord) : $P \approx \SI{800}{--}\SI{1200}{mm}$, saison sèche $5\text{--}7$ mois.
        \item Steppe sahélienne (Extrême-Nord) : $P < \SI{800}{mm}$, saison sèche $> 7$ mois.
    \end{itemize}
\end{methode}

\begin{exoresolu}[Analyse de fiches climatiques et végétales de deux localités camerounaises]
    \textbf{Énoncé :} Le tableau ci-dessous présente des données climatiques et végétales pour deux localités A et B du Cameroun.
    \begin{center}
    \begin{tabularx}{\linewidth}{|l|X|X|}\hline
    \rowcolor{popblueL} \textbf{Paramètre} & \textbf{Localité A} & \textbf{Localité B} \\ \hline
    Pluviométrie annuelle & $\SI{1850}{mm}$ & $\SI{950}{mm}$ \\ \hline
    Température moyenne annuelle & $\SI{24,5}{\celsius}$ & $\SI{27,0}{\celsius}$ \\ \hline
    Durée de la saison sèche & $\SI{1,5}{mois}$ & $\SI{6}{mois}$ \\ \hline
    Végétation dominante & Arbres $> \SI{30}{m}$, canopée fermée, lianes, épiphytes, feuilles larges persistantes & Arbres $\SI{10}{--}\SI{15}{m}$ (Isoberlinia), écorce épaisse, feuilles caduques, strate herbacée dense (Andropogon) \\ \hline
    Type de sol & Ferrallitique lessivé, humifère en surface & Ferrugineux tropical, peu humifère \\ \hline
    \end{tabularx}
    \end{center}
    \textbf{Questions :}
    \begin{enumerate}
        \item Identifier le type d'écosystème pour chaque localité.
        \item Justifier votre réponse en utilisant \textbf{deux} paramètres climatiques et \textbf{deux} caractéristiques végétales pour chaque localité.
        \item Citer une espèce animale emblématique attendue dans chaque écosystème.
    \end{enumerate}
    \tcblower
    \textbf{Solution détaillée :}
    \begin{enumerate}
        \item \textbf{Localité A :} \cle{Forêt dense humide} (type forêt équatoriale / guinéenne).\\
              \textbf{Localité B :} \cle{Savane boisée} (type savane soudanaenne).
        \item \textbf{Justification Localité A (Forêt) :}
              \begin{itemize}
                  \item \textit{Climatiques :} Pluviométrie élevée ($\SI{1850}{mm} > \SI{1500}{mm}$) et saison sèche très courte ($\SI{1,5}{mois} < 2$ mois) : conditions thermohygrométriques favorables à la croissance continue.
                  \item \textit{Végétales :} Canopée fermée (strates superposées, fort LAI) limitant la lumière au sol ; feuilles mésophylles persistantes (adaptation à l'humidité constante, pas de stress hydrique).
              \end{itemize}
              \textbf{Justification Localité B (Savane) :}
              \begin{itemize}
                  \item \textit{Climatiques :} Pluviométrie modérée ($\SI{950}{mm}$) et saison sèche longue ($\SI{6}{mois}$) : stress hydrique marqué sélectionnant des espèces xérophiles.
                  \item \textit{Végétales :} Arbres à écorce épaisse (protection feu/feu) et feuilles caduques microphylles (réduction transpiration en saison sèche) ; strate herbacée continue dominée par des graminées pérennes (Andropogon) adaptées au feu.
              \end{itemize}
        \item \textbf{Espèces emblématiques :}
              \begin{itemize}
                  \item Forêt (A) : Gorille des plaines de l'Ouest (\emph{Gorilla gorilla}), Éléphant de forêt (\emph{Loxodonta cyclotis}), Chimpanzé (\emph{Pan troglodytes}), Bongo (\emph{Tragelaphus eurycerus}), Grand Calao (\emph{Ceratogymna bucinator}).
                  \item Savane (B) : Éléphant de savane (\emph{Loxodonta africana}), Girafe (\emph{Giraffa camelopardalis peralta} -- Nord), Kob de Buffon (\emph{Kobus kob}), Lion (\emph{Panthera leo}), Autruche (\emph{Struthio camelus} -- Nord), Tisserin (\emph{Ploceus cucullatus}).
              \end{itemize}
    \end{enumerate}
    \textbf{Conclusion :} La localité A correspond à la zone forestière du Sud-Cameroun (ex: Ebolowa, Sangmélima) et la localité B à la zone soudanaenne du Nord-Cameroun (ex: Garoua, Maroua). La pluviométrie et la durée de la saison sèche sont les facteurs climatiques discriminants majeurs.
\end{exoresolu}

\begin{methode}[Construire et interpréter une chaîne / un réseau trophique (Interdépendance)]
    \textbf{Objectif :} À partir d'une liste d'êtres vivants et de leurs régimes alimentaires, représenter les relations proies-prédateurs et expliquer la circulation de la matière et de l'énergie.
    \begin{enumerate}
        \item \textbf{Classer les êtres vivants par niveau trophique :}
              \begin{itemize}
                  \item \textbf{Niveau 1 (Producteurs) :} Êtres chlorophylliens (arbres, herbes, phytoplancton) $\rightarrow$ Photosynthèse ($\ce{CO2 + H2O ->[Lumiere] MO + O2}$).
                  \item \textbf{Niveau 2 (Consommateurs primaires / Herbivores) :} Mangeurs de producteurs.
                  \item \textbf{Niveau 3 (Consommateurs secondaires / Carnivores) :} Mangeurs d'herbivores.
                  \item \textbf{Niveau 4 (Consommateurs tertiaires / Superprédateurs) :} Mangeurs de carnivores.
                  \item \textbf{Décomposeurs (Bactéries, Champignons) :} Recyclent la M.O. de \textbf{tous} les niveaux $\rightarrow$ Minéralisation ($\ce{MO -> Sels + CO2 + H2O}$).
              \end{itemize}
        \item \textbf{Tracer les flèches :} Flèche \textbf{partant de la proie vers le prédateur} (sens du \textbf{transfert de matière et d'énergie}). \textit{Astuce : « La flèche va dans le ventre de celui qui mange ».}
        \item \textbf{Identifier les relations :} Compétition (même proie), Prédateur-Proie, Décomposition.
        \item \textbf{Analyser les conséquences d'une perturbation :} Disparition d'un maillon $\rightarrow$ Effet en cascade (pullulation proie, famine prédateur).
    \end{enumerate}
    \textbf{Rappel :} L'énergie se perd (chaleur) à chaque transfert (rendement $\approx 10\%$), la matière se recycle (cycle du carbone, de l'azote, de l'eau).
\end{methode}

\begin{exoresolu}[Construction et analyse d'un réseau trophique en savane boisée]
    \textbf{Énoncé :} Dans une savane boisée de l'Adamaoua, on observe les êtres vivants suivants et leurs régimes :
    \begin{itemize}
        \item Herbes (Andropogon, Hyparrhenia), Arbres (Isoberlinia, Acacia), Arbustes.
        \item Sauterelles, Chenilles (larves de lépidoptères), Rongeurs (rats des champs, écureuils), Kob de Buffon, Phacochère.
        \item Lézards, Mangouste, Python, Aigle martial, Hyène tachetée, Chacal.
        \item Bactéries du sol, Champignons (termites cultivant des champignons).
    \end{itemize}
    \textbf{Régimes :} Sauterelles/Chenilles $\rightarrow$ Herbes/Feuilles ; Rongeurs $\rightarrow$ Graines/Racines/Herbes ; Kob/Phacochère $\rightarrow$ Herbes/Feuilles/Tubercules ; Lézards $\rightarrow$ Insectes ; Mangouste $\rightarrow$ Rongeurs/Insectes/Œufs/Reptiles ; Python $\rightarrow$ Rongeurs/Kob (jeunes)/Mangouste ; Aigle martial $\rightarrow$ Rongeurs/Kob (jeunes)/Mangouste/Reptiles ; Hyène/Chacal $\rightarrow$ Charognes/Kob (malades/jeunes)/Déchets ; Décomposeurs $\rightarrow$ M.O. morte (tous niveaux).
    \textbf{Questions :}
    \begin{enumerate}
        \item Classer ces êtres vivants dans un tableau à 5 colonnes : Producteurs, Consom. 1\up{er}, Consom. 2\up{nd}, Consom. 3\up{me}, Décomposeurs.
        \item Représenter \textbf{deux chaînes trophiques} distinctes aboutissant à l'Aigle martial.
        \item L'homme pratique le braconnage intensif sur le Kob de Buffon. Prédire les conséquences à court et moyen terme sur : (a) la strate herbacée, (b) l'Aigle martial, (c) la Hyène.
        \item Quel est le rôle des termites et champignons dans cet écosystème ?
    \end{enumerate}
    \tcblower
    \textbf{Solution détaillée :}
    \begin{enumerate}
        \item \textbf{Tableau de classification :}
        \begin{center}
        \begin{tabularx}{\linewidth}{|l|X|X|X|X|X|}\hline
        \rowcolor{popgreenL} \textbf{Producteurs} & \textbf{Cons. 1\up{er} (Herbivores)} & \textbf{Cons. 2\up{nd} (Carniv. 1\up{er})} & \textbf{Cons. 3\up{me} (Carniv. 2\up{nd})} & \textbf{Cons. 4\up{me} (Superpréd.)} & \textbf{Décomposeurs} \\ \hline
        Herbes, Arbres, Arbustes & Sauterelles, Chenilles, Rongeurs, Kob, Phacochère & Lézards, Mangouste & Python, Aigle martial, Hyène, Chacal & (Aigle martial, Hyène -- sommet) & Bactéries, Champignons, Termites \\ \hline
        \end{tabularx}
        \end{center}
        \textit{Note :} L'Aigle martial et la Hyène peuvent occuper le niveau 3 ou 4 selon la proie (ex: manger un python = niveau 4). On les place au plus haut niveau atteint.
        \item \textbf{Chaînes trophiques vers l'Aigle martial :}
              \begin{enumerate}
                  \item Herbes $\rightarrow$ Rongeurs $\rightarrow$ Aigle martial.
                  \item Herbes $\rightarrow$ Kob (jeune) $\rightarrow$ Aigle martial.
                  \item Arbres $\rightarrow$ Chenilles $\rightarrow$ Lézards $\rightarrow$ Mangouste $\rightarrow$ Aigle martial. (Plus longue, montre la complexité).
              \end{enumerate}
        \item \textbf{Conséquences du braconnage sur le Kob :}
              \begin{itemize}
                  \item \textbf{(a) Strate herbacée :} Diminution de la pression de broutage $\rightarrow$ \textbf{Pullulation / Biomasse accrue} de la strate herbacée (risque incendie accru en saison sèche).
                  \item \textbf{(b) Aigle martial :} Perte d'une proie majeure (jeunes kobs) $\rightarrow$ \textbf{Diminution de l'effectif} (famine, baisse reproduction, émigration) ou changement de régime (pression accrue sur rongeurs/mangoustes).
                  \item \textbf{(c) Hyène :} Perte de charognes (kobs braconnés non récupérés ou tués par hyènes) et proies vivantes $\rightarrow$ \textbf{Baisse de l'effectif} ou rapprochement des villages (conflit homme-faune).
              \end{itemize}
        \item \textbf{Rôle des termites/champignons :} \textbf{Décomposeurs / Recycleurs}. Ils dégradent la cellulose/lignine (bois mort, fèces, herbes sèches) que peu d'autres dégradent. Ils minéralisent la M.O. $\rightarrow$ libèrent sels minéraux ($\ce{NH4+}$, $\ce{NO3-}$, $\ce{PO4^{3-}}$, $\ce{K+}$) $\rightarrow$ \textbf{Fertilité du sol} pour les producteurs. Les termites aèrent aussi le sol (galeries).
    \end{enumerate}
    \textbf{Conclusion :} Le braconnage déséquilibre le réseau trophique : pullulation des producteurs, effondrement des prédateurs spécialisés, risque d'incendie et conflit homme-faune. La conservation des herbivores régule la végétation et maintient les prédateurs.
\end{exoresolu}

\begin{methode}[Analyser l'impact d'une activité humaine sur un écosystème]
    \textbf{Objectif :} À partir d'un document (texte, photo satellite, graphique d'évolution de la couverture forestière), identifier l'activité, le mécanisme de destruction et les conséquences écologiques.
    \begin{enumerate}
        \item \textbf{Identifier l'activité :} Déforestation (exploitation forestière, agriculture sur brûlis, plantations), Feu de brousse (chasse, pâturage, mégot), Braconnage (viande de brousse, ivoire, écailles de pangolin), Urbanisation/Infrastructures, Pollution (pesticides, déchets plastiques, hydrocarbures).
        \item \textbf{Expliquer le mécanisme physique/biologique :}
              \begin{itemize}
                  \item \textit{Déforestation :} Suppression couvert $\rightarrow$ Érosion hydrique/éolienne (sol nu), Perte habitat, Rupture cycle de l'eau (évapotranspiration $\downarrow$, ruissellement $\uparrow$), Émission $\ce{CO2}$ (stock carbone libéré).
                  \item \textit{Feu de brousse :} Combustion biomasse $\rightarrow$ $\ce{CO2}$, $\ce{CH4}$, particules ; Destruction graines/jeunes plants/faune lente ; Sélection espèces pyrophiles (herbes) $\rightarrow$ Régression ligneuse (savanisation).
                  \item \textit{Braconnage :} Prélèvement $>$ Taux de renouvellement $\rightarrow$ Déclin populations $\rightarrow$ Rupture chaînes trophiques (espèces clés).
              \end{itemize}
        \item \textbf{Citer les conséquences :} Perte biodiversité (espèces endémiques), Dégradation sols (stérilité, latéritisation), Changement microclimat (sécheresse locale), Perte services écosystémiques (bois, plantes médicinales, régulation climat, tourisme).
        \item \textbf{Proposer une mesure d'atténuation adaptée.}
    \end{enumerate}
\end{methode}

\begin{exoresolu}[Étude de cas : L'exploitation forestière et l'agriculture sur brûlis au Sud-Cameroun]
    \textbf{Énoncé :} Un rapport de l'Observatoire des Forêts d'Afrique Centrale (OFAC) indique qu'entre 2000 et 2020, la région du Sud a perdu $12\%$ de sa forêt dense primaire. Les causes principales citées sont l'exploitation forestière industrielle (grumes pour export) et l'agriculture itinérante sur brûlis (cacao, vivriers). Une image satellite montre des « trous » dans le couvert forestier le long des routes forestières.
    \textbf{Questions :}
    \begin{enumerate}
        \item Citer les deux activités humaines responsables et les distinguer (mécanisme, acteurs).
        \item Expliquer pourquoi l'ouverture de routes forestières accélère la déforestation (effet « osmose » ou « fishbone »).
        \item L'exploitation forestière « sélective » ne prélève que quelques arbres/ha. Pourquoi cause-t-elle tout de même une perte massive de biodiversité ?
        \item Calculer : Si la forêt stockait $\SI{250}{t\,C/ha}$ (carbone dans la biomasse aérienne) et que $50\%$ du carbone est émis lors de la conversion en champ de cacao (brûlis + décomposition résidus), quelle quantité de $\ce{CO2}$ est émise par hectare converti ? ($M_{\ce{C}}=\SI{12}{g/mol}$, $M_{\ce{O}}=\SI{16}{g/mol}$).
        \item Proposer deux mesures concrètes pour concilier exploitation du cacao et conservation forestière au Cameroun.
    \end{enumerate}
    \tcblower
    \textbf{Solution détaillée :}
    \begin{enumerate}
        \item \textbf{Activité 1 : Exploitation forestière industrielle.} Mécanisme : Coupe sélective d'essences commerciales (Ayous, Sapelli, Azobé). Acteurs : Sociétés concessionnaires (souvent étrangères), État (permis). Impact direct : Ouverture trouées, pistes, dommages collatéraux (arbres abîmés).
              \textbf{Activité 2 : Agriculture itinérante sur brûlis (cacaoculture vivrière).} Mécanisme : Abattis-brûlis de la forêt secondaire ou primaire $\rightarrow$ Culture 2-3 ans $\rightarrow$ Jachère courte (insuffisante) $\rightarrow$ Nouveau défrichage. Acteurs : Paysans locaux, migrants. Impact direct : Conversion définitive forêt $\rightarrow$ agroécosystème simplifié.
        \item \textbf{Effet « Fishbone » (arête de poisson) :} Les routes pénètrent la forêt (artère principale). Les pistes secondaires (branches) permettent l'accès aux paysans (agriculture) et braconniers. La déforestation suit le réseau routier. La route fragmente l'habitat (effet de lisière : vent, lumière, chasse) et isole les populations animales (barrière génétique).
        \item \textbf{Impact exploitation sélective :}
              \begin{itemize}
                  \item Dommages collatéraux : Pour 1 arbre prélevé, 10 à 20 arbres détruits (chute, treuillage).
                  \item Ouverture canopée $\rightarrow$ Invasion espèces pionnières/lianes $\rightarrow$ Régression espèces climax (ombre).
                  \item Pistes = accès braconnage (vidange faunistique / « empty forest syndrome »).
                  \item Fragmentation : Populations isolées $\rightarrow$ Consanguinité, extinction locale.
              \end{itemize}
        \item \textbf{Calcul émission $\ce{CO2}$ :}
              \begin{align*}
                  \text{Carbone stocké initial} &= \SI{250}{t\,C/ha} \\
                  \text{Carbone émis (50\%)} &= \SI{125}{t\,C/ha} \\
                  \text{Masse molaire } \ce{CO2} &= 12 + 2\times 16 = \SI{44}{g/mol} \\
                  \text{Masse molaire C} &= \SI{12}{g/mol} \\
                  \text{Masse } \ce{CO2} &= \SI{125}{t\,C} \times \frac{44}{12} = 125 \times 3,666... \\
                  &\approx \mathbf{\SI{458,3}{t\,CO2/ha}} \quad (\text{arrondi } \SI{458}{t\,CO2/ha})
              \end{align*}
        \item \textbf{Mesures de conciliation (Agroforesterie / Zéro déforestation) :}
              \begin{enumerate}
                  \item \textbf{Cacao sous couvert (Agroforesterie) :} Planter le cacaoyer sous l'ombrage d'arbres forestiers conservés (espèces à usage : bois d'œuvre, fruits, engrais verts type \emph{Albizia}, \emph{Gliricidia}). Maintient carbone, biodiversité, fertilité, revenus diversifiés.
                  \item \textbf{Intensification durable / Zéro déforestation :} Améliorer rendements ha existants (variétés améliorées, engrais organiques, taille, lutte bio) pour stopper l'extension. Certification (Rainforest Alliance, Fairtrade) + primes pour zéro déforestation. Sécurisation foncière pour investir long terme.
              \end{enumerate}
    \end{enumerate}
    \textbf{Conclusion :} La déforestation au Sud-Cameroun est un couplage route-industrie-agriculture. L'agroforesterie cacaoyère est la solution technique majeure pour sortir de l'abattis-brûlis tout en maintenant la production nationale (1\up{er} exportateur africain).
\end{exoresolu}

\begin{methode}[Planifier et justifier une opération de reboisement (TP N°10)]
    \textbf{Objectif :} Décrire les étapes techniques du reboisement, du choix de l'essence au suivi, en justifiant chaque choix par l'écologie de l'espèce et le but visé (bois d'œuvre, énergie, protection sol, biodiversité).
    \begin{enumerate}
        \item \textbf{Diagnostic du site :} Type de sol (profondeur, texture, pH), Climat (pluviométrie, saison sèche), Topographie (pente, exposition), Pression anthropique (pâturage, feu), Ancien usage.
        \item \textbf{Choix de l'essence (Espèce adaptée = « Bonne essence au bon endroit ») :}
              \begin{itemize}
                  \item \textit{Zone dégradée/sol nu/pente forte :} Espèces pionnières, rustiques, fixatrices d'azote (\emph{Acacia mangium}, \emph{Leucaena leucocephala}, \emph{Calliandra calothyrsus}) $\rightarrow$ Reconstruction fertilité.
                  \item \textit{Zone forestière enrichissement :} Essences commerciales locales à croissance rapide (\emph{Terminalia superba} / Fraké, \emph{Triplochiton scleroxylon} / Ayous, \emph{Entandrophragma} / Acajou).
                  \item \textit{Zone sahélienne (Nord) :} \emph{Acacia senegal} (gomme arabique), \emph{Balanites aegyptiaca}, \emph{Azadirachta indica} (Neem), \emph{Prosopis juliflora} (attention invasive).
              \end{itemize}
        \item \textbf{Production plants (Pépinière) :} Graines (provenance certifiée) ou boutures. Substrat (terre de surface + compost + sable). Ombrière (50\% ombre $\rightarrow$ durcissement plein soleil 1 mois avant plantation). Arrosage régulier. \textbf{Critère qualité plant :} Hauteur $\SI{30}{--}\SI{50}{cm}$, collet $\varnothing > \SI{5}{mm}$, système racinaire fibré (non en chignon), sans maladie.
        \item \textbf{Préparation du terrain :} Débroussaillage, traçage (alignement, écartement), trou de plantation ($\SI{40}{\cm}\times\SI{40}{\cm}\times\SI{40}{\cm}$ standard, séparer terre superficielle / profonde). \textit{Écartement :} $\SI{3}{\m}\times\SI{3}{\m}$ (1111 plants/ha) bois d'œuvre ; $\SI{2}{\m}\times\SI{2}{\m}$ (2500 plants/ha) bois énergie/protection ; $\SI{4}{\m}\times\SI{4}{\m}$ enrichissement.
        \item \textbf{Plantation (Début saison des pluies : mars-avril Sud, juin-juillet Nord) :} Repiquage motte intacte, collet au niveau sol, tassement modéré, arrosage si possible (bassin).
        \item \textbf{Entretien / Suivi (Critique 1-3 ans) :} Désherbage manuel (rayon $\SI{50}{\cm}$) 2-3 fois/an, protection feu (pare-feu $\SI{5}{--}\SI{10}{\m}$), protection pâturage (clôture, gardien), remplacement plants morts (repiquage « beating up » $\rightarrow$ taux de réussite visé $> 80\%$), élagage de formation (bois d'œuvre).
    \end{enumerate}
    \textbf{Indicateurs de succès :} Taux de reprise ($> 80\%$), Hauteur moyenne à 1 an, Diamètre collet, Recouvrement du sol.
\end{methode}

\begin{exoresolu}[Conduite d'un projet de reboisement sur un versant dégradé de l'Adamaoua]
    \textbf{Énoncé :} Une ONG locale souhaite reboiser $\SI{5}{hectares}$ d'un versant nu, érodé (ravines), sol ferrugineux tropical lessivé, pluviométrie $\SI{1400}{mm/an}$, saison sèche $\SI{5}{mois}$, pâturage libre intense. L'objectif double est : (1) stopper l'érosion et restaurer la fertilité, (2) fournir du bois énergie aux villages riverains à horizon 5-7 ans.
    \textbf{Questions :}
    \begin{enumerate}
        \item Proposer \textbf{deux essences principales} adaptées, justifier par leurs traits écologiques (racines, azote, croissance, feu).
        \item Calculer le nombre de plants à commander pour $\SI{5}{ha}$ avec un écartement $\SI{2}{\m}\times\SI{2}{\m}$, en prévoyant un taux de mortalité de $20\%$ au repiquage (plants de remplacement).
        \item Décrire le calendrier des opérations sur 12 mois (mois par mois) pour la région de l'Adamaoua (saison des pluies mai-octobre).
        \item L'ONG a un budget limité. Proposer une technique low-cost pour protéger les jeunes plants du pâturage sans clôture métallique coûteuse.
        \item Citer deux indicateurs de suivi à mesurer à 6 mois et à 2 ans pour valider la réussite.
    \end{enumerate}
    \tcblower
    \textbf{Solution détaillée :}
    \begin{enumerate}
        \item \textbf{Essence 1 : \emph{Acacia mangium} (ou \emph{A. auriculiformis}).}
              \textit{Justification :} Pionnière héliophile, croissance très rapide (bois énergie 5-7 ans), \textbf{fixatrice d'azote} (bactéries \emph{Bradyrhizobium} : enrichit sol pauvre), système racinaire pivotant profond + latéraux (stabilise sol, puise eau nappe), tolère sols acides/ferrugineux, résiste feu (écorce épaisse à l'âge adulte), bois calorifique élevé.
              \textbf{Essence 2 : \emph{Gliricidia sepium} (ou \emph{Leucaena leucocephala} -- attention toxicité graines bétail).}
              \textit{Justification :} Légumineuse fixatrice d'azote, multiplication facile par \textbf{boutures} (pas de pépinière coûteuse, gain temps), repoussent vigoureuses après coupe (taillis), feuillage fourrager (valeur ajoutée si pâturage contrôlé), racines denses superficielles (contrôle érosion surface). \textit{Association :} Lignes \emph{Acacia} (bois d'œuvre/énergie long terme) + Haies \emph{Gliricidia} (contour lignes niveau, fertilité, piquets).
        \item \textbf{Calcul plants :}
              \begin{align*}
                  \text{Densité nominale} &= \frac{\SI{10000}{m^2/ha}}{\SI{2}{m} \times \SI{2}{m}} = \SI{2500}{plants/ha} \\
                  \text{Total nominal 5 ha} &= 2500 \times 5 = \SI{12500}{plants} \\
                  \text{Taux mortalité prévu} &= 20\% \\
                  \text{Plants de remplacement} &= 12500 \times 0,20 = \SI{2500}{plants} \\
                  \text{Total à commander} &= 12500 + 2500 = \mathbf{\SI{15000}{plants}}
              \end{align*}
        \item \textbf{Calendrier Adamaoua (Pluies mai-oct) :}
              \begin{itemize}
                  \item \textbf{Nov-Déc :} Récolte graines/boutures, installation pépinière (ombrière, substrat), semis/bouturage.
                  \item \textbf{Janv-Mars :} Entretien pépinière (arrosage, désherbage, fertilisation), durcissement (levée ombrière mars).
                  \item \textbf{Avril :} Préparation terrain (débroussaillage, traçage, trous), \textbf{pare-feu} ($\SI{5}{\m}$ large brûlé contrôlé ou débroussaillé).
                  \item \textbf{Mai (1\up{re} pluies) :} \textbf{Plantation} (repiquage motte), arrosage d'appoint si trouée sèche.
                  \item \textbf{Juin-Juillet-Août :} \textbf{Désherbage 1 \& 2} (manuel, rayon $\SI{50}{\cm}$), surveillance feu/pâturage, remplacement morts (beating up).
                  \item \textbf{Sept-Oct :} Désherbage 3 (fin saison), dernier remplacement.
                  \item \textbf{Nov-Avril (Saison sèche) :} Surveillance pare-feu, protection pâturage, pas d'arrosage (racines profondes).
              \end{itemize}
        \item \textbf{Protection low-cost anti-pâturage :}
              \begin{itemize}
                  \item \textbf{Haie vive « morte » ou vive :} Plantation dense ($\SI{50}{\cm}$) de \emph{Jatropha curcas} (toxique, non brouté) ou \emph{Euphorbia} spp ou bambous épineux le long de la limite + brins secs (branchages épineux d'acacia) tissés sur piquets bois (clôture « bush fence »). Coût quasi nul (main d'œuvre locale).
                  \item \textbf{Contrat social / Gestion communautaire :} Accord avec éleveurs (chefs traditionnels) : interdiction pâturage 3 ans contre droit de coupe herbe/foin sur pare-feu ou jachères voisines. Gardien unique payé par village (moins cher que clôture).
              \end{itemize}
        \item \textbf{Indicateurs de suivi :}
              \begin{itemize}
                  \item \textbf{À 6 mois (Fin 1\up{re} saison) :} \textbf{Taux de reprise} (nombre plants vivants / plants mis en place $\times 100$). Objectif $> 80\%$. Hauteur moyenne.
                  \item \textbf{À 2 ans (Fin 2\up{e} saison) :} \textbf{Hauteur moyenne} (objectif $> \SI{2,5}{\m}$ pour \emph{Acacia}), \textbf{Diamètre au collet}, \textbf{Taux de recouvrement du sol} (\% sol nu vs litière/herbes). Absence de ravines actives.
              \end{itemize}
    \end{enumerate}
    \textbf{Conclusion :} Le succès repose sur le choix d'espèces fixatrices d'azote rustiques (\emph{Acacia}, \emph{Gliricidia}), une plantation précoce (mai), un désherbage rigoureux les 2 premières années et une gouvernance locale pour exclure le feu et le bétail.
\end{exoresolu}

\begin{methode}[Argumenter pour la conservation de la biodiversité (Développement durable)]
    \textbf{Objectif :} Rédiger un texte argumenté (plaidoyer, rapport, réponse à une question de synthèse) justifiant la création d'une aire protégée ou une mesure de conservation, en articulant valeurs intrinsèques, services écosystémiques et enjeux socio-économiques locaux.
    \begin{enumerate}
        \item \textbf{Définir l'enjeu :} Espèce/Écosystème menacé, Pression identifiée (braconnage, habitat loss, climat).
        \item \textbf{Argument 1 : Valeur intrinsèque / Éthique :} Droit à l'existence, patrimoine naturel national/mondial (UNESCO), responsabilité intergénérationnelle.
        \item \textbf{Argument 2 : Services écosystémiques (Valeur utilitaire) :}
              \begin{itemize}
                  \item \textit{Approvisionnement :} Bois, plantes médicinales (\emph{Prunus africana}, \emph{Voacanga}), aliments (champignons, chenilles, fruits), ressources génétiques (amélioration variétés).
                  \item \textit{Régulation :} Climat (puits carbone), Eau (filtration, régime hydrologique rivières Sanaga, Bénoué, Logone), Sols (anti-érosion), Pollinisation (agriculture), Régulation maladies (dilution pathogènes).
                  \item \textit{Culturel :} Tourisme (revenus devises, emplois guides/hôtels), Sacralité (forêts sacrées), Éducation/Recherche.
              \end{itemize}
        \item \textbf{Argument 3 : Approche participative / Droits humains :} Implication populations riveraines (Gestion communautaire, ZICGC - Zones d'Intérêt Cynégétique à Gestion Communautaire), Partage bénéfices (redevances, emplois éco-gardes), Alternative revenus (apiculture, écotourisme, agroforesterie) $\rightarrow$ Réduction braconnage/feu.
        \item \textbf{Conclusion synthétique :} « Conserver n'est pas interdire, c'est gérer durablement pour le présent et le futur. »
    \end{enumerate}
    \textbf{Mots-clés à utiliser :} \cle{Services écosystémiques}, \cle{Développement durable}, \cle{Participation communautaire}, \cle{Valeur d'option}, \cle{Équité intergénérationnelle}.
\end{methode}

\begin{exoresolu}[Rédaction d'un plaidoyer pour la création d'un corridor écologique entre le Parc National de la Bénoué et la Réserve de Faune du Faro (Nord-Cameroun)]
    \textbf{Énoncé :} Les populations de grands mammifères (Éléphant, Girafe, Lycaon, Éland de Derby) sont isolées dans deux aires protégées séparées par une zone de chasse villageoise (ZIC 4) et des couloirs de transhumance. Le gouvernement envisage de classer un « Corridor Écologique » de $\SI{50}{km}$ de large. Vous êtes consultant pour le MINEPDED. Rédigez les arguments scientifiques et socio-économiques clés (max 20 lignes) pour justifier ce classement auprès des populations locales et des décideurs.
    \tcblower
    \textbf{Solution détaillée (Modèle de réponse type BEPC - Synthèse) :}
    \textbf{Plaidoyer pour le Corridor Bénoué-Faro}
    \par\medskip
    \textbf{1. Urgence écologique : Connectivité génétique et survie des espèces clés.}
    L'isolement des populations d'éléphants (\emph{Loxodonta africana}), de girafes (\emph{Giraffa camelopardalis peralta}) et de lycaons (\emph{Lycaon pictus}) dans les deux aires protégées entraîne une \textbf{dérive génétique} et une \textbf{consanguinité} (effet d'îlot), menaçant leur viabilité à long terme. Le corridor restaure la \cle{connectivité paysagère}, permettant les migrations saisonnières (eau, pâturages) et les flux géniques indispensables à l'adaptation au changement climatique.
    \par\medskip
    \textbf{2. Sécurisation des services écosystémiques vitaux pour le Nord.}
    Ce corridor protège les \textbf{bassins versants} des rivières Bénoué et Faro (affluents du Niger). La couverture végétale maintient l'\textbf{infiltration de l'eau}, régule les crues (protection rizières de la plaine de la Bénoué) et la recharge des nappes (forage eau potable). Il stocke du \textbf{carbone} (atténuation changement climatique) et abrite des \textbf{pollinisateurs} et auxiliaires de cultures (chauves-souris, oiseaux) pour l'agriculture cotonnière et vivrière environnante.
    \par\medskip
    \textbf{3. Développement économique local par la valorisation durable.}
    Le classement n'interdit pas l'homme, il \textbf{organise l'espace} : Zonage (conservation stricte, usage multiple, pâturage contrôlé). Il ouvre la voie à l'\textbf{écotourisme} (safari photo, ornithologie) générant des emplois locaux (guides, hébergement) et des redevances pour les communes (écoles, santé). L'apiculture et la cueillette durable (karité, néré, tamarin) offrent des revenus immédiats aux femmes et jeunes, alternative au braconnage et à l'agriculture extensive sur brûlis.
    \par\medskip
    \textbf{4. Paix sociale et sécurité : Gestion concertée des ressources.}
    La formalisation du corridor par une \textbf{Convention de Gestion Participative} (État, Communes, Chefs traditionnels, Éleveurs transhumants, ONG) définit des règles claires : couloirs de transhumance balisés, points d'eau partagés, comités de surveillance mixtes. Cela réduit les conflits agriculteurs-éleveurs et l'insécurité liée aux groupes armés qui exploitent le vide juridique des zones non gérées.
    \par\medskip
    \textbf{Conclusion :} Ce corridor est un investissement « Nature-based Solution » : il assure la résilience de la biodiversité \emph{et} la sécurité alimentaire/hydrique des populations du Nord-Cameroun. Refuser ce classement, c'est accepter l'érosion irréversible du capital naturel dont dépend l'économie régionale.
\end{exoresolu}

\begin{methode}[Interpréter des données de suivi écologique (Tableaux, Graphiques, Indices)]
    \textbf{Objectif :} Lire un graphique d'évolution d'effectifs, un tableau d'inventaires (Indice Kilométrique d'Abondance - IKA, Pièges photographiques), ou un calcul de richesse spécifique ($S$) / indice de Shannon ($H'$) simplifié, et en tirer une conclusion sur l'état de conservation.
    \begin{enumerate}
        \item \textbf{Lire les axes / en-têtes :} Variable (Effectif, IKA, $S$, $H'$), Unités, Période, Sites comparés (Protégé vs Non protégé, Avant/Après perturbation).
        \item \textbf{Identifier les tendances :} Croissance, déclin, stabilité, seuils critiques.
        \item \textbf{Calculer si nécessaire :}
              \begin{itemize}
                  \item \textbf{Richesse spécifique $S$ :} Nombre total d'espèces recensées.
                  \item \textbf{Indice de Shannon $H' = -\sum (p_i \ln p_i)$} (souvent donné ou à calculer simple : $p_i = n_i / N$). $H'$ haut = diversité forte + équitabilité.
                  \item \textbf{IKA (Indice Kilométrique d'Abondance) :} $\frac{\text{Nombre d'individus observés (ou indices : crottes, traces)}}{\text{Kilomètres parcourus}}$ ou $\frac{\text{Contacts}}{\text{Effort (j-pièges)}} \times 100$. Permet comparaison espace/temps.
              \end{itemize}
        \item \textbf{Corréler avec pressions :} Braconnage (↓ grands mammifères, ↑ espèces tolérantes), Feux (↓ espèces forestières, ↑ espèces savanicoles), Reboisement (↑ richesse, retour espèces exigeantes).
        \item \textbf{Conclure avec chiffres :} « L'IKA de l'éléphant passe de 1,2 à 0,3 ind/km en 10 ans (-75\%), corrélé à la hausse des saisies d'ivoire. »
    \end{enumerate}
\end{methode}

\begin{exoresolu}[Analyse de données de suivi faunique par pièges photographiques (Camera Traps)]
    \textbf{Énoncé :} Une étude par pièges photographiques (Camera Traps) a été menée sur 60 jours dans deux sites du Massif du Mbam-et-Djerem : Site P (Parc National, protection stricte) et Site Z (Zone de Chasse Villageoise ZICGC, exploitation cynégétique régulée + agriculture). Résultats pour 5 espèces cibles :
    \begin{center}
    \begin{tabular}{|l|cc|cc|}\hline
    \rowcolor{poporangeL} \textbf{Espèce} & \multicolumn{2}{c|}{\textbf{Site P (Parc)}} & \multicolumn{2}{c|}{\textbf{Site Z (ZICGC)}} \\ \hline
    & \textbf{Contacts} & \textbf{IKA (contacts/100 j-pièges)} & \textbf{Contacts} & \textbf{IKA (contacts/100 j-pièges)} \\ \hline
    Éléphant de forêt & 120 & 2,0 & 15 & 0,25 \\ \hline
    Bongo & 45 & 0,75 & 5 & 0,08 \\ \hline
    Céphalophe à dos jaune & 210 & 3,5 & 180 & 3,0 \\ \hline
    Potamochère & 80 & 1,33 & 95 & 1,58 \\ \hline
    Léopard & 12 & 0,20 & 0 & 0,00 \\ \hline
    \textbf{Total Effort (j-pièges)} & \multicolumn{2}{c|}{6000} & \multicolumn{2}{c|}{6000} \\ \hline
    \end{tabular}
    \end{center}
    \textbf{Questions :}
    \begin{enumerate}
        \item Vérifier le calcul de l'IKA pour l'Éléphant au Site P.
        \item Comparer la richesse spécifique ($S$) et l'abondance relative (IKA total) entre les deux sites.
        \item Classer les 5 espèces en \textbf{espèces sensibles à la perturbation anthropique} et \textbf{espèces tolérantes / opportunistes} en vous basant sur les ratios IKA (Site Z / Site P).
        \item Le Léopard n'a pas été photographié au Site Z. Peut-on conclure à son extinction locale ? Justifier (biais méthode).
        \item Proposer une hypothèse expliquant pourquoi la Potamochère a un IKA légèrement supérieur au Site Z.
    \end{enumerate}
    \tcblower
    \textbf{Solution détaillée :}
    \begin{enumerate}
        \item \textbf{Vérification IKA Éléphant Site P :}
              $\text{IKA} = \frac{\text{Contacts}}{\text{Effort (j-pièges)}} \times 100 = \frac{120}{6000} \times 100 = 2,0$. \textbf{Calcul exact.}
        \item \textbf{Richesse spécifique $S$ :} Site P = 5 espèces. Site Z = 4 espèces (Léopard absent). $\rightarrow$ \textbf{Richesse plus élevée au Site P (protection stricte).}
              \textbf{Abondance relative totale (Somme IKA) :}
              Site P : $2,0 + 0,75 + 3,5 + 1,33 + 0,20 = \mathbf{7,78}$.
              Site Z : $0,25 + 0,08 + 3,0 + 1,58 + 0 = \mathbf{4,91}$.
              $\rightarrow$ \textbf{Abondance globale plus forte au Site P.}
        \item \textbf{Calcul Ratio $R = \frac{\text{IKA Site Z}}{\text{IKA Site P}}$ :}
              \begin{itemize}
                  \item Éléphant : $0,25 / 2,0 = \mathbf{0,125}$ (Effondrement $87,5\%$).
                  \item Bongo : $0,08 / 0,75 = \mathbf{0,11}$ (Effondrement $89\%$).
                  \item Léopard : $0 / 0,20 = \mathbf{0}$ (Disparu des détections).
                  \item Céphalophe : $3,0 / 3,5 = \mathbf{0,86}$ (Maintien relatif).
                  \item Potamochère : $1,58 / 1,33 = \mathbf{1,19}$ (Légère augmentation).
              \end{itemize}
              \textbf{Sensibles (Ratio $< 0,5$) :} Éléphant, Bongo, Léopard (Grands mammifères, grande aire de vie, braconnage ciblé, faible taux reproduction).
              \textbf{Tolérantes / Opportunistes (Ratio $\approx 1$ ou $> 1$) :} Céphalophe (petit, cryptique, reproduction rapide), Potamochère (omnivore, s'adapte cultures/champs, groupe social).
        \item \textbf{Absence de détection $\neq$ Extinction.}
              \textbf{Biais « Camera Traps » :}
              \begin{itemize}
                  \item Effort d'échantillonnage fini (6000 j-pièges) : Espèce rare + grande aire de vie = probabilité détection faible.
                  \item Léopard : Solitaire, territorial, densité naturellement très faible ($< 1/100\,\mathrm{km^2}$). Non-détection probable = « Faux négatif ».
                  \item Nécessite : Augmenter effort, utiliser indices (traces, crottes, griffades), enquêtes locales, pièges à poils (ADN).
              \end{itemize}
        \item \textbf{Hypothèse Potamochère (Site Z $>$ Site P) :}
              \begin{itemize}
                  \item \textbf{Effet « Bordure / Ressource » :} Site Z = Mosaïque forêt/champs/jachères. Potamochère affectionne les écotones et pille les cultures (manioc, maïs, patates) $\rightarrow$ Ressource alimentaire abondante, prévisible, riche en énergie.
                  \item \textbf{Relâchement prédation :} Disparition Léopard (prédateur principal potamochères adultes) au Site Z $\rightarrow$ Mortalité naturelle diminuée.
                  \item \textbf{Comportement :} Groupes sonneurs bruyants, peu farouches $\rightarrow$ Détectabilité plus élevée près des pièges placés sur sentiers/champs.
              \end{itemize}
    \end{enumerate}
    \textbf{Conclusion :} La protection stricte (Site P) maintient l'intégrité de la communauté (grands mammifères sensibles). L'exploitation cynégétique/agricole (Site Z) filtre la communauté : perte des espèces « parapluie » (Éléphant, Léopard, Bongo), maintien/augmentation d'espèces communes/adaptées (Céphalophe, Potamochère). Le corridor et le renforcement de la surveillance dans la ZICGC sont nécessaires pour les espèces sensibles.
\end{exoresolu}

\begin{attention}[Les 5 erreurs qui coûtent des points au BEPC]
    \begin{enumerate}
        \item \textbf{Confondre « Facteur abiotique » et « Être vivant » / « Relation trophique ».}
              \textit{Exemple :} « Le soleil mange l'herbe » ou « La pluie est un consommateur ». \textbf{Correction :} Le soleil (énergie), l'eau, le sol, la température, la lumière sont des \textbf{facteurs abiotiques (milieu)}. Seuls les êtres vivants (producteurs, consommateurs, décomposeurs) entrent dans les chaînes trophiques. La flèche va \textbf{de la proie vers le prédateur} (sens du transfert matière/énergie).
        \item \textbf{Oublier les Décomposeurs dans le cycle de la matière / Dire que l'énergie « recycle ».}
              \textit{Erreur :} « L'énergie revient au sol par les décomposeurs pour les plantes. » \textbf{Correction :} L'énergie \textbf{ne recycle pas}, elle se dissipe en \textbf{chaleur} à chaque niveau (2\up{nd} principe thermodynamique). Seule la \textbf{MATIÈRE} (C, N, P, K, eau) est recyclée par les décomposeurs (minéralisation $\rightarrow$ sels minéraux $\rightarrow$ producteurs). Toujours distinguer \textbf{Flux d'énergie (linéaire, entrant/sortant)} et \textbf{Cycle de la matière (circulaire)}.
        \item \textbf{Généraliser « Feu = Mauvais » ou « Homme = Destructeur » sans nuance scientifique.}
              \textit{Erreur :} « Il faut interdire tout feu de brousse. » \textbf{Correction :} Le feu est un \textbf{outil de gestion} en savane (régénération pâturages, lutte tiques, stimulation germination certaines espèces pyrophiles). Le problème est le \textbf{feu tardif, trop fréquent, non contrôlé} (destruction graines, jeunes plants, matière organique sol, faune). De même, l'homme peut être \textbf{gestionnaire durable} (agroforesterie, ZICGC, corridors). Nuancer : \textit{Pression} vs \textit{Gestion durable}.
        \item \textbf{Calculer un pourcentage d'évolution ou un IKA sans respecter la formule ou les unités.}
              \textit{Erreurs classiques :} $\frac{V_f - V_i}{V_f}$ au lieu de $V_i$ ; oublier $\times 100$ pour le \%; inverser numérateur/dénominateur IKA (km/contacts au lieu de contacts/km) ; ne pas mettre l'unité finale (\% ou ind/km ou contacts/100 j-pièges). \textbf{Réflexe :} Écrire la formule, identifier $V_i$ et $V_f$, calculer, écrire phrase réponse avec unité.
        \item \textbf{Répondre « Pour protéger les animaux » sans argumenter par les Services Écosystémiques ni l'Homme.}
              \textit{Réponse insuffisante (0,5 pt / 3 pts) :} « Il faut créer un parc pour sauver les éléphants. » \textbf{Attendu (Note max) :} Argumenter en 3 piliers : (1) \textbf{Écologique :} Espèce clé / parapluie, connectivité, cycle eau/carbone. (2) \textbf{Économique :} Écotourisme (devises, emplois), ressources (médicaments, bois, miel), prévention coûts dégâts (érosion, conflits). (3) \textbf{Social / Éthique :} Droits populations autochtones, patrimoine culturel, équité intergénérationnelle. \textbf{Mots-clés obligatoires :} \cle{Services écosystémiques}, \cle{Développement durable}, \cle{Participation communautaire}.
    \end{enumerate}
\end{attention}

\newpage

\coursec{Exercices}

\courssub{Je vérifie mes connaissances}

\begin{exercice}[QCM : Caractéristiques des écosystèmes camerounais]
Parmi les affirmations suivantes, une seule est correcte. La recopier et la justifier brièvement.
\begin{enumerate}
    \item La forêt équatoriale camerounaise se caractérise par une faible pluviométrie (inférieure à \SI{1000}{mm/an}) et une strate herbacée dense.
    \item La savane arbustive du Nord-Cameroun présente une biodiversité spécifique plus élevée que la forêt dense humide du Sud.
    \item Dans un écosystème, les décomposeurs (champignons, bactéries) assurent la \cle{minéralisation} de la matière organique, restituant les \cle{éléments minéraux} aux producteurs (végétaux chlorophylliens).
    \item Le braconnage n'affecte que les populations des espèces chassées sans conséquence sur la structure de l'écosystème.
\end{enumerate}
\end{exercice}

\begin{exercice}[Vrai ou Faux : Interdépendances et menaces]
Dire si les affirmations sont vraies (V) ou fausses (F). Justifier chaque réponse par une phrase précise.
\begin{enumerate}
    \item Les réseaux trophiques sont plus complexes et plus stables dans la savane que dans la forêt équatoriale.
    \item La déforestation pour l'agriculture sur brûlis entraîne une minéralisation rapide de l'humus par le feu, suivie d'une érosion des sols et d'une perte durable de fertilité.
    \item La disparition des éléphants de forêt n'a pas d'impact sur la dissémination des graines des grands arbres forestiers.
    \item Le reboisement avec une seule essence d'arbre (monoculture) restaure pleinement la biodiversité originelle d'une forêt détruite.
\end{enumerate}
\end{exercice}

\begin{exercice}[Définitions et vocabulaire]
Définir rigoureusement les termes suivants en utilisant le vocabulaire scientifique du programme :
\begin{enumerate}
    \item \cle{Écosystème}
    \item \cle{Biodiversité} (ou diversité biologique)
    \item \cle{Niche écologique}
    \item \cle{Pyrophyte} (donner un exemple d'espèce camerounaise)
    \item \cle{Corridor biologique}
\end{enumerate}
\end{exercice}

\begin{exercice}[Calcul direct : Taux de déforestation]
D'après les données du Global Forest Watch, la superficie forestière du Cameroun était d'environ \SI{22,5}{Mha} (\SI{22,5}{\mega\hectare}) en 2000 et d'environ \SI{20,1}{Mha} en 2020.
\begin{enumerate}
    \item Calculer la surface forestière perdue sur cette période (en \si{\hectare}).
    \item Calculer le taux de perte annuel moyen (en \si{\hectare/an}).
    \item Exprimer cette perte en pourcentage de la surface initiale de 2000.
\end{enumerate}
\end{exercice}

\courssub{Je m'entraîne}

\begin{exercice}[Comparaison de la biodiversité : Forêt vs Savane]
Le tableau ci-dessous présente des relevés simplifiés d'espèces végétales et animales observées sur deux parcelles de \SI{1}{\hectare} chacune, l'une dans la forêt dense du Parc National de Campo Ma'an (Sud), l'autre dans la savane du Parc National de Waza (Extrême-Nord).

\begin{center}
\begin{tabularx}{\linewidth}{|l|X|X|}\hline
\rowcolor{popblueL}
\textbf{Groupe taxonomique} & \textbf{Forêt dense (Campo Ma'an)} & \textbf{Savane (Waza)} \\ \hline
Arbres (espèces) & 120 & 15 \\
Arbustes/Lianes (espèces) & 85 & 25 \\
Herbacées (espèces) & 40 & 60 \\
Mammifères (espèces) & 45 & 30 \\
Oiseaux (espèces) & 280 & 180 \\
Insectes (espèces estimées) & 1500 & 400 \\ \hline
\end{tabularx}
\end{center}

\begin{enumerate}
    \item Calculer la richesse spécifique totale (somme des espèces) pour chaque milieu.
    \item Comparer la biodiversité des deux écosystèmes. Lequel est le plus riche ? Justifier par les chiffres.
    \item Expliquer deux facteurs abiotiques principaux qui justifient cette différence de biodiversité au Cameroun.
    \item Pourquoi la strate herbacée est-elle plus diversifiée en savane qu'en forêt dense ?
\end{enumerate}
\end{exercice}

\begin{exercice}[Construction et analyse d'un réseau trophique savanic]
On considère les espèces suivantes de la savane de l'Adamaoua et leurs régimes alimentaires :
\begin{itemize}
    \item Herbes (\emph{Andropogon}, \emph{Hyparrhenia}) : Producteurs.
    \item Sauterelles : Phytophages (mangent les herbes).
    \item Rats des champs : Omnivores (graines, herbes, insectes).
    \item Lézards : Insectivores (mangent sauterelles).
    \item Serpents : Carnivores (mangent rats, lézards).
    \item Rapaces (Faucons) : Carnivores (mangent rats, serpents, lézards).
    \item Chacals : Charognards / Carnivores opportunistes (mangent tout cadavre, petits mammifères).
    \item Bactéries du sol / Champignons : Décomposeurs.
\end{itemize}

\begin{enumerate}
    \item Construire le réseau trophique correspondant (flèches \emph{proie $\rightarrow$ prédateur}).
    \item Identifier une chaîne trophique à 4 maillons et une à 5 maillons.
    \item Prédire et expliquer les conséquences sur les populations de rats, de serpents et d'herbes si les rapaces disparaissent totalement (braconnage, empoisonnement).
    \item Quel rôle jouent les décomposeurs dans la pérennité de ce réseau ? Que se passerait-il s'ils disparaissaient ?
\end{enumerate}
\end{exercice}

\begin{exercice}[Analyse de l'impact d'un feu de brousse sur le sol]
Un feu de brousse tardif (fin de saison sèche, mars) ravage une parcelle de savane boisée destinée à la culture du maïs l'année suivante. Le feu est intense (vent fort, herbes sèches hautes).

\begin{enumerate}
    \item Décrire l'effet immédiat du feu sur la matière organique superficielle (litière, humus) et sur la faune du sol (vers de terre, termites, micro-organismes).
    \item Expliquer le mécanisme par lequel le feu entraîne, à moyen terme (prochaine saison des pluies), une érosion hydrique accrue.
    \item L'agriculteur observe une baisse de rendement du maïs de \num{40}\% par rapport aux années sans feu. Relier cette baisse à l'appauvrissement du sol (azote, phosphore, structure).
    \item Proposer deux alternatives culturales au brûlis pour préparer le champ (ex : agriculture de conservation, agroforesterie).
\end{enumerate}
\end{exercice}

\begin{exercice}[Étude de cas : Le braconnage de l'éléphant de forêt et la régénération]
Dans le Parc National de Lobéké (Sud-Est), les inventaires montrent une baisse de \num{60}\% de la population d'éléphants de forêt (\emph{Loxodonta cyclotis}) en 10 ans due au braconnage pour l'ivoire. Parallèlement, on observe une diminution de la densité de jeunes plants d'espèces d'arbres à gros fruits (ex : \emph{Irvingia gabonensis} - andok, \emph{Detarium macrocarpum}).

\begin{enumerate}
    \item Rappeler le mode de dissémination des graines de ces arbres (zoochorie). Pourquoi l'éléphant est-il un disséminateur « irremplaçable » pour ces grosses graines ?
    \item Expliquer la chaîne de causalité : Braconnage $\rightarrow$ Baisse éléphants $\rightarrow$ Baisse dissémination $\rightarrow$ Baisse régénération $\rightarrow$ Modification structure forêt.
    \item Citer deux autres conséquences écologiques de la disparition des éléphants (création de clairières, points d'eau, dispersion nutriments).
    \item Proposer une mesure de conservation efficace et une mesure de sensibilisation des populations riveraines.
\end{enumerate}
\end{exercice}

\courssub{Je me prépare au BEPC}

\begin{exercice}[Sujet type BEPC : Évolution de la couverture forestière et stratégies de conservation]
\textbf{Document 1 :} Évolution de la superficie forestière au Cameroun (Source : FAO / MINEPDED).
\begin{center}
\begin{tabular}{|c|c|c|c|c|}\hline
\rowcolor{popblueL}
\textbf{Année} & 1990 & 2000 & 2010 & 2020 \\ \hline
\textbf{Surface forestière (millions ha)} & 22,7 & 22,5 & 21,0 & 20,1 \\ \hline
\end{tabular}
\end{center}

\textbf{Document 2 :} Principales causes directes de déforestation identifiées au Cameroun (ordre d'importance) :
\begin{enumerate}
    \item Agriculture itinérante sur brûlis (vivrier : manioc, maïs, plantain).
    \item Exploitation forestière illégale ou non durable (bois d'œuvre).
    \item Expansion des agro-industries (palmier à huile, hévéa, cacao).
    \item Prélèvement de bois de feu et charbon de bois (énergie domestique).
    \item Feux de brousse non maîtrisés.
\end{enumerate}

\textbf{Questions :}
\begin{enumerate}
    \item \textbf{Analyse quantitative :} À l'aide du Document 1, calculer la surface forestière totale perdue entre 1990 et 2020. Calculer le taux de déforestation annuel moyen sur cette période (en \si{ha/an} et en \%/an par rapport à 1990).
    \item \textbf{Interprétation :} En vous appuyant sur le Document 2, expliquer pourquoi l'agriculture sur brûlis est la première cause de déforestation au Cameroun. Relier à la démographie et aux pratiques culturales.
    \item \textbf{Conséquences écologiques :} Citer et expliquer trois conséquences majeures de cette perte de forêt sur : (a) la biodiversité, (b) le cycle de l'eau / climat local, (c) les sols.
    \item \textbf{Mesures de conservation :} Proposer \textbf{trois} mesures concrètes et applicables au contexte camerounais pour freiner cette déforestation. Pour chaque mesure, préciser l'acteur principal (État, ONG, populations locales, secteur privé) et le principe écologique mobilisé (ex : aire protégée, agroforesterie, reboisement, énergie alternative).
    \item \textbf{Synthèse :} Rédiger un paragraphe argumenté (10-12 lignes) expliquant pourquoi la conservation des forêts du Bassin du Congo (dont le Cameroun) est un enjeu planétaire (climat, biodiversité, populations autochtones).
\end{enumerate}
\end{exercice}

\begin{exercice}[Sujet type BEPC : Réhabilitation d'un écosystème dégradé par le reboisement]
Une école secondaire de la région de l'Ouest (climat tropical d'altitude, deux saisons des pluies) souhaite réhabiliter une colline dégradée (\SI{2}{\hectare}), sujette à l'érosion, en y créant un bois scolaire pédagogique. Le sol est argilo-ferrugineux, peu profond, acide (pH \num{4,5}), pauvre en matière organique. L'établissement dispose d'une pépinière scolaire.

\textbf{Questions :}
\begin{enumerate}
    \item \textbf{Choix des essences :} Proposer \textbf{trois} essences d'arbres (noms scientifiques et vernaculaires) adaptées à ce milieu (altitude, sol acide, pente) et utiles pour l'école (ombrage, bois d'œuvre, fruitiers, médicinales, fixation d'azote). Justifier chaque choix par une adaptation écologique ou un usage.
    \item \textbf{Calendrier et protocole :} Établir le calendrier détaillé des opérations de reboisement sur une année (mois par mois) : préparation du terrain (désherbage, traçage, trouage), production/achat plants, plantation, entretien (désherbage, tuteurage, pare-feu), surveillance. Justifier la période de plantation choisie par rapport au cycle de l'eau.
    \item \textbf{Technique de plantation :} Décrire la technique du « trou de plantation » (dimensions, amendement, positionnement du plant, arrosage, buttage). Pourquoi ne pas enterrer le collet ?
    \item \textbf{Évaluation du succès :} Définir deux indicateurs de réussite du reboisement à 6 mois et à 2 ans. Proposer une action corrective si le taux de reprise est inférieur à \num{70}\%.
    \item \textbf{Valeur pédagogique :} Expliquer comment ce bois scolaire peut servir de support aux enseignements de SVT (écologie, botanique, sol, climat) et d'éducation à la citoyenneté environnementale.
\end{enumerate}
\end{exercice}

\begin{exercice}[Sujet type BEPC : Gestion durable d'une aire protégée face aux pressions anthropiques]
Le Parc National de la Benoué (Nord, savane soudanienne, \SI{180 000}{\hectare}) fait face à des pressions croissantes : transhumance massive de troupeaux (zébus) en saison sèche, braconnage commercial (viande de brousse), orpaillage illégal (mercure dans la rivière), et agriculture aux abords.

\textbf{Document :} Extrait du plan de gestion 2020-2024.
\begin{itemize}
    \item Zone cœur (intégrale) : \SI{60 000}{\hectare} (interdiction totale activité humaine).
    \item Zone tampon (multiple usage) : \SI{80 000}{\hectare} (chasse villageoise encadrée, cueillette, pâturage contrôlé).
    \item Zone de transition (périphérie) : \SI{40 000}{\hectare} (agriculture, habitat, corridors).
\end{itemize}

\textbf{Questions :}
\begin{enumerate}
    \item \textbf{Zonage et objectifs :} Expliquer la logique écologique du zonage en trois zones (cœur, tampon, transition). Pourquoi la zone cœur doit-elle être la plus grande possible pour les grands mammifères (lion, girafe, damalisque) ?
    \item \textbf{Conflit d'usage :} Analyser le conflit entre les éleveurs transhumants (besoin de pâture/eau en saison sèche) et la conservation (surpâturage, piétinement, compétition faune sauvage / bétail, maladies). Proposer \textbf{deux} solutions concertées (ex : couloirs de transhumance, points d'eau hors parc, vaccination).
    \item \textbf{Pollution par le mercure :} L'orpaillage utilise le mercure pour l'amalgamation de l'or. Expliquer le phénomène de \textbf{bioamplification} (bioaccumulation le long de la chaîne trophique) du mercure (méthylmercure) depuis l'eau/sédiments $\rightarrow$ poissons $\rightarrow$ oiseaux piscivores / homme. Quel risque sanitaire pour les populations riveraines ?
    \item \textbf{Corridors écologiques :} Le parc est isolé par l'agriculture. Définir un corridor écologique. Pourquoi est-il vital de connecter la Benoué au Parc de Bouba Ndjidda (au Nord) et à la zone de chasse de la Bénoué (au Sud) ? Citer deux espèces « parapluie » dont la survie dépend de ces corridors.
    \item \textbf{Évaluation de la gestion :} On observe une augmentation de \num{15}\% des indices de présence des grands mammifères en 4 ans dans la zone cœur, mais une baisse de \num{30}\% dans la zone tampon. Interpréter ces résultats contrastés et proposer un ajustement de la gestion de la zone tampon.
\end{enumerate}
\end{exercice}

\newpage

\coursec{Corrigés des exercices}

\begin{corrige}[Exercice 1 — QCM : Caractéristiques des écosystèmes camerounais]
\textbf{Réponse correcte : 3.}

\textbf{Justification de la réponse 3 :} Les \cle{décomposeurs} (champignons, bactéries) dégradent la matière organique morte (litière, cadavres, déjections) et la transforment en \cle{éléments minéraux} (nitrates, phosphates...) assimilables par les végétaux chlorophylliens : c'est la \cle{minéralisation}, qui boucle le cycle de la matière dans l'écosystème.

\textbf{Pourquoi les autres sont fausses :}
\begin{enumerate}
    \item \textbf{Fausse :} la forêt équatoriale reçoit une pluviométrie élevée (supérieure à \SI{1500}{mm/an}, souvent \SI{2000}{mm/an} au Sud) ; la strate herbacée y est au contraire peu développée car la lumière n'atteint pas le sol sous la canopée dense.
    \item \textbf{Fausse :} c'est la forêt dense humide du Sud qui présente la biodiversité spécifique la plus élevée (climat chaud et humide toute l'année, grande variété de milieux) ; la savane du Nord, soumise à une longue saison sèche, est moins riche en espèces.
    \item \textbf{Fausse :} le braconnage a des conséquences en cascade sur tout l'écosystème : disparition des disséminateurs de graines, déséquilibre des réseaux trophiques (prolifération des proies, régression des prédateurs).
\end{enumerate}
\end{corrige}

\begin{corrige}[Exercice 2 — Vrai ou Faux : Interdépendances et menaces]
\begin{enumerate}
    \item \textbf{Faux.} Les réseaux trophiques sont plus complexes et plus stables dans la \textbf{forêt équatoriale}, car la richesse spécifique y est beaucoup plus grande : plus il y a d'espèces et de liens alimentaires, plus la disparition d'une espèce peut être compensée, donc plus l'écosystème est stable.
    \item \textbf{Vrai.} Le feu détruit la litière et l'humus (minéralisation brutale de la matière organique) ; le sol, nu et sans structure, est battu par les pluies puis emporté par l'érosion hydrique, d'où une perte durable de fertilité.
    \item \textbf{Faux.} L'éléphant de forêt est un disséminateur majeur (\cle{zoochorie} endozoochore) : il avale les gros fruits et rejette les graines dans ses crottes, parfois à plusieurs kilomètres. Sa disparition réduit la régénération des grands arbres à gros fruits.
    \item \textbf{Faux.} Une monoculture ne restaure qu'une faible partie de la biodiversité : la forêt originelle comporte des centaines d'espèces végétales et animales associées. Un reboisement restauratif doit utiliser un \textbf{mélange d'essences locales diversifiées}.
\end{enumerate}
\end{corrige}

\begin{corrige}[Exercice 3 — Définitions et vocabulaire]
\begin{enumerate}
    \item \textbf{Écosystème :} ensemble formé par une communauté d'êtres vivants (\cle{biocénose}) et son milieu de vie (\cle{biotope}, facteurs abiotiques : climat, sol, eau, lumière), ainsi que les interactions qui s'établissent entre eux.
    \item \textbf{Biodiversité :} diversité des êtres vivants à tous les niveaux : diversité des espèces (richesse spécifique), diversité génétique au sein de chaque espèce et diversité des écosystèmes.
    \item \textbf{Niche écologique :} rôle fonctionnel d'une espèce dans son écosystème : ce qu'elle consomme, la place qu'elle occupe dans le réseau trophique, ses conditions de vie et ses relations avec les autres espèces.
    \item \textbf{Pyrophyte :} plante adaptée aux feux de brousse, capable de résister au feu ou de repartir après son passage (écorce épaisse, bourgeons souterrains, repousse à partir de la souche). Exemple camerounais : le karité (\emph{Vitellaria paradoxa}) ou le néré (\emph{Parkia biglobosa}) des savanes du Nord.
    \item \textbf{Corridor biologique :} bande de milieu naturel reliant deux habitats fragmentés (ex : deux forêts ou deux parcs), permettant le déplacement, la reproduction et les échanges génétiques entre populations d'espèces sauvages.
\end{enumerate}
\end{corrige}

\begin{corrige}[Exercice 4 — Calcul direct : Taux de déforestation]
\begin{enumerate}
    \item \textbf{Surface perdue :}
    \[ \SI{22,5}{Mha} - \SI{20,1}{Mha} = \SI{2,4}{Mha} = \SI{2400000}{\hectare} \]
    \item \textbf{Taux annuel moyen} (période 2000--2020, soit 20 ans) :
    \[ \frac{\SI{2400000}{\hectare}}{\SI{20}{ans}} = \SI{120000}{\hectare/an} \]
    \item \textbf{Pourcentage de perte par rapport à 2000 :}
    \[ \frac{2,4}{22,5} \times 100 \approx 10,7\,\% \]
\end{enumerate}
\textbf{Conclusion :} en 20 ans, le Cameroun a perdu environ \SI{2,4}{Mha} de forêt, soit près de \num{11}\% de sa couverture forestière de 2000, à un rythme moyen de \SI{120000}{\hectare/an}.

\textbf{Point de vigilance :} bien convertir les millions d'hectares en hectares ($\SI{1}{Mha} = \SI{1000000}{\hectare}$) et indiquer les unités à chaque étape.
\end{corrige}

\begin{corrige}[Exercice 5 — Comparaison de la biodiversité : Forêt vs Savane]
\begin{enumerate}
    \item \textbf{Richesse spécifique totale :}
    \begin{itemize}
        \item Forêt (Campo Ma'an) : $120 + 85 + 40 + 45 + 280 + 1500 = \textbf{2070 espèces}$.
        \item Savane (Waza) : $15 + 25 + 60 + 30 + 180 + 400 = \textbf{710 espèces}$.
    \end{itemize}
    \item \textbf{Comparaison :} la forêt dense est nettement plus riche : \num{2070} espèces contre \num{710}, soit environ \textbf{2,9 fois plus} d'espèces ($2070/710 \approx 2,9$). L'écart est considérable pour les arbres ($120$ contre $15$) et les insectes ($1500$ contre $400$).
    \item \textbf{Facteurs abiotiques :}
    \begin{itemize}
        \item \textbf{La pluviométrie :} la forêt du Sud reçoit plus de \SI{1500}{mm/an} répartis presque toute l'année, alors que la savane de l'Extrême-Nord subit une longue saison sèche (7 à 8 mois) avec moins de \SI{1000}{mm/an}.
        \item \textbf{La température et l'humidité permanentes :} le climat équatorial chaud et humide en continu permet une croissance végétale ininterrompue et une grande productivité, base d'une chaîne alimentaire riche.
    \end{itemize}
    \item \textbf{Strate herbacée :} en forêt dense, la canopée fermée intercepte la lumière : le sous-bois est sombre, ce qui limite les herbacées. En savane, les arbres sont clairsemés, la lumière atteint le sol : les herbacées (graminées) y sont abondantes et diversifiées.
\end{enumerate}
\end{corrige}

\begin{corrige}[Exercice 6 — Construction et analyse d'un réseau trophique savanicole]
\begin{enumerate}
    \item \textbf{Réseau trophique} (flèche = « est mangé par ») :

\begin{popfigure}
\begin{center}
\begin{tikzpicture}[>=Stealth, node distance=1.1cm, every node/.style={font=\small}]
\node[draw, fill=popgreenL, rounded corners] (herbes) {Herbes};
\node[draw, fill=popgoldL, rounded corners, above left=1cm and 1.6cm of herbes] (saut) {Sauterelles};
\node[draw, fill=popgoldL, rounded corners, above right=1cm and 1.6cm of herbes] (rats) {Rats};
\node[draw, fill=poporangeL, rounded corners, above=1.1cm of saut] (lez) {Lézards};
\node[draw, fill=poporangeL, rounded corners, above=2.2cm of herbes] (serp) {Serpents};
\node[draw, fill=poppinkL, rounded corners, above=1.1cm of rats] (rap) {Rapaces};
\node[draw, fill=poppinkL, rounded corners, right=1.4cm of rap] (chac) {Chacals};
\node[draw, fill=popblueL, rounded corners, below=1cm of herbes] (dec) {Décomposeurs};
\draw[->, thick, popdark] (herbes) -- (saut);
\draw[->, thick, popdark] (herbes) -- (rats);
\draw[->, thick, popdark] (saut) -- (lez);
\draw[->, thick, popdark] (saut) -- (rats);
\draw[->, thick, popdark] (lez) -- (serp);
\draw[->, thick, popdark] (lez) -- (rap);
\draw[->, thick, popdark] (rats) -- (serp);
\draw[->, thick, popdark] (rats) -- (rap);
\draw[->, thick, popdark] (rats) -- (chac);
\draw[->, thick, popdark] (serp) -- (rap);
\draw[->, thick, popdark, dashed] (herbes) -- (dec);
\draw[->, thick, popdark, dashed] (serp) to[bend left=15] (dec);
\draw[->, thick, popdark, dashed] (chac) to[bend right=20] (dec);
\draw[->, thick, popgreen, dashed] (dec) -- (herbes);
\end{tikzpicture}
\end{center}
\legende{Réseau trophique de la savane de l'Adamaoua. Flèches pleines : relations proie $\rightarrow$ prédateur ; flèches pointillées : matière organique morte vers les décomposeurs, qui restituent les éléments minéraux aux producteurs.}
\end{popfigure}

    \item \textbf{Chaînes trophiques :}
    \begin{itemize}
        \item 4 maillons : Herbes $\rightarrow$ Sauterelles $\rightarrow$ Lézards $\rightarrow$ Serpents (ou $\rightarrow$ Rapaces).
        \item 5 maillons : Herbes $\rightarrow$ Sauterelles $\rightarrow$ Rats $\rightarrow$ Serpents $\rightarrow$ Rapaces (ou Herbes $\rightarrow$ Sauterelles $\rightarrow$ Lézards $\rightarrow$ Serpents $\rightarrow$ Rapaces).
    \end{itemize}
    \item \textbf{Disparition des rapaces :}
    \begin{itemize}
        \item Les \textbf{serpents}, libérés d'un prédateur, \textbf{prolifèrent} dans un premier temps.
        \item Les \textbf{rats}, plus chassés par les serpents devenus nombreux, \textbf{diminuent}.
        \item Les \textbf{herbes}, moins broutées par les rats (et sauterelles réduites par les lézards), peuvent \textbf{augmenter} temporairement.
    \end{itemize}
    Cela montre qu'un réseau trophique est un système en équilibre : la suppression d'un maillon perturbe tous les autres.
    \item \textbf{Rôle des décomposeurs :} ils recyclent la matière organique morte (cadavres, déjections, végétaux secs) en \cle{éléments minéraux} réutilisables par les herbes : ils assurent le \textbf{cycle de la matière} et la fertilité du sol. Sans eux, la matière organique s'accumulerait, les éléments minéraux s'épuiseraient, les producteurs dépériraient et tout le réseau s'effondrerait.
\end{enumerate}
\end{corrige}

\begin{corrige}[Exercice 7 — Analyse de l'impact d'un feu de brousse sur le sol]
\begin{enumerate}
    \item \textbf{Effet immédiat :} le feu intense brûle la litière et l'humus : la matière organique est \textbf{minéralisée brutalement} (transformée en cendres). La faune du sol est détruite ou chassée : vers de terre et micro-organismes des premiers centimètres sont tués par la chaleur ; les termites, plus profonds, sont partiellement épargnés mais leur activité diminue.
    \item \textbf{Érosion hydrique à la saison des pluies :} le sol, privé de couvert végétal et d'humus, est nu et sa structure est dégradée (plus de galeries de vers ni de racines vivantes). Les gouttes de pluie battent directement le sol, l'imperméabilisent (croûte de battance) ; l'eau ne s'infiltre plus et ruisselle en emportant les particules fines : c'est l'\textbf{érosion hydrique}, qui enlève la couche fertile.
    \item \textbf{Baisse de rendement (\num{40}\%) :} l'azote de la matière organique part en fumée lors de la combustion ; le phosphore et le potassium des cendres sont lessivés par les pluies ; la structure dégradée limite l'enracinement et la rétention d'eau. Le maïs, gros consommateur d'azote, souffre de cette faim d'éléments minéraux : le rendement chute.
    \item \textbf{Alternatives au brûlis :}
    \begin{itemize}
        \item \textbf{Agriculture de conservation :} couper les herbes et les laisser au sol en \textbf{paillage} (mulch), qui protège le sol, garde l'humidité et se décompose en humus.
        \item \textbf{Agroforesterie :} associer au maïs des arbres fixateurs d'azote (ex : \emph{Calliandra}, \emph{Gliricidia}) dont la biomasse enrichit le sol et dont les racines limitent l'érosion.
    \end{itemize}
\end{enumerate}
\end{corrige}

\begin{corrige}[Exercice 8 — Étude de cas : Le braconnage de l'éléphant de forêt et la régénération]
\begin{enumerate}
    \item \textbf{Mode de dissémination :} ces arbres se reproduisent par \cle{zoochorie}, plus précisément \textbf{endozoochorie} : l'éléphant consomme les gros fruits, les graines traversent son tube digestif et sont rejetées intactes (voire mieux aptes à germer) dans les crottes, loin de l'arbre parent. L'éléphant est « irremplaçable » car il est le \textbf{seul animal assez grand} pour avaler ces très grosses graines et les transporter sur plusieurs kilomètres ; les autres frugivores ne peuvent pas les ingérer.
    \item \textbf{Chaîne de causalité :} le braconnage pour l'ivoire fait chuter la population d'éléphants (\num{-60}\%) ; moins d'éléphants signifie moins de graines disséminées loin des arbres parents ; les graines tombées sous l'arbre parent germent mal (compétition, prédateurs de graines concentrés) ; la régénération des essences à gros fruits diminue ; à terme, la composition et la structure de la forêt se modifient (vieillissement des peuplements, appauvrissement en essences à gros fruits).
    \item \textbf{Autres conséquences :}
    \begin{itemize}
        \item Disparition des \textbf{clairières} entretenues par les éléphants, pourtant utiles à de nombreuses espèces (bongos, buffles, gorilles qui y recherchent les sels minéraux).
        \item Moindre \textbf{dispersion des nutriments} : les crottes d'éléphants fertilisent le sol et nourrissent de nombreux insectes (bousiers).
    \end{itemize}
    \item \textbf{Mesures :}
    \begin{itemize}
        \item \textbf{Conservation :} renforcer les brigades de \textbf{lutte anti-braconnage} (écogardes armés, patrouilles, surveillance du parc) et réprimer le commerce de l'ivoire.
        \item \textbf{Sensibilisation :} éduquer les populations riveraines (écoles, radios locales, chefs traditionnels) sur le rôle de l'éléphant et développer des revenus alternatifs (écotourisme, apiculture) pour réduire la dépendance au braconnage.
    \end{itemize}
\end{enumerate}
\end{corrige}

\begin{corrige}[Exercice 9 — Sujet type BEPC : Évolution de la couverture forestière et stratégies de conservation]
\textbf{Barème indicatif (sur 20) :} Q1 : 4 pts ; Q2 : 3 pts ; Q3 : 5 pts ; Q4 : 4 pts ; Q5 : 4 pts.

\begin{enumerate}
    \item \textbf{Analyse quantitative :}
    \begin{itemize}
        \item Surface perdue entre 1990 et 2020 : $22,7 - 20,1 = \SI{2,6}{Mha} = \SI{2600000}{\hectare}$.
        \item Taux annuel moyen (30 ans) : $\dfrac{\SI{2600000}{\hectare}}{30} \approx \SI{86700}{\hectare/an}$.
        \item En pourcentage par rapport à 1990 : $\dfrac{2,6}{22,7} \times 100 \approx 11,5\,\%$ sur 30 ans, soit environ $\dfrac{11,5}{30} \approx 0,38\,\%$/an.
    \end{itemize}
    \item \textbf{Interprétation :} l'agriculture sur brûlis est la première cause car la population camerounaise croît rapidement et une grande partie des ruraux pratique une agriculture vivrière itinérante : on défriche et brûle une parcelle de forêt, on la cultive quelques années, puis on l'abandonne quand la fertilité s'épuise pour en défricher une nouvelle. La croissance démographique accélère la rotation et grignote la forêt.
    \item \textbf{Conséquences écologiques :}
    \begin{itemize}
        \item \textbf{(a) Biodiversité :} destruction des habitats $\rightarrow$ disparition des espèces (plantés, okapis, gorilles...), fragmentation des populations, appauvrissement génétique.
        \item \textbf{(b) Cycle de l'eau / climat :} moins d'évapotranspiration $\rightarrow$ diminution des pluies locales, assèchement, hausse des températures ; moins de stockage du carbone $\rightarrow$ aggravation du réchauffement climatique.
        \item \textbf{(c) Sols :} sols nus exposés à l'érosion hydrique et à la latérisation, perte d'humus et de fertilité, ensablement des cours d'eau.
    \end{itemize}
    \item \textbf{Mesures de conservation (exemples) :}
    \begin{itemize}
        \item \textbf{Création et gestion d'aires protégées} (parcs nationaux, réserves) — acteur : \textbf{État} (MINFOF) ; principe : sanctuarisation des habitats.
        \item \textbf{Agroforesterie cacaoyère} (cacao sous couvert forestier) — acteurs : \textbf{populations locales et ONG} ; principe : concilier production agricole et maintien du couvert arboré.
        \item \textbf{Reboisement des zones dégradées et promotion des foyers améliorés / gaz domestique} — acteurs : \textbf{État, secteur privé, populations} ; principe : restaurer les sols et réduire la pression sur le bois de feu.
    \end{itemize}
    \item \textbf{Synthèse (modèle de paragraphe) :} Les forêts du Bassin du Congo, deuxième massif forestier tropical du monde après l'Amazonie, constituent un enjeu planétaire. Elles stockent d'immenses quantités de carbone dans leur biomasse et leurs sols : leur destruction libère du dioxyde de carbone et aggrave le réchauffement climatique qui menace toute l'humanité. Elles abritent une biodiversité exceptionnelle (gorilles, chimpanzés, éléphants de forêt, milliers d'espèces végétales), patrimoine irremplaçable dont la disparition serait définitive. Elles régulent le cycle de l'eau à l'échelle du continent africain en alimentant les pluies bien au-delà de leurs limites. Enfin, elles sont le milieu de vie de populations autochtones (Baka, Bagyeli...) dont les savoirs traditionnels et les droits doivent être respectés. Protéger ces forêts, c'est donc agir à la fois pour le climat, pour la vie et pour la dignité humaine : c'est une responsabilité partagée entre le Cameroun et la communauté internationale.
\end{enumerate}

\textbf{Points de vigilance :} exiger les unités dans les calculs ; en Q3, ne pas se contenter de citer : \emph{expliquer} le mécanisme ; en Q4, chaque mesure doit être accompagnée de l'acteur et du principe écologique ; la synthèse doit mobiliser les trois axes demandés (climat, biodiversité, populations).
\end{corrige}

\begin{corrige}[Exercice 10 — Sujet type BEPC : Réhabilitation d'un écosystème dégradé par le reboisement]
\textbf{Barème indicatif (sur 20) :} Q1 : 4 pts ; Q2 : 5 pts ; Q3 : 4 pts ; Q4 : 3 pts ; Q5 : 4 pts.

\begin{enumerate}
    \item \textbf{Choix des essences (exemples justifiés) :}
    \begin{itemize}
        \item \textbf{\emph{Acacia mearnsii} (acacia noir) :} légumineuse \textbf{fixatrice d'azote}, tolère les sols acides et pauvres ; enrichit le sol, fournit bois de feu et tanin.
        \item \textbf{\emph{Prunus africana} (pygeum, « bome ») :} essence d'\textbf{altitude} de l'Ouest, adaptée au climat tropical montagnard ; arbre médicinal précieux (écorce utilisée en pharmacopée).
        \item \textbf{\emph{Persea americana} (avocatier, « pear ») :} \textbf{fruitier} utile à la cantine scolaire, s'adapte aux sols argileux d'altitude, fournit ombrage et revenus.
    \end{itemize}
    (Autres choix acceptables : \emph{Grevillea robusta} pour le bois d'œuvre et la haie brise-vent, \emph{Calliandra calothyrsus} comme fixateur d'azote fourrager.)
    \item \textbf{Calendrier (région de l'Ouest, deux saisons des pluies : grande saison mars--juin, petite saison septembre--novembre) :}
    \begin{itemize}
        \item \textbf{Janvier--février :} désherbage, traçage des lignes de plantation (courbes de niveau sur la pente pour limiter l'érosion), ouverture des trous (\SI{40}{cm} $\times$ \SI{40}{cm} $\times$ \SI{40}{cm}), confection des pare-feu.
        \item \textbf{Février--mars :} production des plants en pépinière scolaire (semis en sachets) ou achat ; apport de compost dans les trous.
        \item \textbf{Mars--avril :} \textbf{plantation au début de la grande saison des pluies}, pour que les jeunes plants disposent d'une humidité régulière assurant la reprise des racines avant la saison sèche.
        \item \textbf{Avril--juin :} arrosages d'appoint si besoin, désherbage autour des plants, tuteurage.
        \item \textbf{Juillet--août (petite saison sèche) :} surveillance, paillage au pied des plants pour conserver l'humidité.
        \item \textbf{Septembre--novembre :} deuxième désherbage, regarnissage (remplacement des plants morts) pendant la petite saison des pluies.
        \item \textbf{Novembre--décembre :} entretien des pare-feu avant la saison sèche, bilan du taux de reprise.
    \end{itemize}
    \textbf{Justification :} planter en début de saison des pluies synchronise la reprise racinaire avec la disponibilité en eau (cycle de l'eau) ; planter en saison sèche condamnerait les plants sans irrigation coûteuse.
    \item \textbf{Technique du trou de plantation :} trou de \SI{40}{cm} de côté et de profondeur ; séparer la terre de surface (riche) de la terre profonde ; mélanger la terre de surface avec du compost ; déchirer le sachet sans casser la motte ; placer le plant \textbf{à la verticale, le collet (jonction tige--racines) au niveau du sol} ; reboucher avec la terre enrichie, tasser légèrement, arroser, butter légèrement. \textbf{Ne pas enterrer le collet} car il pourrit au contact de l'humidité du sol (fonte des semis) et le plant meurt ; ne pas le laisser trop haut non plus, sinon les racines s'assèchent.
    \item \textbf{Indicateurs de réussite :}
    \begin{itemize}
        \item À 6 mois : \textbf{taux de reprise} (pourcentage de plants vivants) $\geq$ \num{80}\% ; émission de feuilles nouvelles.
        \item À 2 ans : \textbf{hauteur moyenne} des plants (ex : $\geq$ \SI{1,5}{m}) et couverture du sol limitant l'érosion visible (disparition des rigoles de ruissellement).
        \item \textbf{Action corrective} si reprise $<$ \num{70}\% : regarnissage massif à la saison des pluies suivante avec des plants plus âgés (plus résistants), paillage systématique et recherche de la cause de mortalité (sécheresse, feu, divagation animale $\rightarrow$ clôture).
    \end{itemize}
    \item \textbf{Valeur pédagogique :} le bois scolaire est un laboratoire vivant : observation directe des êtres vivants et de leurs relations (écologie), identification des essences (botanique), étude du sol et de l'érosion, mesures de pluie et de température (climat). Il éduque à la citoyenneté environnementale : les élèves deviennent acteurs de la protection de la nature et relaient les bonnes pratiques (anti-feux de brousse, reboisement) dans leurs familles.
\end{enumerate}

\textbf{Points de vigilance :} donner les noms scientifiques soulignés et un usage par essence ; le calendrier doit être cohérent avec le régime des pluies de l'Ouest ; la question du collet est un piège classique.
\end{corrige}

\begin{corrige}[Exercice 11 — Sujet type BEPC : Gestion durable d'une aire protégée face aux pressions anthropiques]
\textbf{Barème indicatif (sur 20) :} Q1 : 4 pts ; Q2 : 4 pts ; Q3 : 4 pts ; Q4 : 4 pts ; Q5 : 4 pts.

\begin{enumerate}
    \item \textbf{Logique du zonage :} la \textbf{zone cœur} est un sanctuaire intégral où les écosystèmes fonctionnent sans perturbation humaine : elle garantit la reproduction et la survie des espèces sensibles. La \textbf{zone tampon} absorbe les pressions humaines (usages encadrés) et protège le cœur. La \textbf{zone de transition} accueille les activités des populations (agriculture, habitat) et les associe à la conservation. Les grands mammifères (lion, girafe, damalisque) ont besoin de \textbf{vastes territoires} pour chasser, se nourrir et se reproduire, et de populations suffisamment grandes pour éviter la consanguinité : la zone cœur doit donc être la plus grande possible.
    \item \textbf{Conflit éleveurs / conservation :} en saison sèche, les troupeaux de zébus pénètrent le parc pour l'eau et l'herbe ; le \textbf{surpâturage} et le piétinement dégradent la végétation, le bétail entre en \textbf{compétition} avec la faune sauvage pour les ressources et peut lui transmettre des \textbf{maladies}. \textbf{Solutions concertées :}
    \begin{itemize}
        \item Aménager des \textbf{couloirs de transhumance} balisés hors zone cœur et des \textbf{points d'eau en périphérie} du parc, définis avec les éleveurs.
        \item Organiser la \textbf{vaccination du bétail} et un calendrier de passage concerté, avec des comités mixtes éleveurs--écogardes.
    \end{itemize}
    \item \textbf{Bioamplification du mercure :} le mercure rejeté par les orpailleurs se transforme en méthylmercure dans les sédiments ; il est absorbé par les petits organismes aquatiques, puis concentré dans les petits poissons, puis encore davantage dans les gros poissons et les oiseaux piscivores : à chaque maillon de la chaîne trophique, la concentration augmente car le mercure s'accumule dans les graisses sans être éliminé. L'\textbf{homme}, en tête de chaîne (consommation de poissons), reçoit les doses les plus fortes : risque d'\textbf{intoxication chronique} (atteintes du système nerveux, malformations chez les enfants).
    \item \textbf{Corridors écologiques :} un corridor écologique est une bande de milieu naturel reliant deux habitats isolés, permettant les déplacements et les échanges entre populations. Connecter la Benoué à Bouba Ndjidda et à la zone de chasse du Sud est vital : cela évite l'\textbf{isolement génétique} (consanguinité), permet les migrations saisonnières vers l'eau et la nourriture et la recolonisation des zones où une espèce a disparu. \textbf{Espèces « parapluie » :} l'\textbf{éléphant de savane} et le \textbf{lion} (leur protection, exigeant de grands espaces, protège de fait tout l'écosystème).
    \item \textbf{Interprétation :} la hausse de \num{15}\% dans la zone cœur montre que la protection intégrale fonctionne. La baisse de \num{30}\% en zone tampon révèle que les pressions (braconnage, pâturage, orpaillage) y sont trop fortes ou mal contrôlées : la zone tampon ne joue plus son rôle et menace le cœur à terme. \textbf{Ajustement :} renforcer la surveillance et l'encadrement des usages en zone tampon (patrouilles, quotas de chasse villageoise stricts, limitation du pâturage), voire réduire temporairement les activités autorisées le temps que les populations de grands mammifères se reconstituent.
\end{enumerate}

\textbf{Points de vigilance :} en Q3, expliciter le mécanisme d'accumulation le long de la chaîne (c'est le cœur de la notion de bioamplification) ; en Q4, bien définir le corridor avant de justifier ; en Q5, l'interprétation doit être \emph{contrastée} (succès du cœur / échec du tampon) et l'ajustement proposé doit être réaliste et applicable.
\end{corrige}

\newpage

\coursec{Fiche bilan}

\begin{popfigure}
\begin{tikzpicture}[
    mindmap,
    grow cyclic,
    every node/.style={concept, execute at begin node=\hskip0pt},
    concept color=popblue!80,
    level 1/.append style={level distance=4.5cm, sibling angle=360/6, font=\small\bfseries},
    level 2/.append style={level distance=3cm, sibling angle=360/4, font=\footnotesize},
    branch1/.style={concept color=popgreen!70},
    branch2/.style={concept color=poporange!70},
    branch3/.style={concept color=poppurple!70},
    branch4/.style={concept color=poppink!70},
    branch5/.style={concept color=popgold!70},
    branch6/.style={concept color=popblue!70},
    text=popdark
]
\node[root concept, font=\large\bfseries, text width=3.5cm, align=center] {Écosystèmes\\Forêt \& Savane}
    child[branch1] { node[concept, text width=2.8cm, align=center] {Notion d'écosystème}
        child { node {Biocénose + Biotope} }
        child { node {Facteurs biotiques/abiotiques} }
        child { node {Milieu forestier / sahélien} }
    }
    child[branch2] { node[concept, text width=2.8cm, align=center] {Grands écosystèmes du Cameroun}
        child { node {Forêt dense humide (Sud)} }
        child { node {Savane boisée / herbeuse (Nord)} }
        child { node {Mangrove (Littoral)} }
        child { node {Montagne (Mont Cameroun)} }
    }
    child[branch3] { node[concept, text width=2.8cm, align=center] {Biodiversité spécifique}
        child { node {Forêt : strates, endémisme fort} }
        child { node {Savane : adaptations xérophytes} }
        child { node {Richesse spécifique \& génétique} }
        child { node {Services écosystémiques} }
    }
    child[branch4] { node[concept, text width=2.8cm, align=center] {Interdépendance}
        child { node {Réseaux trophiques (P/C/D)} }
        child { node {Cycles de la matière (C, N, H₂O)} }
        child { node {Flux d'énergie (productivité)} }
        child { node {Relations symbiotiques} }
    }
    child[branch5] { node[concept, text width=2.8cm, align=center] {Menaces anthropiques}
        child { node {Déforestation (bois, agriculture)} }
        child { node {Feux de brousse (départ tardif)} }
        child { node {Braconnage / Surpêche} }
        child { node {Pollution / Urbanisation} }
    }
    child[branch6] { node[concept, text width=2.8cm, align=center] {Conservation \& Restauration}
        child { node {Aires protégées (Parcs, Réserves)} }
        child { node {Reboisement (TP N°10)} }
        child { node {Législation (Loi 94/01)} }
        child { node {Éducation / Écocitoyenneté} }
    };
\end{tikzpicture}
\legende{Carte mentale récapitulative de la leçon : structure, biodiversité, fonctionnement, menaces et solutions pour les écosystèmes forestiers et savanicoles du Cameroun.}
\end{popfigure}

\begin{bilan}

\courssub{Définitions clés à connaître par cœur}

\begin{itemize}
    \item \textbf{Écosystème} : Ensemble formé par une communauté d'êtres vivants (\textbf{biocénose}) et son milieu de vie physique et chimique (\textbf{biotope}), en interaction dynamique.
    \item \textbf{Biodiversité} : Variabilité des organismes vivants (diversité spécifique, génétique, écosystémique). La forêt dense humide du bassin du Congo (Sud-Cameroun) est un \emph{hotspot} mondial.
    \item \textbf{Producteurs (autotrophes)} : Êtres vivants (végétaux chlorophylliens, cyanobactéries) synthétisant leur matière organique par photosynthèse à partir de CO₂, H₂O, lumière. Base des réseaux trophiques.
    \item \textbf{Consommateurs (hétérotrophes)} : Primaires (herbivores), secondaires (carnivores), tertiaires (superprédateurs). Se nourrissent d'autres êtres vivants.
    \item \textbf{Décomposeurs (réducteurs)} : Bactéries, champignons. Minéralisent la matière organique morte $\rightarrow$ sels minéraux réutilisables par les producteurs (fermeture du cycle).
    \item \textbf{Chaîne trophique} : Suite linéaire « mangé par » (ex. Herbe $\rightarrow$ Gazelle $\rightarrow$ Lion $\rightarrow$ Bactéries). \textbf{Réseau trophique} : Ensemble imbriqué des chaînes trophiques d'un écosystème.
    \item \textbf{Pyramide des biomasses / de l'énergie} : L'énergie diminue d'un niveau trophique au suivant (rendement \SI{\sim 10}{\%}), limitant le nombre de niveaux (généralement 4--5).
\end{itemize}

\courssub{Les deux grands écosystèmes terrestres du Cameroun (comparatif)}

\begin{center}
\begin{tabularx}{\linewidth}{|l|X|X|}
\hline
\rowcolor{popblueL}
\textbf{Critère} & \textbf{Forêt dense humide (Sud, Est)} & \textbf{Savane (Nord, Adamaoua)} \\
\hline
Climat & Équatorial : chaud, humide, pluies abondantes (\SI{>1500}{mm/an}) & Tropical sec (soudanien) : saison sèche marquée, \SI{800-1200}{mm/an} \\
\hline
Végétation (Flore) & Strates multiples (émergents \SI{>40}{m}, canopée, sous-bois, herbacée). Lianes, épiphytes. Essences : Ayous, Sapelli, Iroko, Moabi. & Formation herbacée continue (graminées) + arbres/arbustes épars (xérophytes : karité, néré, fromager). Adaptations : racines profondes, feuilles coriaces, écorces épaisses. \\
\hline
Faune caractéristique & Primates (gorille, chimpanzé), éléphant de forêt, bongo, oiseaux (calao), grande diversité insectes/amphibiens. & Grands ongulés (kob, buffle, girafe), prédateurs (lion, guépard, hyène), avifaune nombreuse, termites (constructeurs). \\
\hline
Sol & Ferrallitique, lessivé, pauvre en bases, humus de surface mince (cyclage rapide). & Ferrugineux tropical, lessivé en saison humide, durcissement en saison sèche (croûte). \\
\hline
Productivité primaire & Très élevée (biomasse importante, croissance continue). & Saisonnière (pic en saison des pluies, arrêt en saison sèche). Biomasse aérienne plus faible. \\
\hline
\end{tabularx}
\end{center}

\courssub{Interdépendance et fonctionnement}

\begin{itemize}
    \item \textbf{Cycles biogéochimiques} : Carbone (photosynthèse/respiration), Azote (fixation $\rightarrow$ assimilation $\rightarrow$ ammonification $\rightarrow$ nitrification $\rightarrow$ dénitrification), Eau (évapotranspiration, ruissellement, infiltration). \emph{Rôle clé des décomposeurs}.
    \item \textbf{Flux d'énergie} : Unidirectionnel (Soleil $\rightarrow$ Producteurs $\rightarrow$ Consommateurs $\rightarrow$ Chaleur dissipée). Non recyclable.
    \item \textbf{Régulations naturelles} : Proies-Prédateurs (Lotka-Volterra simplifié), compétition, symbiose (mycorhizes forêt, acacia-fourmis savane), portance du milieu (K).
\end{itemize}

\courssub{Activités humaines destructrices (contexte camerounais)}

\begin{enumerate}
    \item \textbf{Déforestation} : Exploitation forestière (industrielle/artisanale), agriculture sur brûlis (cacao, huile de palme, vivrier), coupe de bois de feu/charbon. \emph{Conséquence} : Perte d'habitat, érosion, perturbation cycle de l'eau, émission CO₂.
    \item \textbf{Feux de brousse} : Pratiques culturales, chasse au feu, négligence. \emph{Conséquence} : Destruction graines/jeunes plants, érosion, émission gaz à effet de serre, régression forêt $\rightarrow$ savane anthropique.
    \item \textbf{Braconnage} : Viande de brousse (commerce urbain), ivoire, écailles de pangolin. \emph{Conséquence} : Défaunation (« forêts vides »), rupture réseaux trophiques, disparition espèces clés (éléphant = dispersant de graines).
    \item \textbf{Autres} : Pollution (pesticides coton Nord, déchets urbains Douala/Yaoundé), urbanisation galopante, espèces invasives.
\end{enumerate}

\courssub{Restauration et conservation de la biodiversité}

\begin{itemize}
    \item \textbf{In situ} : Réseau d'aires protégées (\SI{\sim 20}{\%} territoire) : Parcs Nationaux (Waza, Bouba Ndjida, Korup, Lobéké, Campo Ma'an), Réserves de Faune (Dja - Patrimoine UNESCO), Zones de Chasse. Gestion participative (COVAREF).
    \item \textbf{Ex situ} : Jardins botaniques (Limbe), zoos (Mvog-Betsi), banques de gènes (IRAD).
    \item \textbf{Reboïsement / Restauration (TP N°10)} : Choix essences locales (adaptées, à croissance rapide ou à haute valeur), pépinière (sacs, substrat, ombrière, arrosage), plantation (trous \SI{40x40x40}{cm}, amendement), \textbf{suivi-protection} (pare-feu, garde, remplacement plants morts).
    \item \textbf{Cadre légal} : Loi n°94/01 (forêt, faune, pêche), Décret 95/531 (EIES), Convention CITES, CBD. Lutte contre exploitation illégale (certification FSC/PEFC).
    \item \textbf{Éducation / Écocitoyenneté} : Clubs nature, Journée Mondiale de l'Environnement (5 juin), Agroforesterie (cacao sous ombrage, parcs à karité/néré).
\end{itemize}

\courssub{Méthode express : TP N°10 – Pratique du reboisement (étapes clés)}

\begin{enumerate}
    \item \textbf{Diagnostic site} : Type sol, pente, pluviométrie, essences endémiques voisines.
    \item \textbf{Sélection essences} : Rustiques, à croissance rapide (Acacia, Leucaena, Gmelina) \emph{ou} à haute valeur (Moabi, Ayous, Karité) selon objectif (bois d'œuvre, fertilité, fruitier).
    \item \textbf{Pépinière} : Remplissage sachets (terre + fumier décomposé), semis/bouturage, ombrière (\SI{50}{\%} ombre), arrosage régulier, désherbage, durcissement (réduction ombre/arrosage avant plantation).
    \item \textbf{Préparation terrain} : Traçage (alignement/contour), piquetage, ouverture trous \SI{40x40x40}{cm} (terre de surface / profondeur séparées), amendement (fumier/compost).
    \item \textbf{Plantation} : Saison des pluies (juin-août Sud, mai-juin Nord). Retrait sachet délicat, placement motte, bourrage terre fine, tasseau, arrosage d'appoint (\SI{10}{L}).
    \item \textbf{Entretien \& Suivi} : Désherbage \SI{1}{m} autour (3--4/an les 2 premières années), pare-feu (\SI{5}{m} large), remplacement mortalité (\textsuperscript{é}claircie), mesure croissance (hauteur, diamètre \SI{1,30}{m}).
\end{enumerate}

\end{bilan}

\begin{aretenir}[Checklist avant le BEPC]
\begin{itemize}
    \item $\square$ Définir écosystème, biocénose, biotope et citer les composants abiotiques/biotiques.
    \item $\square$ Localiser et caractériser (climat, sol, flore, faune) la forêt dense humide et la savane au Cameroun.
    \item $\square$ Expliquer la stratification de la forêt et les adaptations xérophytes de la savane.
    \item $\square$ Construire et interpréter une chaîne et un réseau trophique (identifier P, C1, C2, C3, D).
    \item $\square$ Expliquer le rôle des décomposeurs dans le recyclage de la matière (cycles C, N, eau).
    \item $\square$ Justifier la forme pyramidale de la pyramide des biomasses et de l'énergie (rendement \SI{\sim 10}{\%}).
    \item $\square$ Citer 4 activités humaines destructrices et expliquer leurs conséquences sur la biodiversité et le fonctionnement écosystémique.
    \item $\square$ Différencier conservation \emph{in situ} (aires protégées) et \emph{ex situ} ; citer 3 parcs nationaux camerounais.
    \item $\square$ Décrire les 6 étapes techniques du reboisement (TP N°10) de la pépinière au suivi.
    \item $\square$ Argumenter l'importance de la biodiversité (services écosystémiques : approvisionnement, régulation, support, culture).
    \item $\square$ Analyser un document (carte, graphique, texte) sur l'évolution de la couverture forestière ou la dynamique d'une population animale.
    \item $\square$ Proposer des mesures concrètes de gestion durable (agroforesterie, écotourisme, lutte anti-braconnage, éducation).
\end{itemize}
\end{aretenir}

\findecours

\end{document}
